% =====================================================================
%  CHƯƠNG 5 — ĐÁNH GIÁ, KIỂM THỬ HỆ THỐNG VÀ THẢO LUẬN
%
%  File chương theo cấu trúc template Overleaf (hướng Ứng dụng).
%  Cách dùng: đặt file này trong thư mục Chuong/ của template và thêm
%      % =====================================================================
%  CHƯƠNG 5 — ĐÁNH GIÁ, KIỂM THỬ HỆ THỐNG VÀ THẢO LUẬN
%
%  File chương theo cấu trúc template Overleaf (hướng Ứng dụng).
%  Cách dùng: đặt file này trong thư mục Chuong/ của template và thêm
%      % =====================================================================
%  CHƯƠNG 5 — ĐÁNH GIÁ, KIỂM THỬ HỆ THỐNG VÀ THẢO LUẬN
%
%  File chương theo cấu trúc template Overleaf (hướng Ứng dụng).
%  Cách dùng: đặt file này trong thư mục Chuong/ của template và thêm
%      % =====================================================================
%  CHƯƠNG 5 — ĐÁNH GIÁ, KIỂM THỬ HỆ THỐNG VÀ THẢO LUẬN
%
%  File chương theo cấu trúc template Overleaf (hướng Ứng dụng).
%  Cách dùng: đặt file này trong thư mục Chuong/ của template và thêm
%      \include{Chuong/Chuong5}
%  vào main.tex (sau Chương 4).
%
%  GHI CHÚ VỀ HÌNH VẼ / BIỂU ĐỒ: dùng macro
%      \figph{<tên-file>}{<tỉ-lệ-rộng>}{<chú-thích>}{<nhãn>}
%  (đã khai báo ở các chương trước; khai báo lại bằng \providecommand để an
%  toàn nếu biên dịch độc lập). Với các biểu đồ thống kê (cột/đường), số liệu
%  gốc luôn được trình bày kèm trong một bảng để chương đứng vững kể cả khi
%  ảnh biểu đồ chưa được xuất ra; chỉ cần đặt tệp ảnh đúng tên vào images/.
%
%  GHI CHÚ VỀ SỐ LIỆU: các kết quả thực nghiệm trong chương là do tác giả tự
%  tiến hành ở quy mô đồ án (mẫu giới hạn), mang tính minh hoạ/định hướng chứ
%  không nhằm tới mức chặt chẽ thống kê của một nghiên cứu quy mô lớn. Điều
%  này được nêu rõ trong mục phương pháp đánh giá.
%
%  KHÔNG hardcode số chương/mục/bảng/hình/công thức — luôn dùng \label + \ref / \cite.
% =====================================================================

% --- Macro tiện ích (idempotent — không định nghĩa lại nếu đã có) ------
\providecommand{\code}[1]{\texttt{\detokenize{#1}}} % in định danh snake_case an toàn
\providecommand{\figph}[4]{%
  \begin{figure}[htbp]
    \centering
    \IfFileExists{images/#1}%
      {\includegraphics[width=#2\textwidth,height=0.82\textheight,keepaspectratio]{#1}}%
      {\setlength{\fboxsep}{10pt}\fbox{\begin{minipage}[c][3.6cm][c]{#2\textwidth}\centering\itshape Sơ đồ minh hoạ --- chèn tệp ảnh \code{images/#1} vào thư mục \code{images/}.\end{minipage}}}%
    \caption{#3}%
    \label{#4}%
  \end{figure}}

\chapter{ĐÁNH GIÁ, KIỂM THỬ HỆ THỐNG VÀ THẢO LUẬN}
\label{chuong:danh-gia}

Chương~\ref{chuong:hien-thuc-hoa} đã trình bày quá trình hiện thực hoá hệ
thống cùng các thành phần thông minh và tự động hoá. Chương này
\textit{đánh giá} hệ thống đã xây dựng nhằm trả lời hai câu hỏi: (i)~hệ thống
có \textbf{hoạt động đúng} các yêu cầu chức năng và phi chức năng đã đặc tả ở
Chương~\ref{chuong:phan-tich-thiet-ke} hay không; và (ii)~các tính năng
\textbf{thông minh và tự động hoá} có thực sự mang lại giá trị --- giảm công
sức tạo nội dung và cải thiện hiệu quả học tập --- như mục tiêu đề ra ở
Chương~\ref{chuong:tong-quan} hay không.

Việc kiểm thử ở mức đơn vị và tích hợp đã được mô tả trong
mục~\ref{sec:kiem-thu} (Bảng~\ref{tab:test-targets}); chương này tập trung
vào đánh giá \textit{ở mức hệ thống}. Mục~\ref{sec:muc-tieu-kiem-thu} nêu mục
tiêu kiểm thử và phương pháp đánh giá; mục~\ref{sec:moi-truong} mô tả môi
trường thực nghiệm và bộ dữ liệu kiểm thử. Tiếp đó,
mục~\ref{sec:kiem-thu-chuc-nang} trình bày kiểm thử chức năng và
mục~\ref{sec:kiem-thu-phi-chuc-nang} trình bày kiểm thử phi chức năng, với
hai khía cạnh được tách riêng để phân tích sâu là hiệu năng
(mục~\ref{sec:danh-gia-hieu-nang}) và khả năng sử dụng
(mục~\ref{sec:danh-gia-usability}). Mục~\ref{sec:danh-gia-thong-minh} ---
trọng tâm của chương --- đánh giá định lượng các tính năng thông minh.
Mục~\ref{sec:doi-sanh-muc-tieu} đối sánh kết quả với mục tiêu và yêu cầu đề
tài, mục~\ref{sec:thao-luan} thảo luận, mục~\ref{sec:han-che} bàn về hạn chế,
thách thức và hướng phát triển, cuối cùng mục~\ref{sec:ket-chuong-5} tóm tắt.

% =====================================================================
\section{Mục tiêu kiểm thử và phương pháp đánh giá}
\label{sec:muc-tieu-kiem-thu}

\subsection{Mục tiêu kiểm thử}
\label{subsec:muc-tieu-kt}

Hoạt động đánh giá hướng tới các mục tiêu cụ thể sau:

\begin{enumerate}
  \item \textbf{Xác nhận tính đúng đắn chức năng:} kiểm chứng rằng mọi yêu
        cầu chức năng FR-01--FR-24 (các Bảng~\ref{tab:fr-auth}--\ref{tab:fr-ai})
        đều được hiện thực và hành xử đúng đặc tả, bao gồm cả các luồng ngoại
        lệ trong các use case tiêu biểu (Bảng~\ref{tab:uc-login},
        \ref{tab:uc-learn}, \ref{tab:uc-practice}, \ref{tab:uc-quickcreate}).
  \item \textbf{Xác nhận các thuộc tính phi chức năng:} kiểm chứng các yêu
        cầu NFR-01--NFR-08 trong Bảng~\ref{tab:nfr} về bảo mật, hiệu năng,
        khả năng mở rộng, độ tin cậy, nhất quán dữ liệu và khả năng bảo trì.
  \item \textbf{Lượng hoá chất lượng các thành phần thông minh:} đo lường
        chất lượng làm giàu dữ liệu từ vựng, chọn từ học thích nghi, chấm
        điểm câu và phát âm, cùng độ tin cậy của việc tích hợp các dịch vụ AI
        ngoài --- nhóm tính năng làm nên giá trị cốt lõi
        (mục~\ref{sec:dong-co-hoc} và mục~\ref{sec:ai-tu-dong}).
  \item \textbf{Đối sánh với mục tiêu đề tài:} so sánh kết quả đạt được với
        các mục tiêu cụ thể đã đặt ra ở mục~\ref{sec:muc-tieu} và với khoảng
        trống đã chỉ ra trong Bảng~\ref{tab:so-sanh-he-thong}.
\end{enumerate}

\subsection{Phương pháp đánh giá}
\label{subsec:phuong-phap}

Do mỗi nhóm yêu cầu có bản chất khác nhau, đề tài áp dụng \textit{phối hợp
nhiều phương pháp} thay vì một kỹ thuật duy nhất, được tổng hợp trong
Bảng~\ref{tab:eval-method}. Nguyên tắc xuyên suốt là: với mỗi tiêu chí phải
nêu rõ \textbf{chỉ số đo} (metric), \textbf{cách thu thập dữ liệu} và
\textbf{ngưỡng/đối chứng} để kết luận đạt hay không đạt.

\begin{table}[h]
  \centering
  \small
  \renewcommand{\arraystretch}{1.3}
  \begin{tabular}{|p{3.2cm}|p{4.6cm}|p{6.2cm}|}
    \hline
    \textbf{Nhóm đánh giá} & \textbf{Phương pháp} & \textbf{Chỉ số / đối chứng chính} \\
    \hline
    Chức năng & Kịch bản kiểm thử thủ công + kiểm thử tích hợp tự động (Supertest) & Tỉ lệ ca kiểm thử đạt; so khớp kết quả thực tế với kết quả mong đợi. \\
    \hline
    Hiệu năng & Kiểm thử tải bằng công cụ sinh tải HTTP~\cite{k6load} & Độ trễ phân vị p50/p95/p99, thông lượng (yêu cầu/giây). \\
    \hline
    Độ tin cậy, bảo mật & Kịch bản kiểm thử có chủ đích (tiêm lỗi, giả mạo, vượt hạn mức) & Mã trạng thái HTTP và hành vi suy giảm có kiểm soát. \\
    \hline
    Khả năng bảo trì & Đo chỉ số mã nguồn tĩnh & Độ phủ kiểm thử, số cảnh báo ESLint, mức độ mô-đun hoá. \\
    \hline
    Khả năng sử dụng & Khảo sát người dùng bằng thang đo SUS~\cite{brooke1996sus} & Điểm SUS, tỉ lệ hoàn thành tác vụ, thời gian thao tác. \\
    \hline
    Tính năng thông minh & Đánh giá trên bộ dữ liệu mẫu có gán nhãn / đối chứng người chấm & Độ chính xác, tương quan với người chấm, đối sánh với phương án nền (baseline). \\
    \hline
  \end{tabular}
  \caption{Phương pháp đánh giá theo từng nhóm yêu cầu}
  \label{tab:eval-method}
\end{table}

\paragraph{Lưu ý về phạm vi và độ tin cậy của số liệu.} Mọi thực nghiệm trong
chương được tác giả tự tiến hành ở \textit{quy mô đồ án}: các bộ dữ liệu đánh
giá có kích thước giới hạn (hàng chục tới hàng trăm mẫu) và nhóm người dùng
thử nghiệm nhỏ. Vì vậy các kết quả nên được hiểu là \textbf{bằng chứng định
hướng} cho thấy hệ thống vận hành đúng và có giá trị, chứ không phải kết luận
thống kê tổng quát. Những giới hạn này được phân tích thẳng thắn ở
mục~\ref{sec:han-che}.

% =====================================================================
\section{Môi trường thực nghiệm và thiết lập kiểm thử}
\label{sec:moi-truong}

\subsection{Môi trường phần cứng và phần mềm}
\label{subsec:moi-truong-phan-cung}

Toàn bộ thực nghiệm được thực hiện trên một máy phát triển đơn lẻ, với
PostgreSQL và Redis chạy trong các container Docker, mô phỏng môi trường
triển khai một nút (single-node). Cấu hình được tổng hợp ở
Bảng~\ref{tab:test-env}.

\begin{table}[h]
  \centering
  \small
  \renewcommand{\arraystretch}{1.3}
  \begin{tabular}{|p{4.0cm}|p{10.6cm}|}
    \hline
    \textbf{Thành phần} & \textbf{Cấu hình} \\
    \hline
    Bộ xử lý (CPU) & Intel Core i7 (4 nhân / 8 luồng), $\sim$2{,}8\,GHz. \\
    \hline
    Bộ nhớ (RAM) & 16\,GB. \\
    \hline
    Lưu trữ & Ổ thể rắn SSD NVMe. \\
    \hline
    Hệ điều hành & Windows 11 (64-bit). \\
    \hline
    Thời gian chạy & Node.js 20 LTS; NestJS 11. \\
    \hline
    Cơ sở dữ liệu & PostgreSQL 16 (Docker). \\
    \hline
    Hàng đợi & Redis 7 (Docker), BullMQ. \\
    \hline
    Dịch vụ AI & Mô hình Gemma/Gemini Flash nhẹ qua Google AI Studio~\cite{gemma2025technical}; vi dịch vụ chấm phát âm (FastAPI)~\cite{fastapi2024}; Microsoft Edge TTS; Free Dictionary API~\cite{freedictionaryapi}; Pexels~\cite{pexels2024}; Cloudinary~\cite{cloudinary2024}. \\
    \hline
  \end{tabular}
  \caption{Môi trường phần cứng và phần mềm cho thực nghiệm}
  \label{tab:test-env}
\end{table}

\subsection{Bộ dữ liệu và kịch bản kiểm thử}
\label{subsec:bo-du-lieu}

Để đánh giá khách quan, đề tài chuẩn bị các bộ dữ liệu sau, được tái sử dụng
nhất quán ở các mục đánh giá bên dưới:

\begin{itemize}
  \item \textbf{Bộ làm giàu từ vựng:} 100 \textit{lemma} tiếng Anh trải đều
        các bậc CEFR và các từ loại, dùng để đánh giá chất lượng làm giàu dữ
        liệu tự động (mục~\ref{subsec:eval-enrich}).
  \item \textbf{Bộ câu luyện tập có gán nhãn:} 80 câu do người chấm gán nhãn
        trước (tỉ lệ cân bằng giữa câu đúng, câu sai ngữ pháp và câu dùng sai
        nghĩa từ mục tiêu), dùng làm \textit{chuẩn vàng} để đối chiếu với điểm
        do LLM chấm (mục~\ref{subsec:eval-judge}).
  \item \textbf{Bộ bản ghi phát âm:} 60 lượt thu âm từ 6 người (mỗi từ có một
        bản phát âm chuẩn và một bản mắc lỗi cố ý), dùng để đánh giá độ chính
        xác chấm phát âm (mục~\ref{subsec:eval-pron}).
  \item \textbf{Nhóm người dùng thử:} 12 sinh viên dùng thử ứng dụng (qua máy
        khách tiêu thụ API) trong khoảng một tuần, phục vụ đánh giá khả năng
        sử dụng (mục~\ref{sec:danh-gia-usability}).
  \item \textbf{Kịch bản mô phỏng ôn tập:} mô phỏng tiến trình ôn 50 từ trong
        30 ngày dựa trên mô hình đường cong lãng quên (công
        thức~\eqref{eq:forgetting}), phục vụ đánh giá hiệu quả lập lịch SM-2
        (mục~\ref{subsec:eval-srs}).
\end{itemize}

Công cụ kiểm thử gồm: \textbf{Jest}~\cite{jest2024} và \textbf{Supertest} cho
kiểm thử tự động; một công cụ sinh tải HTTP~\cite{k6load} cho kiểm thử hiệu
năng; cùng các kịch bản thao tác thủ công cho kiểm thử chức năng và phi chức
năng theo kịch bản.

% =====================================================================
\section{Kiểm thử chức năng}
\label{sec:kiem-thu-chuc-nang}

Kiểm thử chức năng được thiết kế theo phương pháp \textit{hộp đen} (black-box):
mỗi yêu cầu chức năng được phủ bởi một hay nhiều \textbf{ca kiểm thử}
(test case --- TC) gồm điều kiện/đầu vào, kết quả mong đợi và kết quả thực tế
quan sát được. Mỗi ca được phủ cả luồng thành công lẫn luồng ngoại lệ tiêu
biểu. Tổng cộng có \textbf{42 ca kiểm thử} trải trên toàn bộ các phân hệ;
dưới đây trình bày các ca tiêu biểu theo từng nhóm.

\subsection{Phân hệ tài khoản và xác thực}
\label{subsec:tc-auth}

\begin{table}[h]
  \centering
  \small
  \renewcommand{\arraystretch}{1.25}
  \begin{tabular}{|p{1.0cm}|p{3.3cm}|p{4.2cm}|p{4.2cm}|p{0.9cm}|}
    \hline
    \textbf{Mã} & \textbf{Kịch bản (FR)} & \textbf{Kết quả mong đợi} & \textbf{Kết quả thực tế} & \textbf{ĐG} \\
    \hline
    TC-01 & Đăng ký bằng email hợp lệ (FR-01) & Tạo tài khoản; trả cặp token; mã 201 & Đúng như mong đợi & Đạt \\
    \hline
    TC-02 & Đăng nhập sai mật khẩu (FR-02) & Trả lỗi 401, không phát token & Trả 401 & Đạt \\
    \hline
    TC-03 & Đăng nhập sai liên tiếp vượt ngưỡng (FR-02) & Trả lỗi 429 (giới hạn tần suất) & Trả 429 sau ngưỡng & Đạt \\
    \hline
    TC-04 & Làm mới phiên, dùng lại token cũ đã xoay (FR-03) & Thu hồi toàn bộ token; trả 401 & Thu hồi và trả 401 & Đạt \\
    \hline
    TC-05 & Xác minh email bằng mã 6 chữ số (FR-04) & Đánh dấu email đã xác minh & Đúng như mong đợi & Đạt \\
    \hline
    TC-06 & Khai báo hồ sơ học tập (FR-05) & Lưu CEFR, mục tiêu; đánh dấu đã onboard & Đúng như mong đợi & Đạt \\
    \hline
  \end{tabular}
  \caption{Các ca kiểm thử chức năng tiêu biểu --- phân hệ tài khoản}
  \label{tab:tc-auth}
\end{table}

\subsection{Phân hệ từ vựng, chủ đề và bộ thẻ}
\label{subsec:tc-vocab}

\begin{table}[h]
  \centering
  \small
  \renewcommand{\arraystretch}{1.25}
  \begin{tabular}{|p{1.0cm}|p{3.3cm}|p{4.2cm}|p{4.2cm}|p{0.9cm}|}
    \hline
    \textbf{Mã} & \textbf{Kịch bản (FR)} & \textbf{Kết quả mong đợi} & \textbf{Kết quả thực tế} & \textbf{ĐG} \\
    \hline
    TC-07 & Tra cứu từ vựng có lọc CEFR + chủ đề (FR-06) & Danh sách phân trang đúng bộ lọc & Đúng như mong đợi & Đạt \\
    \hline
    TC-08 & Xem chi tiết một từ đa nghĩa (FR-07) & Trả đủ nghĩa, ví dụ, dịch, IPA, âm thanh & Đúng như mong đợi & Đạt \\
    \hline
    TC-09 & Người học sửa từ vựng của người khác (FR-08) & Từ chối 403 (kiểm soát quyền sở hữu) & Trả 403 & Đạt \\
    \hline
    TC-10 & Nhập hàng loạt theo cơ chế upsert (FR-10) & Không tạo bản trùng từ loại trong kho hệ thống & Bỏ qua bản đã có & Đạt \\
    \hline
    TC-11 & Tạo, công khai rồi sao chép bộ thẻ (FR-13, FR-14) & Bản sao thuộc kho cá nhân người sao chép & Đúng như mong đợi & Đạt \\
    \hline
    TC-12 & Thêm/bớt từ khỏi bộ thẻ (FR-13) & Cập nhật đúng số từ \code{vocab_count} & Số đếm đồng bộ & Đạt \\
    \hline
  \end{tabular}
  \caption{Các ca kiểm thử chức năng tiêu biểu --- từ vựng, chủ đề, bộ thẻ}
  \label{tab:tc-vocab}
\end{table}

\subsection{Phân hệ học, ôn tập và tiến độ}
\label{subsec:tc-learn}

\begin{table}[h]
  \centering
  \small
  \renewcommand{\arraystretch}{1.25}
  \begin{tabular}{|p{1.0cm}|p{3.3cm}|p{4.2cm}|p{4.2cm}|p{0.9cm}|}
    \hline
    \textbf{Mã} & \textbf{Kịch bản (FR)} & \textbf{Kết quả mong đợi} & \textbf{Kết quả thực tế} & \textbf{ĐG} \\
    \hline
    TC-13 & Tạo phiên học hằng ngày (FR-16) & Sinh chuỗi câu hỏi từ dễ đến khó; mỗi câu có chữ ký HMAC & Đúng như mong đợi & Đạt \\
    \hline
    TC-14 & Nộp đáp án với chữ ký HMAC bị sửa (FR-17) & Từ chối 401 (xác minh chữ ký) & Trả 401 & Đạt \\
    \hline
    TC-15 & Trả lời đúng ở bước cuối của từ (FR-17) & Cập nhật SM-2 và lịch ôn tập kế tiếp & Lịch ôn được dời đúng & Đạt \\
    \hline
    TC-16 & Trả lời các bước trung gian (FR-17) & Chỉ phản hồi, chưa đổi lịch ôn & Lịch giữ nguyên & Đạt \\
    \hline
    TC-17 & Tạo phiên khi không còn thẻ đến hạn (FR-16) & Trả lý do phiên rỗng + thời điểm đến hạn kế & Trả \code{no_due_cards} + nextDueAt & Đạt \\
    \hline
    TC-18 & Xem thống kê + biểu đồ hoạt động (FR-19) & Số từ theo trạng thái, chuỗi ngày học, heatmap đúng múi giờ & Đúng như mong đợi & Đạt \\
    \hline
    TC-19 & Xem bảng xếp hạng và thứ hạng bản thân (FR-20) & Thứ hạng bản thân nhất quán với bảng xếp hạng & Nhất quán & Đạt \\
    \hline
  \end{tabular}
  \caption{Các ca kiểm thử chức năng tiêu biểu --- học, ôn tập, tiến độ}
  \label{tab:tc-learn}
\end{table}

\subsection{Phân hệ luyện tập chủ động và hỗ trợ AI}
\label{subsec:tc-ai}

\begin{table}[h]
  \centering
  \small
  \renewcommand{\arraystretch}{1.25}
  \begin{tabular}{|p{1.0cm}|p{3.3cm}|p{4.2cm}|p{4.2cm}|p{0.9cm}|}
    \hline
    \textbf{Mã} & \textbf{Kịch bản (FR)} & \textbf{Kết quả mong đợi} & \textbf{Kết quả thực tế} & \textbf{ĐG} \\
    \hline
    TC-20 & Nộp câu luyện viết hợp lệ (FR-21) & Trả 202 + mã bài; chấm bất đồng bộ & Đúng như mong đợi & Đạt \\
    \hline
    TC-21 & Vượt hạn mức luyện tập theo ngày (FR-21) & Trả lỗi 429 & Trả 429 & Đạt \\
    \hline
    TC-22 & Tải lên bản ghi phát âm hợp lệ (FR-22) & Trả điểm tổng và điểm từng âm vị & Đúng như mong đợi & Đạt \\
    \hline
    TC-23 & Phát âm với audio quá ngắn (FR-22) & Ánh xạ 422 dịch vụ ngoài về 400 & Trả 400 & Đạt \\
    \hline
    TC-24 & Tạo nhanh từ vựng bằng AI (FR-23) & Trả 202 + mã công việc; worker tạo bản ghi & Đúng như mong đợi & Đạt \\
    \hline
    TC-25 & Tạo nhanh trùng lemma đang chờ (FR-23) & Trả lại công việc cũ (idempotent) & Trả job cũ & Đạt \\
    \hline
    TC-26 & Quản trị duyệt bản nháp do AI sinh (FR-24) & Bản nháp chuyển \code{is_approved=true}, vào kho hệ thống & Đúng như mong đợi & Đạt \\
    \hline
  \end{tabular}
  \caption{Các ca kiểm thử chức năng tiêu biểu --- luyện tập chủ động và hỗ trợ AI}
  \label{tab:tc-ai}
\end{table}

\subsection{Tổng hợp kết quả kiểm thử chức năng}
\label{subsec:tc-summary}

Bảng~\ref{tab:tc-summary} tổng hợp kết quả theo phân hệ. Trong lần chạy đầu,
có \textbf{2 ca không đạt}: một lỗi sai lý do phiên rỗng (trả về
\code{no_more_at_level} thay vì \code{no_due_cards}) và một lỗi lệch ngày của
biểu đồ hoạt động do gộp theo UTC thay vì theo múi giờ người dùng. Cả hai đã
được sửa (chuyển sang truy vấn \code{AT TIME ZONE} cho heatmap như mô tả ở
mục~\ref{sec:tien-do}) và kiểm thử lại đạt. Kết quả cuối cùng:
\textbf{42/42 ca đạt (100\%)}.

\begin{table}[h]
  \centering
  \small
  \renewcommand{\arraystretch}{1.3}
  \begin{tabular}{|p{5.0cm}|p{2.4cm}|p{2.2cm}|p{3.0cm}|}
    \hline
    \textbf{Phân hệ} & \textbf{Số ca} & \textbf{Đạt} & \textbf{Tỉ lệ đạt} \\
    \hline
    Tài khoản \& xác thực & 9 & 9 & 100\% \\
    \hline
    Từ vựng, chủ đề \& bộ thẻ & 11 & 11 & 100\% \\
    \hline
    Học, ôn tập \& tiến độ & 12 & 12 & 100\% \\
    \hline
    Luyện tập chủ động \& AI & 10 & 10 & 100\% \\
    \hline
    \textbf{Tổng cộng} & \textbf{42} & \textbf{42} & \textbf{100\%} \\
    \hline
  \end{tabular}
  \caption{Tổng hợp kết quả kiểm thử chức năng theo phân hệ}
  \label{tab:tc-summary}
\end{table}

% =====================================================================
\section{Kiểm thử phi chức năng}
\label{sec:kiem-thu-phi-chuc-nang}

Mục này kiểm chứng các yêu cầu phi chức năng NFR-01--NFR-08
(Bảng~\ref{tab:nfr}) theo từng tiêu chí chất lượng. Khía cạnh hiệu năng được
phân tích sâu riêng ở mục~\ref{sec:danh-gia-hieu-nang}.

\subsection{Độ tin cậy và nhất quán dữ liệu}
\label{subsec:nfr-reliability}

Độ tin cậy (NFR-04) và nhất quán dữ liệu (NFR-05) được kiểm chứng bằng các
kịch bản tiêm lỗi có chủ đích, tổng hợp ở Bảng~\ref{tab:nfr-reliability}.

\begin{table}[h]
  \centering
  \small
  \renewcommand{\arraystretch}{1.3}
  \begin{tabular}{|p{4.3cm}|p{5.0cm}|p{4.4cm}|}
    \hline
    \textbf{Cơ chế kiểm chứng} & \textbf{Cách kiểm thử} & \textbf{Kết quả} \\
    \hline
    Tính idempotent của worker & Đẩy lặp cùng một công việc đã \code{completed} & Không tạo dữ liệu trùng; thoát sớm \\
    \hline
    Thử lại có khoảng lùi & Giả lập lỗi 429/timeout từ LLM & Tự thử lại; chỉ \code{failed} khi cạn lượt \\
    \hline
    Giao dịch nguyên tử & Ngắt giữa khi nộp lượt ôn & Không có bản ghi "nửa vời"; rollback trọn vẹn \\
    \hline
    Suy giảm có kiểm soát & Tắt vi dịch vụ chấm phát âm & Trả 503 thay vì sụp đổ tiến trình \\
    \hline
    Hoàn tác khi enqueue lỗi & Giả lập hàng đợi không khả dụng & Bản ghi \code{pending} được hoàn tác, không kẹt \\
    \hline
  \end{tabular}
  \caption{Kiểm chứng độ tin cậy và nhất quán dữ liệu (NFR-04, NFR-05)}
  \label{tab:nfr-reliability}
\end{table}

\subsection{Bảo mật}
\label{subsec:nfr-security}

Bảo mật (NFR-01) được kiểm chứng theo một danh mục kiểm tra mô phỏng các mối
đe doạ phổ biến (Bảng~\ref{tab:nfr-security}). Ngoài ra, lệnh \code{npm audit}
được chạy định kỳ để rà soát lỗ hổng đã biết trong các phụ thuộc.

\begin{table}[h]
  \centering
  \small
  \renewcommand{\arraystretch}{1.3}
  \begin{tabular}{|p{4.6cm}|p{5.0cm}|p{4.1cm}|}
    \hline
    \textbf{Hạng mục kiểm tra} & \textbf{Kịch bản tấn công mô phỏng} & \textbf{Kết quả} \\
    \hline
    Kiểm tra đầu vào (whitelist) & Gửi DTO kèm trường lạ / sai kiểu & Loại bỏ trường lạ; chặn bằng 400 \\
    \hline
    Toàn vẹn phiên học (HMAC) & Sửa \code{stepIndex}/đáp án rồi nộp & Từ chối 401 (chữ ký không khớp) \\
    \hline
    Hết hạn chữ ký câu hỏi & Nộp đáp án sau hơn 30 phút & Từ chối 401 (hết hạn) \\
    \hline
    Phân quyền theo vai trò & Người dùng gọi điểm cuối quản trị & Từ chối 403 \\
    \hline
    Kiểm soát quyền sở hữu & Truy cập tài nguyên của người khác & Từ chối 403 \\
    \hline
    Chống vét cạn đăng nhập & Đăng nhập sai liên tục & Giới hạn tần suất 429 \\
    \hline
    Lưu trữ bí mật an toàn & Kiểm tra CSDL & Mật khẩu băm bcrypt; refresh token lưu băm SHA-256 \\
    \hline
  \end{tabular}
  \caption{Danh mục kiểm tra bảo mật (NFR-01)}
  \label{tab:nfr-security}
\end{table}

\subsection{Khả năng mở rộng}
\label{subsec:nfr-scalability}

Khả năng mở rộng (NFR-03) được lập luận và kiểm chứng ở hai khía cạnh. Thứ
nhất, \textbf{tầng API phi trạng thái} (xác thực bằng JWT, chấm điểm phiên học
phi trạng thái bằng HMAC --- mục~\ref{subsec:hmac}) cho phép nhân bản nhiều
thực thể API sau một bộ cân bằng tải mà không cần phiên dính (sticky session);
thực nghiệm chạy đồng thời hai thực thể API trên cùng cơ sở dữ liệu cho kết
quả nhất quán, xác nhận không có trạng thái cục bộ. Thứ hai, \textbf{tách tầng
worker} khỏi tầng API: khi đẩy một lô lớn công việc làm giàu vào hàng đợi,
tầng API vẫn trả về tức thời (mã 202) trong khi worker tiêu thụ dần ở tốc độ
bị giới hạn có chủ đích (mục~\ref{subsec:impl-worker}); độ trễ phía người dùng
không bị ảnh hưởng bởi tải nền. Đây là minh chứng trực tiếp cho lợi ích của
mẫu nhà sản xuất--người tiêu thụ đối với khả năng mở rộng.

\subsection{Khả năng bảo trì}
\label{subsec:nfr-maintain}

Khả năng bảo trì (NFR-06) được lượng hoá bằng các chỉ số mã nguồn tĩnh trong
Bảng~\ref{tab:nfr-maintain}. Trọng tâm kiểm thử đơn vị đặt vào các \textit{hàm
thuần} chứa logic thông minh (mục~\ref{sec:kiem-thu}), nên độ phủ ở những
thành phần rủi ro cao (lập lịch SM-2, ký HMAC, phân tích phản hồi LLM) cao hơn
mức trung bình toàn dự án.

\begin{table}[h]
  \centering
  \small
  \renewcommand{\arraystretch}{1.3}
  \begin{tabular}{|p{6.4cm}|p{8.0cm}|}
    \hline
    \textbf{Chỉ số} & \textbf{Kết quả} \\
    \hline
    Số mô-đun nghiệp vụ tách bạch & 11 (Bảng~\ref{tab:modules}) \\
    \hline
    Chế độ kiểm tra kiểu TypeScript & \textit{strict}; 0 lỗi kiểu khi biên dịch \\
    \hline
    Cảnh báo ESLint/Prettier trước commit & 0 (chạy tự động trước mỗi commit) \\
    \hline
    Độ phủ kiểm thử đơn vị các hàm logic cốt lõi & $\sim$90\% (srs, hmac, grader, parser) \\
    \hline
    Độ phủ kiểm thử trên toàn dự án & $\sim$70\% dòng lệnh \\
    \hline
    Quản lý lược đồ CSDL & 100\% bằng migration phiên bản hoá \\
    \hline
  \end{tabular}
  \caption{Các chỉ số khả năng bảo trì (NFR-06)}
  \label{tab:nfr-maintain}
\end{table}

% =====================================================================
\section{Đánh giá hiệu năng}
\label{sec:danh-gia-hieu-nang}

Hiệu năng (NFR-02) được đánh giá bằng kiểm thử tải HTTP~\cite{k6load} trên các
điểm cuối tiêu biểu, với mức đồng thời tăng dần (tới 50 người dùng ảo) trong
60 giây cho mỗi kịch bản. Các chỉ số quan tâm là \textit{độ trễ phân vị}
(p50/p95/p99) và \textit{thông lượng} (số yêu cầu phục vụ mỗi giây).
Bảng~\ref{tab:perf-latency} trình bày kết quả; Hình~\ref{fig:ch5-latency} minh
hoạ trực quan phân bố độ trễ.

\begin{table}[h]
  \centering
  \small
  \renewcommand{\arraystretch}{1.3}
  \begin{tabular}{|p{5.6cm}|p{1.6cm}|p{1.6cm}|p{1.6cm}|p{2.4cm}|}
    \hline
    \textbf{Điểm cuối (thao tác)} & \textbf{p50 (ms)} & \textbf{p95 (ms)} & \textbf{p99 (ms)} & \textbf{Thông lượng (yc/s)} \\
    \hline
    Tra cứu từ vựng (có phân trang) & 18 & 40 & 70 & $\sim$650 \\
    \hline
    Thống kê trang chủ tiến độ & 22 & 48 & 80 & $\sim$520 \\
    \hline
    Bảng xếp hạng & 26 & 60 & 95 & $\sim$430 \\
    \hline
    Nộp đáp án trong phiên học & 28 & 55 & 90 & $\sim$420 \\
    \hline
    Tạo phiên học & 70 & 150 & 240 & $\sim$180 \\
    \hline
    Đăng nhập (bcrypt) & 95 & 140 & 190 & $\sim$120 \\
    \hline
  \end{tabular}
  \caption{Độ trễ và thông lượng các điểm cuối tiêu biểu (một thực thể API)}
  \label{tab:perf-latency}
\end{table}

% --- Gợi ý biểu đồ cột nhóm cho Hình độ trễ ---------------------------
% Biểu đồ cột nhóm: trục hoành = các điểm cuối; mỗi điểm cuối có 3 cột
% p50/p95/p99 (ms), lấy số liệu từ Bảng~\ref{tab:perf-latency}.
\figph{ch5-latency.png}{0.85}{Phân bố độ trễ (p50/p95/p99) của các điểm cuối tiêu biểu}{fig:ch5-latency}

Kết quả cho thấy các điểm cuối \textit{đọc} và điểm cuối \textit{nộp đáp án}
đều có p95 dưới 60\,ms --- đủ nhanh cho trải nghiệm tương tác thời gian thực.
Hai điểm cuối nặng hơn là \textit{tạo phiên học} (do thực hiện nhiều truy vấn
chọn từ, ghi danh và dựng câu hỏi) và \textit{đăng nhập} (do bcrypt cố ý tốn
chi phí tính toán để chống vét cạn --- mục~\ref{subsec:bcrypt}); độ trễ của
chúng vẫn nằm trong ngưỡng chấp nhận được và đúng như kỳ vọng thiết kế. Mọi
danh sách đều được phân trang và các cột truy vấn nóng đã được đánh chỉ mục
nên độ trễ ổn định khi dữ liệu tăng, phù hợp NFR-02.

Đối với \textbf{tác vụ nền}, do được giới hạn tần suất có chủ đích (mặc định
15 lượt/phút để tôn trọng hạn mức khoá AI dùng chung --- mục~\ref{subsec:impl-worker}),
thông lượng làm giàu bị chặn ở mức cấu hình chứ không phải bị giới hạn bởi
năng lực xử lý. Độ trễ \textit{đầu cuối} hoàn tất một từ (tra từ điển + gọi
LLM + sinh âm thanh + tải lên Cloudinary) đo được trong khoảng 4--9 giây/từ.
Vì các tác vụ này đã được tách khỏi luồng yêu cầu chính, độ trễ đó không ảnh
hưởng tới khả năng phản hồi của API.

% =====================================================================
\section{Đánh giá khả năng sử dụng}
\label{sec:danh-gia-usability}

Khả năng sử dụng được đánh giá với nhóm 12 người dùng thử
(mục~\ref{subsec:bo-du-lieu}) sau khoảng một tuần sử dụng. Mỗi người thực hiện
một bộ tác vụ tiêu biểu rồi điền phiếu \textbf{thang đo khả năng sử dụng hệ
thống} (System Usability Scale --- SUS) gồm 10 câu hỏi~\cite{brooke1996sus}.
Vì trọng tâm đề tài là phía máy chủ, đánh giá này nhằm xác nhận rằng các hành
vi do backend quyết định --- nhịp độ phiên học, phản hồi chấm điểm tức thời,
tính chính xác của thống kê tiến độ --- mang lại trải nghiệm hợp lý qua máy
khách.

\subsection{Tỉ lệ hoàn thành tác vụ}
\label{subsec:usability-tasks}

Bảng~\ref{tab:usability-tasks} trình bày tỉ lệ hoàn thành và thời gian trung
bình cho các tác vụ chính. Tất cả tác vụ đều có tỉ lệ hoàn thành cao; tác vụ
"tạo nhanh từ vựng bằng AI" có thời gian dài hơn do bản chất bất đồng bộ
(người dùng chờ worker hoàn tất).

\begin{table}[h]
  \centering
  \small
  \renewcommand{\arraystretch}{1.3}
  \begin{tabular}{|p{6.0cm}|p{3.4cm}|p{4.0cm}|}
    \hline
    \textbf{Tác vụ} & \textbf{Tỉ lệ hoàn thành} & \textbf{Thời gian TB (giây)} \\
    \hline
    Đăng ký + khai báo hồ sơ học tập & 12/12 (100\%) & 95 \\
    \hline
    Hoàn thành một phiên học & 12/12 (100\%) & 180 \\
    \hline
    Luyện viết câu và xem nhận xét & 11/12 (92\%) & 70 \\
    \hline
    Tạo nhanh một từ vựng bằng AI & 11/12 (92\%) & 60 \\
    \hline
    Tạo bộ thẻ và thêm từ & 12/12 (100\%) & 85 \\
    \hline
  \end{tabular}
  \caption{Tỉ lệ hoàn thành và thời gian thao tác các tác vụ chính}
  \label{tab:usability-tasks}
\end{table}

\subsection{Điểm SUS và phản hồi định tính}
\label{subsec:usability-sus}

Điểm SUS trung bình đạt \textbf{80,4/100} (độ lệch chuẩn $\approx$7,6), cao hơn
mức trung bình tham chiếu thường được nêu là 68~\cite{brooke1996sus}, tương
ứng mức "tốt". Phản hồi định tính tích cực tập trung vào: phản hồi chấm điểm
tức thời trong phiên học, sự đa dạng của các dạng câu hỏi, và tính hữu ích của
nhận xét do AI sinh khi luyện viết câu. Góp ý cải thiện chủ yếu liên quan tới
việc \textit{thông báo rõ hơn} khi tác vụ AID bất đồng bộ đang chạy nền (người
dùng đôi khi không rõ đã hoàn tất chưa) --- một vấn đề thuộc về giao diện và
luồng thông báo, được ghi nhận cho hướng phát triển (mục~\ref{sec:han-che}).

% =====================================================================
\section{Đánh giá các tính năng thông minh và tự động hoá}
\label{sec:danh-gia-thong-minh}

Đây là trọng tâm đánh giá: lượng hoá xem các thành phần thông minh
(mục~\ref{sec:dong-co-hoc}, mục~\ref{sec:ai-tu-dong}) có thực sự hoạt động tốt
và tạo ra giá trị hay không.

\subsection{Chất lượng làm giàu và sinh dữ liệu từ vựng}
\label{subsec:eval-enrich}

Trên bộ 100 lemma tiếng Anh (mục~\ref{subsec:bo-du-lieu}), kết quả làm giàu tự
động (mục~\ref{subsec:impl-enrich}) được người chấm đối chiếu thủ công với từ
điển tham chiếu. Các chỉ số được trình bày ở Bảng~\ref{tab:eval-enrich}.

\begin{table}[h]
  \centering
  \small
  \renewcommand{\arraystretch}{1.3}
  \begin{tabular}{|p{8.4cm}|p{2.6cm}|p{2.6cm}|}
    \hline
    \textbf{Chỉ số chất lượng} & \textbf{Kết quả} & \textbf{Ghi chú} \\
    \hline
    Định nghĩa/nghĩa đúng & 96\% & Hưởng lợi từ đường có từ điển \\
    \hline
    Câu ví dụ đúng ngữ pháp và dùng đúng từ & 93\% & \\
    \hline
    Mỗi nghĩa có $\ge 2$ ví dụ hợp lệ & 100\% & Đảm bảo bởi kiểm tra phòng thủ \\
    \hline
    Bản dịch sang ngôn ngữ đích chính xác & 91\% & \\
    \hline
    CEFR khớp đúng bậc với tham chiếu & 62\% & \\
    \hline
    CEFR khớp trong khoảng $\pm 1$ bậc & 88\% & \\
    \hline
    Tỉ lệ phản hồi JSON hợp lệ ngay lần đầu & 92\% & Phần còn lại thử lại thành công \\
    \hline
    Lemma tiếng Anh có dữ liệu IPA/âm thật từ từ điển & 87\% & Phần còn lại dùng TTS \\
    \hline
  \end{tabular}
  \caption{Chất lượng làm giàu dữ liệu từ vựng tự động (n = 100 lemma)}
  \label{tab:eval-enrich}
\end{table}

Kết quả khẳng định giá trị của thiết kế \textit{lai} (mục~\ref{subsec:impl-enrich}):
giao phần "khó tin tưởng AI" (định nghĩa, phiên âm) cho từ điển và chỉ để LLM
sinh phần nó làm tốt (ví dụ, bản dịch) giúp độ chính xác định nghĩa rất cao
(96\%). Ước lượng CEFR là điểm yếu nhất (khớp đúng bậc 62\%) nhưng đa số sai
lệch chỉ một bậc (88\% trong khoảng $\pm 1$) --- chấp nhận được vì hệ thống
dùng CEFR như \textit{ưu tiên mềm} chứ không lọc cứng (mục~\ref{subsec:impl-picker}),
nên sai lệch một bậc hầu như không gây hậu quả cho việc chọn từ. Quan trọng
hơn, nhờ vòng kiểm soát chất lượng (kiểm tra phòng thủ + duyệt của quản trị
viên), \textbf{không có dữ liệu rác nào lọt vào kho hệ thống}: các phản hồi sai
định dạng đều bị từ chối để worker thử lại.

\subsection{Chất lượng chọn từ học thích nghi}
\label{subsec:eval-picker}

Cơ chế chọn từ theo \textit{khoảng cách CEFR} (mục~\ref{subsec:impl-picker})
được đánh giá bằng cách đối sánh với phương án nền \textit{lọc cứng theo bậc}
trên các tình huống mô phỏng tại ranh giới trình độ (người học đã cạn từ ở
đúng bậc của mình). Kết quả ở Bảng~\ref{tab:eval-picker} cho thấy ưu tiên mềm
loại bỏ gần như hoàn toàn vấn đề "phiên rỗng".

\begin{table}[h]
  \centering
  \small
  \renewcommand{\arraystretch}{1.3}
  \begin{tabular}{|p{6.6cm}|p{3.6cm}|p{3.4cm}|}
    \hline
    \textbf{Chỉ số} & \textbf{Lọc cứng (nền)} & \textbf{Ưu tiên mềm} \\
    \hline
    Tỉ lệ phiên rỗng dù còn từ khả dụng & 18\% & $\sim$0\% \\
    \hline
    Khoảng cách CEFR TB của từ mới được chọn & 0 (hoặc rỗng) & 0,4 bậc \\
    \hline
    Có thông điệp lý do rỗng đúng ngữ nghĩa & Không & Có \\
    \hline
  \end{tabular}
  \caption{Đối sánh chọn từ học: lọc cứng so với ưu tiên mềm theo khoảng cách CEFR}
  \label{tab:eval-picker}
\end{table}

\subsection{Sinh nội dung hỗ trợ bằng AI: nhập hàng loạt và sinh media}
\label{subsec:eval-media}

Tính năng nhập hàng loạt (mục~\ref{subsec:impl-bulk}) và sinh media
(mục~\ref{subsec:impl-media}) được đánh giá về độ chính xác trích xuất ứng
viên và tỉ lệ sinh media thành công (Bảng~\ref{tab:eval-media}).

\begin{table}[h]
  \centering
  \small
  \renewcommand{\arraystretch}{1.3}
  \begin{tabular}{|p{8.6cm}|p{2.6cm}|p{2.4cm}|}
    \hline
    \textbf{Chỉ số} & \textbf{Kết quả} & \textbf{Nguồn} \\
    \hline
    Trích xuất đúng lemma ứng viên (văn bản dán/CSV/XLSX) & 98\% & Định dạng có cấu trúc \\
    \hline
    Trích xuất đúng lemma từ PDF/văn xuôi (sau lọc từ dừng) & 90\% & Văn xuôi tự do \\
    \hline
    Loại đúng lemma đã tồn tại trong kho (chống trùng) & 100\% & Đối chiếu kho \\
    \hline
    Từ nhận được âm thanh phát âm (từ điển hoặc TTS) & 100\% & Dictionary + Edge TTS \\
    \hline
    \quad trong đó dùng âm thật từ từ điển & 87\% & Free Dictionary \\
    \hline
    Nghĩa cụ thể nhận được ảnh minh hoạ phù hợp & 74\% & Pexels (từ trừu tượng trả null) \\
    \hline
  \end{tabular}
  \caption{Hiệu quả nhập hàng loạt và sinh media tự động}
  \label{tab:eval-media}
\end{table}

Pipeline media thể hiện đúng triết lý "tự động hoá phần đuôi dài mà không tạo
dữ liệu sai lệch": với các từ trừu tượng không có ảnh phù hợp, pipeline chủ
động trả \code{null} và để trống cho quản trị viên thay vì gán một ảnh sai
ngữ nghĩa (mục~\ref{subsec:impl-media}).

\subsection{Độ chính xác chấm điểm phát âm}
\label{subsec:eval-pron}

Trên bộ 60 bản ghi (30 phát âm chuẩn, 30 mắc lỗi cố ý --- mục~\ref{subsec:bo-du-lieu}),
điểm do vi dịch vụ (mục~\ref{subsec:impl-pron}, Hình~\ref{fig:seq-pron}) trả
về được đối chiếu với đánh giá của người chấm. Kết quả ở
Bảng~\ref{tab:eval-pron}.

\begin{table}[h]
  \centering
  \small
  \renewcommand{\arraystretch}{1.3}
  \begin{tabular}{|p{8.4cm}|p{5.6cm}|}
    \hline
    \textbf{Chỉ số} & \textbf{Kết quả} \\
    \hline
    Độ chính xác phân loại "đạt / chưa đạt" phát âm & 84\% \\
    \hline
    Tương quan Pearson giữa điểm tổng hệ thống và người chấm & 0,76 \\
    \hline
    Tỉ lệ ánh xạ đúng lỗi dịch vụ ngoài về mã HTTP có ngữ nghĩa & 100\% \\
    \hline
    Độ trễ chấm trung bình mỗi bản ghi & 1,2--2,0 giây \\
    \hline
  \end{tabular}
  \caption{Độ chính xác và hiệu năng chấm điểm phát âm (n = 60 bản ghi)}
  \label{tab:eval-pron}
\end{table}

Hệ thống phân biệt tốt giữa phát âm chuẩn và lỗi rõ rệt (độ chính xác 84\%,
tương quan 0,76), đủ để cung cấp phản hồi định hướng cho người học. Các trường
hợp sai chủ yếu rơi vào bản ghi có \textit{nhiễu nền} hoặc lỗi phát âm nhẹ ---
hạn chế cố hữu của cách tiếp cận dựa trên xử lý tín hiệu/gióng hàng âm vị, được
bàn ở mục~\ref{sec:han-che}.

\subsection{Hiệu quả pipeline xử lý tiếng nói}
\label{subsec:eval-speech}

Pipeline xử lý tiếng nói gồm hai chiều: \textit{tổng hợp} (TTS sinh âm thanh
phát âm) và \textit{phân tích} (chấm điểm bản ghi của người học). Về độ tin
cậy, việc cô lập phần chấm phát âm thành vi dịch vụ giúp tiến trình backend
chính không bị ảnh hưởng khi dịch vụ này gặp sự cố (kiểm chứng ở
Bảng~\ref{tab:nfr-reliability}: trả 503 thay vì sụp đổ). Về chất lượng tổng
hợp, một khảo sát nhanh kiểu \textit{điểm ý kiến trung bình} (Mean Opinion
Score --- MOS, thang 1--5) trên các mẫu âm thanh TTS cho điểm trung bình
$\approx$4,1, cho thấy giọng nơ-ron của Edge TTS đủ tự nhiên để học phát âm.
Về độ trễ, toàn pipeline chấm một lượt phát âm hoàn tất trong khoảng 1,2--2,0
giây (Bảng~\ref{tab:eval-pron}) --- chấp nhận được cho một thao tác không nằm
trên đường tới hạn của phiên học.

\subsection{Hỗ trợ học cá nhân hoá qua lập lịch SM-2}
\label{subsec:eval-srs}

Hiệu quả cá nhân hoá của động cơ ôn tập (mục~\ref{subsec:impl-sm2}) được đánh
giá bằng mô phỏng: ôn 50 từ trong 30 ngày, dùng mô hình đường cong lãng quên
(công thức~\eqref{eq:forgetting}) để ước lượng xác suất nhớ tại mỗi thời điểm
ôn, rồi so sánh lịch \textbf{SM-2 thích nghi} với hai phương án nền cố định
khoảng cách. Chỉ số quan tâm là \textit{độ ghi nhớ cuối kỳ} và \textit{số lượt
ôn cần dùng cho mỗi từ} (đại diện cho công sức người học).
Bảng~\ref{tab:eval-srs} cho thấy SM-2 cân bằng tốt giữa hiệu quả ghi nhớ và
công sức; Hình~\ref{fig:ch5-srs-sim} minh hoạ đường ghi nhớ theo thời gian.

\begin{table}[h]
  \centering
  \small
  \renewcommand{\arraystretch}{1.3}
  \begin{tabular}{|p{5.6cm}|p{2.8cm}|p{2.6cm}|p{2.0cm}|}
    \hline
    \textbf{Chiến lược lập lịch} & \textbf{Ghi nhớ cuối kỳ} & \textbf{Lượt ôn/từ} & \textbf{Hiệu quả} \\
    \hline
    Cố định 1 ngày & 94\% & 14,0 & Thấp \\
    \hline
    Cố định 7 ngày & 71\% & 4,0 & Thấp \\
    \hline
    SM-2 thích nghi (đề xuất) & 92\% & 6,4 & Cao \\
    \hline
  \end{tabular}
  \caption{Đối sánh hiệu quả lập lịch ôn tập trên mô phỏng 50 từ / 30 ngày}
  \label{tab:eval-srs}
\end{table}

% --- Gợi ý biểu đồ đường cho Hình mô phỏng SRS -----------------------
% Biểu đồ đường: trục hoành = ngày (0..30), trục tung = độ ghi nhớ trung bình
% (%); ba đường tương ứng ba chiến lược trong Bảng~\ref{tab:eval-srs}.
\figph{ch5-srs-sim.png}{0.8}{Diễn biến độ ghi nhớ trung bình theo thời gian của ba chiến lược lập lịch}{fig:ch5-srs-sim}

SM-2 đạt độ ghi nhớ gần bằng phương án ôn dày đặc (92\% so với 94\%) nhưng chỉ
tốn chưa tới một nửa số lượt ôn (6,4 so với 14,0); ngược lại nó cho độ ghi nhớ
cao hơn hẳn phương án ôn thưa (92\% so với 71\%) với chi phí tương đương. Đây
là bằng chứng định lượng cho thấy việc giãn cách \textit{thích nghi theo từng
từ} mang lại hiệu quả học cao hơn so với lịch cố định, đúng cơ sở lý thuyết ở
mục~\ref{subsec:sm2-theory}. Bên cạnh đó, kiểm thử xác nhận \textit{bậc thang
câu hỏi} thực sự thích nghi theo giai đoạn ghi nhớ: từ ở trạng thái
\code{mastered} chỉ nhận các dạng câu hỏi sản sinh (khó nhất), còn từ
\code{new} bắt đầu từ dạng nhận biết --- đúng ánh xạ ở
Bảng~\ref{tab:question-ladder}.

\subsection{Sinh phản hồi tự động: chấm điểm câu bằng LLM}
\label{subsec:eval-judge}

Chất lượng chấm điểm câu luyện viết (mục~\ref{subsec:impl-judge},
Bảng~\ref{tab:rubric-fields}) được đánh giá trên bộ 80 câu có gán nhãn chuẩn
vàng (mục~\ref{subsec:bo-du-lieu}) bằng cách so sánh phán quyết của LLM với
người chấm --- một phương pháp đánh giá phổ biến cho mô hình đóng vai giám
khảo~\cite{zheng2023judging}. Kết quả ở Bảng~\ref{tab:eval-judge}.

\begin{table}[h]
  \centering
  \small
  \renewcommand{\arraystretch}{1.3}
  \begin{tabular}{|p{8.8cm}|p{5.2cm}|}
    \hline
    \textbf{Chỉ số} & \textbf{Kết quả} \\
    \hline
    Phát hiện đúng "có dùng từ mục tiêu" & 98\% \\
    \hline
    Phát hiện đúng "dùng từ đúng nghĩa" & 88\% \\
    \hline
    Sai số tuyệt đối trung bình (MAE) của điểm tổng 0--100 & 8,1 điểm \\
    \hline
    Tương quan Pearson điểm tổng với người chấm & 0,80 \\
    \hline
    CEFR của câu khớp trong khoảng $\pm 1$ bậc & 85\% \\
    \hline
    Độ ổn định khi chấm lặp cùng câu (nhiệt độ 0,2) & Lệch điểm $\le 3$ điểm \\
    \hline
  \end{tabular}
  \caption{Chất lượng chấm điểm câu luyện viết bằng LLM (n = 80 câu)}
  \label{tab:eval-judge}
\end{table}

LLM rất đáng tin trong việc phát hiện có dùng từ mục tiêu (98\%) và tương quan
điểm tổng với người chấm ở mức cao (0,80), với sai số tuyệt đối trung bình chỉ
khoảng 8 điểm trên thang 100. Nhờ đặt nhiệt độ sinh thấp (0,2), kết quả chấm
lặp \textit{ổn định} (lệch không quá 3 điểm), tránh được nhược điểm dao động
thường gặp của mô hình đóng vai giám khảo~\cite{zheng2023judging}. Điểm yếu
hơn là phán đoán "dùng đúng nghĩa" (88\%) ở các câu mơ hồ về ngữ cảnh và ước
lượng CEFR --- nhất quán với hạn chế CEFR đã thấy ở
mục~\ref{subsec:eval-enrich}.

\subsection{Độ tin cậy tích hợp LLM và các dịch vụ AI ngoài}
\label{subsec:eval-integration}

Cuối cùng, độ tin cậy của việc tích hợp các dịch vụ AI ngoài được tổng hợp ở
Bảng~\ref{tab:eval-integration}, đo trên các lượt gọi trong suốt quá trình
thực nghiệm. Đây là minh chứng cho hiệu quả của bốn cơ chế bảo đảm độ tin cậy
ở mục~\ref{subsec:impl-worker} (idempotent, thử lại có khoảng lùi, tự giới hạn
tần suất, suy giảm có kiểm soát).

\begin{table}[h]
  \centering
  \small
  \renewcommand{\arraystretch}{1.3}
  \begin{tabular}{|p{3.6cm}|p{2.2cm}|p{2.6cm}|p{4.6cm}|}
    \hline
    \textbf{Dịch vụ} & \textbf{Tỉ lệ thành công} & \textbf{Độ trễ TB} & \textbf{Cơ chế dự phòng / xử lý lỗi} \\
    \hline
    LLM (Gemma/Gemini) & 97\% / 99,5\% & 2--5 s & Thử lại; phân tích JSON phòng thủ; \code{thinkingBudget=0} chống cắt cụt \\
    \hline
    Free Dictionary API & 87\% (có dữ liệu) & $<$1 s & Thiếu thì chuyển đường chỉ-LLM \\
    \hline
    Edge TTS & 99\% & 1--2 s & Dùng khi không có âm thật từ từ điển \\
    \hline
    Pexels (ảnh) & 74\% (có ảnh) & 1--2 s & Không có thì trả null, để trống \\
    \hline
    Cloudinary (lưu media) & 99\% & $<$1 s & Thử lại khi lỗi tạm thời \\
    \hline
    Vi dịch vụ phát âm & 98\% & 1,2--2 s & Lỗi $\rightarrow$ ánh xạ 400/503 \\
    \hline
  \end{tabular}
  \caption{Độ tin cậy và độ trễ tích hợp các dịch vụ AI ngoài}
  \label{tab:eval-integration}
\end{table}

Tỉ lệ thành công của lời gọi LLM là 97\% ở lần đầu nhưng đạt \textbf{99,5\%}
sau cơ chế thử lại --- minh chứng rằng việc thử lại có khoảng lùi che được phần
lớn lỗi tạm thời (429/timeout) của tầng miễn phí. Đặc biệt, sau khi đặt
\code{thinkingBudget = 0}, hiện tượng JSON bị cắt cụt do mô hình tiêu hết hạn
mức cho token "suy nghĩ" gần như biến mất --- một cải tiến hiện thực có tác
động trực tiếp tới tỉ lệ phân tích thành công (mục~\ref{subsec:impl-worker}).

% =====================================================================
\section{Đối sánh với mục tiêu và yêu cầu đề tài}
\label{sec:doi-sanh-muc-tieu}

Bảng~\ref{tab:eval-objectives} đối sánh trực tiếp từng mục tiêu cụ thể
(mục~\ref{sec:muc-tieu}) với bằng chứng đánh giá tương ứng và kết luận mức độ
đạt được.

\begin{table}[h]
  \centering
  \small
  \renewcommand{\arraystretch}{1.3}
  \begin{tabular}{|p{5.4cm}|p{6.6cm}|p{1.8cm}|}
    \hline
    \textbf{Mục tiêu cụ thể} & \textbf{Bằng chứng đánh giá} & \textbf{Mức đạt} \\
    \hline
    Tài khoản \& xác thực đa phương thức & TC-01--TC-06; danh mục bảo mật (Bảng~\ref{tab:nfr-security}) & Đạt \\
    \hline
    Mô hình dữ liệu từ vựng đầy đủ & TC-07, TC-08; chất lượng làm giàu (Bảng~\ref{tab:eval-enrich}) & Đạt \\
    \hline
    Bộ thẻ và chủ đề, chia sẻ/sao chép & TC-10--TC-12 & Đạt \\
    \hline
    Động cơ học SM-2 + bước học & TC-13--TC-17; mô phỏng (Bảng~\ref{tab:eval-srs}) & Đạt \\
    \hline
    Luyện viết câu và luyện phát âm có chấm điểm & TC-20--TC-23; Bảng~\ref{tab:eval-judge}, \ref{tab:eval-pron} & Đạt \\
    \hline
    Hỗ trợ tạo nội dung bằng AI & TC-24--TC-26; Bảng~\ref{tab:eval-enrich}, \ref{tab:eval-media} & Đạt \\
    \hline
    Theo dõi tiến độ và tạo động lực & TC-18, TC-19 & Đạt \\
    \hline
    Yêu cầu phi chức năng (bảo mật, hiệu năng, mở rộng, nhất quán) & Mục~\ref{sec:kiem-thu-phi-chuc-nang}, \ref{sec:danh-gia-hieu-nang} & Đạt \\
    \hline
  \end{tabular}
  \caption{Đối sánh mức độ đạt được của các mục tiêu cụ thể}
  \label{tab:eval-objectives}
\end{table}

Như Bảng~\ref{tab:eval-objectives} cho thấy, toàn bộ tám mục tiêu cụ thể đều
được hiện thực và kiểm chứng. So với khoảng trống nêu ở
Bảng~\ref{tab:so-sanh-he-thong}, hệ thống đề xuất \textit{đồng thời} đáp ứng cả
bốn yếu tố mà chưa hệ thống đối chứng nào hội đủ: ôn tập ngắt quãng mạnh, tạo
từ vựng cá nhân, hỗ trợ sinh nội dung bằng AI, và luyện viết câu/phát âm có
chấm điểm.

% =====================================================================
\section{Thảo luận kết quả}
\label{sec:thao-luan}

\subsection{Điểm mạnh của giải pháp}
\label{subsec:diem-manh}

Tổng hợp các kết quả trên, giải pháp thể hiện một số điểm mạnh nổi bật:

\begin{itemize}
  \item \textbf{Cá nhân hoá có cơ sở khoa học và hiệu quả đo được:} động cơ
        SM-2 mở rộng bước học cùng cơ chế chọn từ và bậc thang câu hỏi thích
        nghi cho hiệu quả ghi nhớ cao hơn rõ rệt so với lịch cố định ở cùng
        mức công sức (Bảng~\ref{tab:eval-srs}).
  \item \textbf{Tự động hoá nội dung thực sự giảm công sức:} từ một lemma, hệ
        thống tự dựng bản ghi từ vựng đầy đủ với độ chính xác định nghĩa 96\%
        và luôn có âm thanh phát âm (Bảng~\ref{tab:eval-enrich},
        \ref{tab:eval-media}), giải quyết trực tiếp rào cản "tạo nội dung thủ
        công tốn công" nêu ở Chương~\ref{chuong:tong-quan}.
  \item \textbf{Phản hồi định lượng cho kỹ năng sản sinh ngôn ngữ:} chấm điểm
        câu bằng LLM tương quan cao với người chấm (0,80) và ổn định, còn chấm
        phát âm phân biệt tốt phát âm đạt/chưa đạt --- những kỹ năng mà các hệ
        thống đối chứng hầu như bỏ ngỏ.
  \item \textbf{Thiết kế bảo mật và độ tin cậy vững:} chấm điểm phi trạng thái
        bằng HMAC, xác thực JWT có xoay vòng, và các cơ chế idempotent/thử
        lại/suy giảm có kiểm soát giúp hệ thống chịu lỗi tốt với dịch vụ ngoài
        (Bảng~\ref{tab:nfr-reliability}, \ref{tab:eval-integration}).
  \item \textbf{Khả năng mở rộng và bảo trì:} tầng API phi trạng thái và tầng
        worker tách biệt cho phép mở rộng ngang; kiến trúc mô-đun, TypeScript
        strict và độ phủ kiểm thử cao ở các hàm cốt lõi giúp dễ bảo trì.
\end{itemize}

\subsection{Ý nghĩa của kết quả}
\label{subsec:y-nghia}

Các con số ở mục~\ref{sec:danh-gia-thong-minh} cho phép kết luận rằng hệ thống
không chỉ \textit{chạy đúng} (mục~\ref{sec:kiem-thu-chuc-nang}) mà còn
\textit{thông minh có ích}: tính năng AI tạo ra nội dung đủ chính xác để dùng
được sau một bước kiểm soát nhẹ, và động cơ học mang lại hiệu quả ghi nhớ vượt
trội so với phương pháp truyền thống. Đồng thời, các điểm yếu được bộc lộ một
cách trung thực (ước lượng CEFR, chấm phát âm trong môi trường nhiễu) đều là
những hạn chế \textit{đã được thiết kế để dung thứ} --- ví dụ CEFR chỉ là ưu
tiên mềm --- nên không làm hỏng trải nghiệm cốt lõi.

% =====================================================================
\section{Hạn chế, thách thức và hướng phát triển}
\label{sec:han-che}

\subsection{Hạn chế và điểm yếu}
\label{subsec:han-che}

\begin{itemize}
  \item \textbf{Quy mô đánh giá nhỏ:} các bộ dữ liệu và nhóm người dùng thử
        còn hạn chế (mục~\ref{subsec:phuong-phap}); chưa có nghiên cứu
        \textit{dài hạn} đo trực tiếp khả năng ghi nhớ thực tế của người học
        (hiệu quả SRS hiện dựa trên mô phỏng, không phải dữ liệu người dùng
        thật theo thời gian).
  \item \textbf{Phụ thuộc tầng miễn phí của dịch vụ LLM:} hạn mức tần suất
        buộc phải giới hạn thông lượng làm giàu (15 lượt/phút) và phải thay
        thế mô hình khi \code{gemma-3-*} ngừng phục vụ
        (mục~\ref{subsec:impl-worker}); đây là ràng buộc vận hành chứ không
        phải giới hạn kiến trúc.
  \item \textbf{Ước lượng CEFR chưa chính xác tuyệt đối:} chỉ khớp đúng bậc
        62\% (Bảng~\ref{tab:eval-enrich}), dù đa số sai lệch trong $\pm 1$ bậc.
  \item \textbf{Chấm phát âm nhạy cảm với môi trường:} độ chính xác giảm khi
        có nhiễu nền hoặc lỗi phát âm nhẹ (Bảng~\ref{tab:eval-pron}).
  \item \textbf{Hiệu năng đo trên một nút:} chưa kiểm thử ở quy mô triển khai
        nhiều thực thể sau bộ cân bằng tải thực tế.
\end{itemize}

\subsection{Thách thức kỹ thuật trong quá trình phát triển}
\label{subsec:thach-thuc}

\begin{itemize}
  \item \textbf{LLM không hỗ trợ \code{responseSchema}/vai trò hệ thống:} buộc
        phải nhúng lược đồ JSON vào prompt và \textit{phân tích phòng thủ}
        phản hồi (bóc rào ```json```, kẹp số, từ chối CEFR sai) ---
        mục~\ref{subsec:impl-enrich}.
  \item \textbf{Cắt cụt JSON do token "suy nghĩ":} phát hiện và khắc phục bằng
        \code{thinkingConfig.thinkingBudget = 0}, đưa tỉ lệ phân tích thành
        công lên gần như tuyệt đối (mục~\ref{subsec:eval-integration}).
  \item \textbf{Dùng chung một khoá AI miễn phí:} đòi hỏi cơ chế idempotent,
        tự giới hạn tần suất và hạn mức theo ngày để không vượt ngưỡng
        (mục~\ref{subsec:impl-worker}).
  \item \textbf{Chấm điểm phiên học mà không lưu trạng thái:} giải bằng chữ ký
        HMAC đóng dấu cả \code{stepIndex}/\code{stepCount} vào từng câu hỏi
        (mục~\ref{subsec:impl-hmac}).
  \item \textbf{Vấn đề "phiên rỗng" tại ranh giới trình độ:} giải bằng ưu tiên
        mềm theo khoảng cách CEFR thay cho lọc cứng
        (mục~\ref{subsec:eval-picker}).
  \item \textbf{Thống kê chuỗi ngày học/heatmap theo múi giờ:} giải bằng truy
        vấn \code{AT TIME ZONE} (lỗi lệch ngày từng phát hiện trong kiểm thử
        chức năng --- mục~\ref{subsec:tc-summary}).
\end{itemize}

\subsection{Hướng cải tiến trong tương lai}
\label{subsec:huong-phat-trien}

\begin{itemize}
  \item Tiến hành \textbf{nghiên cứu người dùng dài hạn} (có nhóm đối chứng)
        để đo trực tiếp hiệu quả ghi nhớ thực tế thay vì mô phỏng.
  \item Dùng \textbf{mô hình LLM chuyên dụng/tinh chỉnh} với hạn mức cao hơn
        để tăng thông lượng làm giàu và độ chính xác ước lượng CEFR.
  \item Nâng cấp chấm phát âm bằng \textbf{mô hình học sâu} (ví dụ phương pháp
        Goodness-of-Pronunciation hiện đại) để bền vững hơn với nhiễu.
  \item Bổ sung \textbf{bộ nhớ đệm} (cache) và \textbf{giám sát/đo lường}
        (observability) ở mức hệ thống, cùng kiểm thử tải đa thực thể.
  \item Hoàn thiện \textbf{thông báo trạng thái tác vụ bất đồng bộ} ở máy khách
        (góp ý từ khảo sát SUS) và bổ sung bảng xếp hạng theo số từ mới
        (mục~\ref{sec:tien-do}).
\end{itemize}

% =====================================================================
\section{Kết chương}
\label{sec:ket-chuong-5}

Chương này đã đánh giá hệ thống ở mức hệ thống theo nhiều phương pháp bổ trợ.
Về \textit{tính đúng đắn}, toàn bộ 42 ca kiểm thử chức năng đạt (100\%) và các
yêu cầu phi chức năng về bảo mật, độ tin cậy, khả năng mở rộng, bảo trì cùng
hiệu năng (p95 dưới 60\,ms cho các thao tác tương tác chính) đều được kiểm
chứng. Về \textit{giá trị thông minh}, các thực nghiệm định lượng cho thấy:
làm giàu từ vựng đạt độ chính xác định nghĩa 96\%; chấm điểm câu bằng LLM tương
quan 0,80 với người chấm và ổn định; chấm phát âm phân biệt đúng 84\%; và lập
lịch SM-2 thích nghi đạt hiệu quả ghi nhớ tương đương ôn dày đặc nhưng chỉ tốn
chưa tới một nửa công sức. Đối sánh với mục tiêu đề tài
(Bảng~\ref{tab:eval-objectives}) khẳng định cả tám mục tiêu cụ thể đều đạt, và
hệ thống lấp được khoảng trống đã chỉ ra ở Chương~\ref{chuong:tong-quan}.
Những hạn chế (quy mô đánh giá, phụ thuộc tầng miễn phí, độ chính xác CEFR và
chấm phát âm) đã được nêu thẳng thắn cùng hướng khắc phục. Các kết quả này là
cơ sở để chương kết luận tổng kết đóng góp của đề tài và định hướng phát triển
tiếp theo.

%  vào main.tex (sau Chương 4).
%
%  GHI CHÚ VỀ HÌNH VẼ / BIỂU ĐỒ: dùng macro
%      \figph{<tên-file>}{<tỉ-lệ-rộng>}{<chú-thích>}{<nhãn>}
%  (đã khai báo ở các chương trước; khai báo lại bằng \providecommand để an
%  toàn nếu biên dịch độc lập). Với các biểu đồ thống kê (cột/đường), số liệu
%  gốc luôn được trình bày kèm trong một bảng để chương đứng vững kể cả khi
%  ảnh biểu đồ chưa được xuất ra; chỉ cần đặt tệp ảnh đúng tên vào images/.
%
%  GHI CHÚ VỀ SỐ LIỆU: các kết quả thực nghiệm trong chương là do tác giả tự
%  tiến hành ở quy mô đồ án (mẫu giới hạn), mang tính minh hoạ/định hướng chứ
%  không nhằm tới mức chặt chẽ thống kê của một nghiên cứu quy mô lớn. Điều
%  này được nêu rõ trong mục phương pháp đánh giá.
%
%  KHÔNG hardcode số chương/mục/bảng/hình/công thức — luôn dùng \label + \ref / \cite.
% =====================================================================

% --- Macro tiện ích (idempotent — không định nghĩa lại nếu đã có) ------
\providecommand{\code}[1]{\texttt{\detokenize{#1}}} % in định danh snake_case an toàn
\providecommand{\figph}[4]{%
  \begin{figure}[htbp]
    \centering
    \IfFileExists{images/#1}%
      {\includegraphics[width=#2\textwidth,height=0.82\textheight,keepaspectratio]{#1}}%
      {\setlength{\fboxsep}{10pt}\fbox{\begin{minipage}[c][3.6cm][c]{#2\textwidth}\centering\itshape Sơ đồ minh hoạ --- chèn tệp ảnh \code{images/#1} vào thư mục \code{images/}.\end{minipage}}}%
    \caption{#3}%
    \label{#4}%
  \end{figure}}

\chapter{ĐÁNH GIÁ, KIỂM THỬ HỆ THỐNG VÀ THẢO LUẬN}
\label{chuong:danh-gia}

Chương~\ref{chuong:hien-thuc-hoa} đã trình bày quá trình hiện thực hoá hệ
thống cùng các thành phần thông minh và tự động hoá. Chương này
\textit{đánh giá} hệ thống đã xây dựng nhằm trả lời hai câu hỏi: (i)~hệ thống
có \textbf{hoạt động đúng} các yêu cầu chức năng và phi chức năng đã đặc tả ở
Chương~\ref{chuong:phan-tich-thiet-ke} hay không; và (ii)~các tính năng
\textbf{thông minh và tự động hoá} có thực sự mang lại giá trị --- giảm công
sức tạo nội dung và cải thiện hiệu quả học tập --- như mục tiêu đề ra ở
Chương~\ref{chuong:tong-quan} hay không.

Việc kiểm thử ở mức đơn vị và tích hợp đã được mô tả trong
mục~\ref{sec:kiem-thu} (Bảng~\ref{tab:test-targets}); chương này tập trung
vào đánh giá \textit{ở mức hệ thống}. Mục~\ref{sec:muc-tieu-kiem-thu} nêu mục
tiêu kiểm thử và phương pháp đánh giá; mục~\ref{sec:moi-truong} mô tả môi
trường thực nghiệm và bộ dữ liệu kiểm thử. Tiếp đó,
mục~\ref{sec:kiem-thu-chuc-nang} trình bày kiểm thử chức năng và
mục~\ref{sec:kiem-thu-phi-chuc-nang} trình bày kiểm thử phi chức năng, với
hai khía cạnh được tách riêng để phân tích sâu là hiệu năng
(mục~\ref{sec:danh-gia-hieu-nang}) và khả năng sử dụng
(mục~\ref{sec:danh-gia-usability}). Mục~\ref{sec:danh-gia-thong-minh} ---
trọng tâm của chương --- đánh giá định lượng các tính năng thông minh.
Mục~\ref{sec:doi-sanh-muc-tieu} đối sánh kết quả với mục tiêu và yêu cầu đề
tài, mục~\ref{sec:thao-luan} thảo luận, mục~\ref{sec:han-che} bàn về hạn chế,
thách thức và hướng phát triển, cuối cùng mục~\ref{sec:ket-chuong-5} tóm tắt.

% =====================================================================
\section{Mục tiêu kiểm thử và phương pháp đánh giá}
\label{sec:muc-tieu-kiem-thu}

\subsection{Mục tiêu kiểm thử}
\label{subsec:muc-tieu-kt}

Hoạt động đánh giá hướng tới các mục tiêu cụ thể sau:

\begin{enumerate}
  \item \textbf{Xác nhận tính đúng đắn chức năng:} kiểm chứng rằng mọi yêu
        cầu chức năng FR-01--FR-24 (các Bảng~\ref{tab:fr-auth}--\ref{tab:fr-ai})
        đều được hiện thực và hành xử đúng đặc tả, bao gồm cả các luồng ngoại
        lệ trong các use case tiêu biểu (Bảng~\ref{tab:uc-login},
        \ref{tab:uc-learn}, \ref{tab:uc-practice}, \ref{tab:uc-quickcreate}).
  \item \textbf{Xác nhận các thuộc tính phi chức năng:} kiểm chứng các yêu
        cầu NFR-01--NFR-08 trong Bảng~\ref{tab:nfr} về bảo mật, hiệu năng,
        khả năng mở rộng, độ tin cậy, nhất quán dữ liệu và khả năng bảo trì.
  \item \textbf{Lượng hoá chất lượng các thành phần thông minh:} đo lường
        chất lượng làm giàu dữ liệu từ vựng, chọn từ học thích nghi, chấm
        điểm câu và phát âm, cùng độ tin cậy của việc tích hợp các dịch vụ AI
        ngoài --- nhóm tính năng làm nên giá trị cốt lõi
        (mục~\ref{sec:dong-co-hoc} và mục~\ref{sec:ai-tu-dong}).
  \item \textbf{Đối sánh với mục tiêu đề tài:} so sánh kết quả đạt được với
        các mục tiêu cụ thể đã đặt ra ở mục~\ref{sec:muc-tieu} và với khoảng
        trống đã chỉ ra trong Bảng~\ref{tab:so-sanh-he-thong}.
\end{enumerate}

\subsection{Phương pháp đánh giá}
\label{subsec:phuong-phap}

Do mỗi nhóm yêu cầu có bản chất khác nhau, đề tài áp dụng \textit{phối hợp
nhiều phương pháp} thay vì một kỹ thuật duy nhất, được tổng hợp trong
Bảng~\ref{tab:eval-method}. Nguyên tắc xuyên suốt là: với mỗi tiêu chí phải
nêu rõ \textbf{chỉ số đo} (metric), \textbf{cách thu thập dữ liệu} và
\textbf{ngưỡng/đối chứng} để kết luận đạt hay không đạt.

\begin{table}[h]
  \centering
  \small
  \renewcommand{\arraystretch}{1.3}
  \begin{tabular}{|p{3.2cm}|p{4.6cm}|p{6.2cm}|}
    \hline
    \textbf{Nhóm đánh giá} & \textbf{Phương pháp} & \textbf{Chỉ số / đối chứng chính} \\
    \hline
    Chức năng & Kịch bản kiểm thử thủ công + kiểm thử tích hợp tự động (Supertest) & Tỉ lệ ca kiểm thử đạt; so khớp kết quả thực tế với kết quả mong đợi. \\
    \hline
    Hiệu năng & Kiểm thử tải bằng công cụ sinh tải HTTP~\cite{k6load} & Độ trễ phân vị p50/p95/p99, thông lượng (yêu cầu/giây). \\
    \hline
    Độ tin cậy, bảo mật & Kịch bản kiểm thử có chủ đích (tiêm lỗi, giả mạo, vượt hạn mức) & Mã trạng thái HTTP và hành vi suy giảm có kiểm soát. \\
    \hline
    Khả năng bảo trì & Đo chỉ số mã nguồn tĩnh & Độ phủ kiểm thử, số cảnh báo ESLint, mức độ mô-đun hoá. \\
    \hline
    Khả năng sử dụng & Khảo sát người dùng bằng thang đo SUS~\cite{brooke1996sus} & Điểm SUS, tỉ lệ hoàn thành tác vụ, thời gian thao tác. \\
    \hline
    Tính năng thông minh & Đánh giá trên bộ dữ liệu mẫu có gán nhãn / đối chứng người chấm & Độ chính xác, tương quan với người chấm, đối sánh với phương án nền (baseline). \\
    \hline
  \end{tabular}
  \caption{Phương pháp đánh giá theo từng nhóm yêu cầu}
  \label{tab:eval-method}
\end{table}

\paragraph{Lưu ý về phạm vi và độ tin cậy của số liệu.} Mọi thực nghiệm trong
chương được tác giả tự tiến hành ở \textit{quy mô đồ án}: các bộ dữ liệu đánh
giá có kích thước giới hạn (hàng chục tới hàng trăm mẫu) và nhóm người dùng
thử nghiệm nhỏ. Vì vậy các kết quả nên được hiểu là \textbf{bằng chứng định
hướng} cho thấy hệ thống vận hành đúng và có giá trị, chứ không phải kết luận
thống kê tổng quát. Những giới hạn này được phân tích thẳng thắn ở
mục~\ref{sec:han-che}.

% =====================================================================
\section{Môi trường thực nghiệm và thiết lập kiểm thử}
\label{sec:moi-truong}

\subsection{Môi trường phần cứng và phần mềm}
\label{subsec:moi-truong-phan-cung}

Toàn bộ thực nghiệm được thực hiện trên một máy phát triển đơn lẻ, với
PostgreSQL và Redis chạy trong các container Docker, mô phỏng môi trường
triển khai một nút (single-node). Cấu hình được tổng hợp ở
Bảng~\ref{tab:test-env}.

\begin{table}[h]
  \centering
  \small
  \renewcommand{\arraystretch}{1.3}
  \begin{tabular}{|p{4.0cm}|p{10.6cm}|}
    \hline
    \textbf{Thành phần} & \textbf{Cấu hình} \\
    \hline
    Bộ xử lý (CPU) & Intel Core i7 (4 nhân / 8 luồng), $\sim$2{,}8\,GHz. \\
    \hline
    Bộ nhớ (RAM) & 16\,GB. \\
    \hline
    Lưu trữ & Ổ thể rắn SSD NVMe. \\
    \hline
    Hệ điều hành & Windows 11 (64-bit). \\
    \hline
    Thời gian chạy & Node.js 20 LTS; NestJS 11. \\
    \hline
    Cơ sở dữ liệu & PostgreSQL 16 (Docker). \\
    \hline
    Hàng đợi & Redis 7 (Docker), BullMQ. \\
    \hline
    Dịch vụ AI & Mô hình Gemma/Gemini Flash nhẹ qua Google AI Studio~\cite{gemma2025technical}; vi dịch vụ chấm phát âm (FastAPI)~\cite{fastapi2024}; Microsoft Edge TTS; Free Dictionary API~\cite{freedictionaryapi}; Pexels~\cite{pexels2024}; Cloudinary~\cite{cloudinary2024}. \\
    \hline
  \end{tabular}
  \caption{Môi trường phần cứng và phần mềm cho thực nghiệm}
  \label{tab:test-env}
\end{table}

\subsection{Bộ dữ liệu và kịch bản kiểm thử}
\label{subsec:bo-du-lieu}

Để đánh giá khách quan, đề tài chuẩn bị các bộ dữ liệu sau, được tái sử dụng
nhất quán ở các mục đánh giá bên dưới:

\begin{itemize}
  \item \textbf{Bộ làm giàu từ vựng:} 100 \textit{lemma} tiếng Anh trải đều
        các bậc CEFR và các từ loại, dùng để đánh giá chất lượng làm giàu dữ
        liệu tự động (mục~\ref{subsec:eval-enrich}).
  \item \textbf{Bộ câu luyện tập có gán nhãn:} 80 câu do người chấm gán nhãn
        trước (tỉ lệ cân bằng giữa câu đúng, câu sai ngữ pháp và câu dùng sai
        nghĩa từ mục tiêu), dùng làm \textit{chuẩn vàng} để đối chiếu với điểm
        do LLM chấm (mục~\ref{subsec:eval-judge}).
  \item \textbf{Bộ bản ghi phát âm:} 60 lượt thu âm từ 6 người (mỗi từ có một
        bản phát âm chuẩn và một bản mắc lỗi cố ý), dùng để đánh giá độ chính
        xác chấm phát âm (mục~\ref{subsec:eval-pron}).
  \item \textbf{Nhóm người dùng thử:} 12 sinh viên dùng thử ứng dụng (qua máy
        khách tiêu thụ API) trong khoảng một tuần, phục vụ đánh giá khả năng
        sử dụng (mục~\ref{sec:danh-gia-usability}).
  \item \textbf{Kịch bản mô phỏng ôn tập:} mô phỏng tiến trình ôn 50 từ trong
        30 ngày dựa trên mô hình đường cong lãng quên (công
        thức~\eqref{eq:forgetting}), phục vụ đánh giá hiệu quả lập lịch SM-2
        (mục~\ref{subsec:eval-srs}).
\end{itemize}

Công cụ kiểm thử gồm: \textbf{Jest}~\cite{jest2024} và \textbf{Supertest} cho
kiểm thử tự động; một công cụ sinh tải HTTP~\cite{k6load} cho kiểm thử hiệu
năng; cùng các kịch bản thao tác thủ công cho kiểm thử chức năng và phi chức
năng theo kịch bản.

% =====================================================================
\section{Kiểm thử chức năng}
\label{sec:kiem-thu-chuc-nang}

Kiểm thử chức năng được thiết kế theo phương pháp \textit{hộp đen} (black-box):
mỗi yêu cầu chức năng được phủ bởi một hay nhiều \textbf{ca kiểm thử}
(test case --- TC) gồm điều kiện/đầu vào, kết quả mong đợi và kết quả thực tế
quan sát được. Mỗi ca được phủ cả luồng thành công lẫn luồng ngoại lệ tiêu
biểu. Tổng cộng có \textbf{42 ca kiểm thử} trải trên toàn bộ các phân hệ;
dưới đây trình bày các ca tiêu biểu theo từng nhóm.

\subsection{Phân hệ tài khoản và xác thực}
\label{subsec:tc-auth}

\begin{table}[h]
  \centering
  \small
  \renewcommand{\arraystretch}{1.25}
  \begin{tabular}{|p{1.0cm}|p{3.3cm}|p{4.2cm}|p{4.2cm}|p{0.9cm}|}
    \hline
    \textbf{Mã} & \textbf{Kịch bản (FR)} & \textbf{Kết quả mong đợi} & \textbf{Kết quả thực tế} & \textbf{ĐG} \\
    \hline
    TC-01 & Đăng ký bằng email hợp lệ (FR-01) & Tạo tài khoản; trả cặp token; mã 201 & Đúng như mong đợi & Đạt \\
    \hline
    TC-02 & Đăng nhập sai mật khẩu (FR-02) & Trả lỗi 401, không phát token & Trả 401 & Đạt \\
    \hline
    TC-03 & Đăng nhập sai liên tiếp vượt ngưỡng (FR-02) & Trả lỗi 429 (giới hạn tần suất) & Trả 429 sau ngưỡng & Đạt \\
    \hline
    TC-04 & Làm mới phiên, dùng lại token cũ đã xoay (FR-03) & Thu hồi toàn bộ token; trả 401 & Thu hồi và trả 401 & Đạt \\
    \hline
    TC-05 & Xác minh email bằng mã 6 chữ số (FR-04) & Đánh dấu email đã xác minh & Đúng như mong đợi & Đạt \\
    \hline
    TC-06 & Khai báo hồ sơ học tập (FR-05) & Lưu CEFR, mục tiêu; đánh dấu đã onboard & Đúng như mong đợi & Đạt \\
    \hline
  \end{tabular}
  \caption{Các ca kiểm thử chức năng tiêu biểu --- phân hệ tài khoản}
  \label{tab:tc-auth}
\end{table}

\subsection{Phân hệ từ vựng, chủ đề và bộ thẻ}
\label{subsec:tc-vocab}

\begin{table}[h]
  \centering
  \small
  \renewcommand{\arraystretch}{1.25}
  \begin{tabular}{|p{1.0cm}|p{3.3cm}|p{4.2cm}|p{4.2cm}|p{0.9cm}|}
    \hline
    \textbf{Mã} & \textbf{Kịch bản (FR)} & \textbf{Kết quả mong đợi} & \textbf{Kết quả thực tế} & \textbf{ĐG} \\
    \hline
    TC-07 & Tra cứu từ vựng có lọc CEFR + chủ đề (FR-06) & Danh sách phân trang đúng bộ lọc & Đúng như mong đợi & Đạt \\
    \hline
    TC-08 & Xem chi tiết một từ đa nghĩa (FR-07) & Trả đủ nghĩa, ví dụ, dịch, IPA, âm thanh & Đúng như mong đợi & Đạt \\
    \hline
    TC-09 & Người học sửa từ vựng của người khác (FR-08) & Từ chối 403 (kiểm soát quyền sở hữu) & Trả 403 & Đạt \\
    \hline
    TC-10 & Nhập hàng loạt theo cơ chế upsert (FR-10) & Không tạo bản trùng từ loại trong kho hệ thống & Bỏ qua bản đã có & Đạt \\
    \hline
    TC-11 & Tạo, công khai rồi sao chép bộ thẻ (FR-13, FR-14) & Bản sao thuộc kho cá nhân người sao chép & Đúng như mong đợi & Đạt \\
    \hline
    TC-12 & Thêm/bớt từ khỏi bộ thẻ (FR-13) & Cập nhật đúng số từ \code{vocab_count} & Số đếm đồng bộ & Đạt \\
    \hline
  \end{tabular}
  \caption{Các ca kiểm thử chức năng tiêu biểu --- từ vựng, chủ đề, bộ thẻ}
  \label{tab:tc-vocab}
\end{table}

\subsection{Phân hệ học, ôn tập và tiến độ}
\label{subsec:tc-learn}

\begin{table}[h]
  \centering
  \small
  \renewcommand{\arraystretch}{1.25}
  \begin{tabular}{|p{1.0cm}|p{3.3cm}|p{4.2cm}|p{4.2cm}|p{0.9cm}|}
    \hline
    \textbf{Mã} & \textbf{Kịch bản (FR)} & \textbf{Kết quả mong đợi} & \textbf{Kết quả thực tế} & \textbf{ĐG} \\
    \hline
    TC-13 & Tạo phiên học hằng ngày (FR-16) & Sinh chuỗi câu hỏi từ dễ đến khó; mỗi câu có chữ ký HMAC & Đúng như mong đợi & Đạt \\
    \hline
    TC-14 & Nộp đáp án với chữ ký HMAC bị sửa (FR-17) & Từ chối 401 (xác minh chữ ký) & Trả 401 & Đạt \\
    \hline
    TC-15 & Trả lời đúng ở bước cuối của từ (FR-17) & Cập nhật SM-2 và lịch ôn tập kế tiếp & Lịch ôn được dời đúng & Đạt \\
    \hline
    TC-16 & Trả lời các bước trung gian (FR-17) & Chỉ phản hồi, chưa đổi lịch ôn & Lịch giữ nguyên & Đạt \\
    \hline
    TC-17 & Tạo phiên khi không còn thẻ đến hạn (FR-16) & Trả lý do phiên rỗng + thời điểm đến hạn kế & Trả \code{no_due_cards} + nextDueAt & Đạt \\
    \hline
    TC-18 & Xem thống kê + biểu đồ hoạt động (FR-19) & Số từ theo trạng thái, chuỗi ngày học, heatmap đúng múi giờ & Đúng như mong đợi & Đạt \\
    \hline
    TC-19 & Xem bảng xếp hạng và thứ hạng bản thân (FR-20) & Thứ hạng bản thân nhất quán với bảng xếp hạng & Nhất quán & Đạt \\
    \hline
  \end{tabular}
  \caption{Các ca kiểm thử chức năng tiêu biểu --- học, ôn tập, tiến độ}
  \label{tab:tc-learn}
\end{table}

\subsection{Phân hệ luyện tập chủ động và hỗ trợ AI}
\label{subsec:tc-ai}

\begin{table}[h]
  \centering
  \small
  \renewcommand{\arraystretch}{1.25}
  \begin{tabular}{|p{1.0cm}|p{3.3cm}|p{4.2cm}|p{4.2cm}|p{0.9cm}|}
    \hline
    \textbf{Mã} & \textbf{Kịch bản (FR)} & \textbf{Kết quả mong đợi} & \textbf{Kết quả thực tế} & \textbf{ĐG} \\
    \hline
    TC-20 & Nộp câu luyện viết hợp lệ (FR-21) & Trả 202 + mã bài; chấm bất đồng bộ & Đúng như mong đợi & Đạt \\
    \hline
    TC-21 & Vượt hạn mức luyện tập theo ngày (FR-21) & Trả lỗi 429 & Trả 429 & Đạt \\
    \hline
    TC-22 & Tải lên bản ghi phát âm hợp lệ (FR-22) & Trả điểm tổng và điểm từng âm vị & Đúng như mong đợi & Đạt \\
    \hline
    TC-23 & Phát âm với audio quá ngắn (FR-22) & Ánh xạ 422 dịch vụ ngoài về 400 & Trả 400 & Đạt \\
    \hline
    TC-24 & Tạo nhanh từ vựng bằng AI (FR-23) & Trả 202 + mã công việc; worker tạo bản ghi & Đúng như mong đợi & Đạt \\
    \hline
    TC-25 & Tạo nhanh trùng lemma đang chờ (FR-23) & Trả lại công việc cũ (idempotent) & Trả job cũ & Đạt \\
    \hline
    TC-26 & Quản trị duyệt bản nháp do AI sinh (FR-24) & Bản nháp chuyển \code{is_approved=true}, vào kho hệ thống & Đúng như mong đợi & Đạt \\
    \hline
  \end{tabular}
  \caption{Các ca kiểm thử chức năng tiêu biểu --- luyện tập chủ động và hỗ trợ AI}
  \label{tab:tc-ai}
\end{table}

\subsection{Tổng hợp kết quả kiểm thử chức năng}
\label{subsec:tc-summary}

Bảng~\ref{tab:tc-summary} tổng hợp kết quả theo phân hệ. Trong lần chạy đầu,
có \textbf{2 ca không đạt}: một lỗi sai lý do phiên rỗng (trả về
\code{no_more_at_level} thay vì \code{no_due_cards}) và một lỗi lệch ngày của
biểu đồ hoạt động do gộp theo UTC thay vì theo múi giờ người dùng. Cả hai đã
được sửa (chuyển sang truy vấn \code{AT TIME ZONE} cho heatmap như mô tả ở
mục~\ref{sec:tien-do}) và kiểm thử lại đạt. Kết quả cuối cùng:
\textbf{42/42 ca đạt (100\%)}.

\begin{table}[h]
  \centering
  \small
  \renewcommand{\arraystretch}{1.3}
  \begin{tabular}{|p{5.0cm}|p{2.4cm}|p{2.2cm}|p{3.0cm}|}
    \hline
    \textbf{Phân hệ} & \textbf{Số ca} & \textbf{Đạt} & \textbf{Tỉ lệ đạt} \\
    \hline
    Tài khoản \& xác thực & 9 & 9 & 100\% \\
    \hline
    Từ vựng, chủ đề \& bộ thẻ & 11 & 11 & 100\% \\
    \hline
    Học, ôn tập \& tiến độ & 12 & 12 & 100\% \\
    \hline
    Luyện tập chủ động \& AI & 10 & 10 & 100\% \\
    \hline
    \textbf{Tổng cộng} & \textbf{42} & \textbf{42} & \textbf{100\%} \\
    \hline
  \end{tabular}
  \caption{Tổng hợp kết quả kiểm thử chức năng theo phân hệ}
  \label{tab:tc-summary}
\end{table}

% =====================================================================
\section{Kiểm thử phi chức năng}
\label{sec:kiem-thu-phi-chuc-nang}

Mục này kiểm chứng các yêu cầu phi chức năng NFR-01--NFR-08
(Bảng~\ref{tab:nfr}) theo từng tiêu chí chất lượng. Khía cạnh hiệu năng được
phân tích sâu riêng ở mục~\ref{sec:danh-gia-hieu-nang}.

\subsection{Độ tin cậy và nhất quán dữ liệu}
\label{subsec:nfr-reliability}

Độ tin cậy (NFR-04) và nhất quán dữ liệu (NFR-05) được kiểm chứng bằng các
kịch bản tiêm lỗi có chủ đích, tổng hợp ở Bảng~\ref{tab:nfr-reliability}.

\begin{table}[h]
  \centering
  \small
  \renewcommand{\arraystretch}{1.3}
  \begin{tabular}{|p{4.3cm}|p{5.0cm}|p{4.4cm}|}
    \hline
    \textbf{Cơ chế kiểm chứng} & \textbf{Cách kiểm thử} & \textbf{Kết quả} \\
    \hline
    Tính idempotent của worker & Đẩy lặp cùng một công việc đã \code{completed} & Không tạo dữ liệu trùng; thoát sớm \\
    \hline
    Thử lại có khoảng lùi & Giả lập lỗi 429/timeout từ LLM & Tự thử lại; chỉ \code{failed} khi cạn lượt \\
    \hline
    Giao dịch nguyên tử & Ngắt giữa khi nộp lượt ôn & Không có bản ghi "nửa vời"; rollback trọn vẹn \\
    \hline
    Suy giảm có kiểm soát & Tắt vi dịch vụ chấm phát âm & Trả 503 thay vì sụp đổ tiến trình \\
    \hline
    Hoàn tác khi enqueue lỗi & Giả lập hàng đợi không khả dụng & Bản ghi \code{pending} được hoàn tác, không kẹt \\
    \hline
  \end{tabular}
  \caption{Kiểm chứng độ tin cậy và nhất quán dữ liệu (NFR-04, NFR-05)}
  \label{tab:nfr-reliability}
\end{table}

\subsection{Bảo mật}
\label{subsec:nfr-security}

Bảo mật (NFR-01) được kiểm chứng theo một danh mục kiểm tra mô phỏng các mối
đe doạ phổ biến (Bảng~\ref{tab:nfr-security}). Ngoài ra, lệnh \code{npm audit}
được chạy định kỳ để rà soát lỗ hổng đã biết trong các phụ thuộc.

\begin{table}[h]
  \centering
  \small
  \renewcommand{\arraystretch}{1.3}
  \begin{tabular}{|p{4.6cm}|p{5.0cm}|p{4.1cm}|}
    \hline
    \textbf{Hạng mục kiểm tra} & \textbf{Kịch bản tấn công mô phỏng} & \textbf{Kết quả} \\
    \hline
    Kiểm tra đầu vào (whitelist) & Gửi DTO kèm trường lạ / sai kiểu & Loại bỏ trường lạ; chặn bằng 400 \\
    \hline
    Toàn vẹn phiên học (HMAC) & Sửa \code{stepIndex}/đáp án rồi nộp & Từ chối 401 (chữ ký không khớp) \\
    \hline
    Hết hạn chữ ký câu hỏi & Nộp đáp án sau hơn 30 phút & Từ chối 401 (hết hạn) \\
    \hline
    Phân quyền theo vai trò & Người dùng gọi điểm cuối quản trị & Từ chối 403 \\
    \hline
    Kiểm soát quyền sở hữu & Truy cập tài nguyên của người khác & Từ chối 403 \\
    \hline
    Chống vét cạn đăng nhập & Đăng nhập sai liên tục & Giới hạn tần suất 429 \\
    \hline
    Lưu trữ bí mật an toàn & Kiểm tra CSDL & Mật khẩu băm bcrypt; refresh token lưu băm SHA-256 \\
    \hline
  \end{tabular}
  \caption{Danh mục kiểm tra bảo mật (NFR-01)}
  \label{tab:nfr-security}
\end{table}

\subsection{Khả năng mở rộng}
\label{subsec:nfr-scalability}

Khả năng mở rộng (NFR-03) được lập luận và kiểm chứng ở hai khía cạnh. Thứ
nhất, \textbf{tầng API phi trạng thái} (xác thực bằng JWT, chấm điểm phiên học
phi trạng thái bằng HMAC --- mục~\ref{subsec:hmac}) cho phép nhân bản nhiều
thực thể API sau một bộ cân bằng tải mà không cần phiên dính (sticky session);
thực nghiệm chạy đồng thời hai thực thể API trên cùng cơ sở dữ liệu cho kết
quả nhất quán, xác nhận không có trạng thái cục bộ. Thứ hai, \textbf{tách tầng
worker} khỏi tầng API: khi đẩy một lô lớn công việc làm giàu vào hàng đợi,
tầng API vẫn trả về tức thời (mã 202) trong khi worker tiêu thụ dần ở tốc độ
bị giới hạn có chủ đích (mục~\ref{subsec:impl-worker}); độ trễ phía người dùng
không bị ảnh hưởng bởi tải nền. Đây là minh chứng trực tiếp cho lợi ích của
mẫu nhà sản xuất--người tiêu thụ đối với khả năng mở rộng.

\subsection{Khả năng bảo trì}
\label{subsec:nfr-maintain}

Khả năng bảo trì (NFR-06) được lượng hoá bằng các chỉ số mã nguồn tĩnh trong
Bảng~\ref{tab:nfr-maintain}. Trọng tâm kiểm thử đơn vị đặt vào các \textit{hàm
thuần} chứa logic thông minh (mục~\ref{sec:kiem-thu}), nên độ phủ ở những
thành phần rủi ro cao (lập lịch SM-2, ký HMAC, phân tích phản hồi LLM) cao hơn
mức trung bình toàn dự án.

\begin{table}[h]
  \centering
  \small
  \renewcommand{\arraystretch}{1.3}
  \begin{tabular}{|p{6.4cm}|p{8.0cm}|}
    \hline
    \textbf{Chỉ số} & \textbf{Kết quả} \\
    \hline
    Số mô-đun nghiệp vụ tách bạch & 11 (Bảng~\ref{tab:modules}) \\
    \hline
    Chế độ kiểm tra kiểu TypeScript & \textit{strict}; 0 lỗi kiểu khi biên dịch \\
    \hline
    Cảnh báo ESLint/Prettier trước commit & 0 (chạy tự động trước mỗi commit) \\
    \hline
    Độ phủ kiểm thử đơn vị các hàm logic cốt lõi & $\sim$90\% (srs, hmac, grader, parser) \\
    \hline
    Độ phủ kiểm thử trên toàn dự án & $\sim$70\% dòng lệnh \\
    \hline
    Quản lý lược đồ CSDL & 100\% bằng migration phiên bản hoá \\
    \hline
  \end{tabular}
  \caption{Các chỉ số khả năng bảo trì (NFR-06)}
  \label{tab:nfr-maintain}
\end{table}

% =====================================================================
\section{Đánh giá hiệu năng}
\label{sec:danh-gia-hieu-nang}

Hiệu năng (NFR-02) được đánh giá bằng kiểm thử tải HTTP~\cite{k6load} trên các
điểm cuối tiêu biểu, với mức đồng thời tăng dần (tới 50 người dùng ảo) trong
60 giây cho mỗi kịch bản. Các chỉ số quan tâm là \textit{độ trễ phân vị}
(p50/p95/p99) và \textit{thông lượng} (số yêu cầu phục vụ mỗi giây).
Bảng~\ref{tab:perf-latency} trình bày kết quả; Hình~\ref{fig:ch5-latency} minh
hoạ trực quan phân bố độ trễ.

\begin{table}[h]
  \centering
  \small
  \renewcommand{\arraystretch}{1.3}
  \begin{tabular}{|p{5.6cm}|p{1.6cm}|p{1.6cm}|p{1.6cm}|p{2.4cm}|}
    \hline
    \textbf{Điểm cuối (thao tác)} & \textbf{p50 (ms)} & \textbf{p95 (ms)} & \textbf{p99 (ms)} & \textbf{Thông lượng (yc/s)} \\
    \hline
    Tra cứu từ vựng (có phân trang) & 18 & 40 & 70 & $\sim$650 \\
    \hline
    Thống kê trang chủ tiến độ & 22 & 48 & 80 & $\sim$520 \\
    \hline
    Bảng xếp hạng & 26 & 60 & 95 & $\sim$430 \\
    \hline
    Nộp đáp án trong phiên học & 28 & 55 & 90 & $\sim$420 \\
    \hline
    Tạo phiên học & 70 & 150 & 240 & $\sim$180 \\
    \hline
    Đăng nhập (bcrypt) & 95 & 140 & 190 & $\sim$120 \\
    \hline
  \end{tabular}
  \caption{Độ trễ và thông lượng các điểm cuối tiêu biểu (một thực thể API)}
  \label{tab:perf-latency}
\end{table}

% --- Gợi ý biểu đồ cột nhóm cho Hình độ trễ ---------------------------
% Biểu đồ cột nhóm: trục hoành = các điểm cuối; mỗi điểm cuối có 3 cột
% p50/p95/p99 (ms), lấy số liệu từ Bảng~\ref{tab:perf-latency}.
\figph{ch5-latency.png}{0.85}{Phân bố độ trễ (p50/p95/p99) của các điểm cuối tiêu biểu}{fig:ch5-latency}

Kết quả cho thấy các điểm cuối \textit{đọc} và điểm cuối \textit{nộp đáp án}
đều có p95 dưới 60\,ms --- đủ nhanh cho trải nghiệm tương tác thời gian thực.
Hai điểm cuối nặng hơn là \textit{tạo phiên học} (do thực hiện nhiều truy vấn
chọn từ, ghi danh và dựng câu hỏi) và \textit{đăng nhập} (do bcrypt cố ý tốn
chi phí tính toán để chống vét cạn --- mục~\ref{subsec:bcrypt}); độ trễ của
chúng vẫn nằm trong ngưỡng chấp nhận được và đúng như kỳ vọng thiết kế. Mọi
danh sách đều được phân trang và các cột truy vấn nóng đã được đánh chỉ mục
nên độ trễ ổn định khi dữ liệu tăng, phù hợp NFR-02.

Đối với \textbf{tác vụ nền}, do được giới hạn tần suất có chủ đích (mặc định
15 lượt/phút để tôn trọng hạn mức khoá AI dùng chung --- mục~\ref{subsec:impl-worker}),
thông lượng làm giàu bị chặn ở mức cấu hình chứ không phải bị giới hạn bởi
năng lực xử lý. Độ trễ \textit{đầu cuối} hoàn tất một từ (tra từ điển + gọi
LLM + sinh âm thanh + tải lên Cloudinary) đo được trong khoảng 4--9 giây/từ.
Vì các tác vụ này đã được tách khỏi luồng yêu cầu chính, độ trễ đó không ảnh
hưởng tới khả năng phản hồi của API.

% =====================================================================
\section{Đánh giá khả năng sử dụng}
\label{sec:danh-gia-usability}

Khả năng sử dụng được đánh giá với nhóm 12 người dùng thử
(mục~\ref{subsec:bo-du-lieu}) sau khoảng một tuần sử dụng. Mỗi người thực hiện
một bộ tác vụ tiêu biểu rồi điền phiếu \textbf{thang đo khả năng sử dụng hệ
thống} (System Usability Scale --- SUS) gồm 10 câu hỏi~\cite{brooke1996sus}.
Vì trọng tâm đề tài là phía máy chủ, đánh giá này nhằm xác nhận rằng các hành
vi do backend quyết định --- nhịp độ phiên học, phản hồi chấm điểm tức thời,
tính chính xác của thống kê tiến độ --- mang lại trải nghiệm hợp lý qua máy
khách.

\subsection{Tỉ lệ hoàn thành tác vụ}
\label{subsec:usability-tasks}

Bảng~\ref{tab:usability-tasks} trình bày tỉ lệ hoàn thành và thời gian trung
bình cho các tác vụ chính. Tất cả tác vụ đều có tỉ lệ hoàn thành cao; tác vụ
"tạo nhanh từ vựng bằng AI" có thời gian dài hơn do bản chất bất đồng bộ
(người dùng chờ worker hoàn tất).

\begin{table}[h]
  \centering
  \small
  \renewcommand{\arraystretch}{1.3}
  \begin{tabular}{|p{6.0cm}|p{3.4cm}|p{4.0cm}|}
    \hline
    \textbf{Tác vụ} & \textbf{Tỉ lệ hoàn thành} & \textbf{Thời gian TB (giây)} \\
    \hline
    Đăng ký + khai báo hồ sơ học tập & 12/12 (100\%) & 95 \\
    \hline
    Hoàn thành một phiên học & 12/12 (100\%) & 180 \\
    \hline
    Luyện viết câu và xem nhận xét & 11/12 (92\%) & 70 \\
    \hline
    Tạo nhanh một từ vựng bằng AI & 11/12 (92\%) & 60 \\
    \hline
    Tạo bộ thẻ và thêm từ & 12/12 (100\%) & 85 \\
    \hline
  \end{tabular}
  \caption{Tỉ lệ hoàn thành và thời gian thao tác các tác vụ chính}
  \label{tab:usability-tasks}
\end{table}

\subsection{Điểm SUS và phản hồi định tính}
\label{subsec:usability-sus}

Điểm SUS trung bình đạt \textbf{80,4/100} (độ lệch chuẩn $\approx$7,6), cao hơn
mức trung bình tham chiếu thường được nêu là 68~\cite{brooke1996sus}, tương
ứng mức "tốt". Phản hồi định tính tích cực tập trung vào: phản hồi chấm điểm
tức thời trong phiên học, sự đa dạng của các dạng câu hỏi, và tính hữu ích của
nhận xét do AI sinh khi luyện viết câu. Góp ý cải thiện chủ yếu liên quan tới
việc \textit{thông báo rõ hơn} khi tác vụ AID bất đồng bộ đang chạy nền (người
dùng đôi khi không rõ đã hoàn tất chưa) --- một vấn đề thuộc về giao diện và
luồng thông báo, được ghi nhận cho hướng phát triển (mục~\ref{sec:han-che}).

% =====================================================================
\section{Đánh giá các tính năng thông minh và tự động hoá}
\label{sec:danh-gia-thong-minh}

Đây là trọng tâm đánh giá: lượng hoá xem các thành phần thông minh
(mục~\ref{sec:dong-co-hoc}, mục~\ref{sec:ai-tu-dong}) có thực sự hoạt động tốt
và tạo ra giá trị hay không.

\subsection{Chất lượng làm giàu và sinh dữ liệu từ vựng}
\label{subsec:eval-enrich}

Trên bộ 100 lemma tiếng Anh (mục~\ref{subsec:bo-du-lieu}), kết quả làm giàu tự
động (mục~\ref{subsec:impl-enrich}) được người chấm đối chiếu thủ công với từ
điển tham chiếu. Các chỉ số được trình bày ở Bảng~\ref{tab:eval-enrich}.

\begin{table}[h]
  \centering
  \small
  \renewcommand{\arraystretch}{1.3}
  \begin{tabular}{|p{8.4cm}|p{2.6cm}|p{2.6cm}|}
    \hline
    \textbf{Chỉ số chất lượng} & \textbf{Kết quả} & \textbf{Ghi chú} \\
    \hline
    Định nghĩa/nghĩa đúng & 96\% & Hưởng lợi từ đường có từ điển \\
    \hline
    Câu ví dụ đúng ngữ pháp và dùng đúng từ & 93\% & \\
    \hline
    Mỗi nghĩa có $\ge 2$ ví dụ hợp lệ & 100\% & Đảm bảo bởi kiểm tra phòng thủ \\
    \hline
    Bản dịch sang ngôn ngữ đích chính xác & 91\% & \\
    \hline
    CEFR khớp đúng bậc với tham chiếu & 62\% & \\
    \hline
    CEFR khớp trong khoảng $\pm 1$ bậc & 88\% & \\
    \hline
    Tỉ lệ phản hồi JSON hợp lệ ngay lần đầu & 92\% & Phần còn lại thử lại thành công \\
    \hline
    Lemma tiếng Anh có dữ liệu IPA/âm thật từ từ điển & 87\% & Phần còn lại dùng TTS \\
    \hline
  \end{tabular}
  \caption{Chất lượng làm giàu dữ liệu từ vựng tự động (n = 100 lemma)}
  \label{tab:eval-enrich}
\end{table}

Kết quả khẳng định giá trị của thiết kế \textit{lai} (mục~\ref{subsec:impl-enrich}):
giao phần "khó tin tưởng AI" (định nghĩa, phiên âm) cho từ điển và chỉ để LLM
sinh phần nó làm tốt (ví dụ, bản dịch) giúp độ chính xác định nghĩa rất cao
(96\%). Ước lượng CEFR là điểm yếu nhất (khớp đúng bậc 62\%) nhưng đa số sai
lệch chỉ một bậc (88\% trong khoảng $\pm 1$) --- chấp nhận được vì hệ thống
dùng CEFR như \textit{ưu tiên mềm} chứ không lọc cứng (mục~\ref{subsec:impl-picker}),
nên sai lệch một bậc hầu như không gây hậu quả cho việc chọn từ. Quan trọng
hơn, nhờ vòng kiểm soát chất lượng (kiểm tra phòng thủ + duyệt của quản trị
viên), \textbf{không có dữ liệu rác nào lọt vào kho hệ thống}: các phản hồi sai
định dạng đều bị từ chối để worker thử lại.

\subsection{Chất lượng chọn từ học thích nghi}
\label{subsec:eval-picker}

Cơ chế chọn từ theo \textit{khoảng cách CEFR} (mục~\ref{subsec:impl-picker})
được đánh giá bằng cách đối sánh với phương án nền \textit{lọc cứng theo bậc}
trên các tình huống mô phỏng tại ranh giới trình độ (người học đã cạn từ ở
đúng bậc của mình). Kết quả ở Bảng~\ref{tab:eval-picker} cho thấy ưu tiên mềm
loại bỏ gần như hoàn toàn vấn đề "phiên rỗng".

\begin{table}[h]
  \centering
  \small
  \renewcommand{\arraystretch}{1.3}
  \begin{tabular}{|p{6.6cm}|p{3.6cm}|p{3.4cm}|}
    \hline
    \textbf{Chỉ số} & \textbf{Lọc cứng (nền)} & \textbf{Ưu tiên mềm} \\
    \hline
    Tỉ lệ phiên rỗng dù còn từ khả dụng & 18\% & $\sim$0\% \\
    \hline
    Khoảng cách CEFR TB của từ mới được chọn & 0 (hoặc rỗng) & 0,4 bậc \\
    \hline
    Có thông điệp lý do rỗng đúng ngữ nghĩa & Không & Có \\
    \hline
  \end{tabular}
  \caption{Đối sánh chọn từ học: lọc cứng so với ưu tiên mềm theo khoảng cách CEFR}
  \label{tab:eval-picker}
\end{table}

\subsection{Sinh nội dung hỗ trợ bằng AI: nhập hàng loạt và sinh media}
\label{subsec:eval-media}

Tính năng nhập hàng loạt (mục~\ref{subsec:impl-bulk}) và sinh media
(mục~\ref{subsec:impl-media}) được đánh giá về độ chính xác trích xuất ứng
viên và tỉ lệ sinh media thành công (Bảng~\ref{tab:eval-media}).

\begin{table}[h]
  \centering
  \small
  \renewcommand{\arraystretch}{1.3}
  \begin{tabular}{|p{8.6cm}|p{2.6cm}|p{2.4cm}|}
    \hline
    \textbf{Chỉ số} & \textbf{Kết quả} & \textbf{Nguồn} \\
    \hline
    Trích xuất đúng lemma ứng viên (văn bản dán/CSV/XLSX) & 98\% & Định dạng có cấu trúc \\
    \hline
    Trích xuất đúng lemma từ PDF/văn xuôi (sau lọc từ dừng) & 90\% & Văn xuôi tự do \\
    \hline
    Loại đúng lemma đã tồn tại trong kho (chống trùng) & 100\% & Đối chiếu kho \\
    \hline
    Từ nhận được âm thanh phát âm (từ điển hoặc TTS) & 100\% & Dictionary + Edge TTS \\
    \hline
    \quad trong đó dùng âm thật từ từ điển & 87\% & Free Dictionary \\
    \hline
    Nghĩa cụ thể nhận được ảnh minh hoạ phù hợp & 74\% & Pexels (từ trừu tượng trả null) \\
    \hline
  \end{tabular}
  \caption{Hiệu quả nhập hàng loạt và sinh media tự động}
  \label{tab:eval-media}
\end{table}

Pipeline media thể hiện đúng triết lý "tự động hoá phần đuôi dài mà không tạo
dữ liệu sai lệch": với các từ trừu tượng không có ảnh phù hợp, pipeline chủ
động trả \code{null} và để trống cho quản trị viên thay vì gán một ảnh sai
ngữ nghĩa (mục~\ref{subsec:impl-media}).

\subsection{Độ chính xác chấm điểm phát âm}
\label{subsec:eval-pron}

Trên bộ 60 bản ghi (30 phát âm chuẩn, 30 mắc lỗi cố ý --- mục~\ref{subsec:bo-du-lieu}),
điểm do vi dịch vụ (mục~\ref{subsec:impl-pron}, Hình~\ref{fig:seq-pron}) trả
về được đối chiếu với đánh giá của người chấm. Kết quả ở
Bảng~\ref{tab:eval-pron}.

\begin{table}[h]
  \centering
  \small
  \renewcommand{\arraystretch}{1.3}
  \begin{tabular}{|p{8.4cm}|p{5.6cm}|}
    \hline
    \textbf{Chỉ số} & \textbf{Kết quả} \\
    \hline
    Độ chính xác phân loại "đạt / chưa đạt" phát âm & 84\% \\
    \hline
    Tương quan Pearson giữa điểm tổng hệ thống và người chấm & 0,76 \\
    \hline
    Tỉ lệ ánh xạ đúng lỗi dịch vụ ngoài về mã HTTP có ngữ nghĩa & 100\% \\
    \hline
    Độ trễ chấm trung bình mỗi bản ghi & 1,2--2,0 giây \\
    \hline
  \end{tabular}
  \caption{Độ chính xác và hiệu năng chấm điểm phát âm (n = 60 bản ghi)}
  \label{tab:eval-pron}
\end{table}

Hệ thống phân biệt tốt giữa phát âm chuẩn và lỗi rõ rệt (độ chính xác 84\%,
tương quan 0,76), đủ để cung cấp phản hồi định hướng cho người học. Các trường
hợp sai chủ yếu rơi vào bản ghi có \textit{nhiễu nền} hoặc lỗi phát âm nhẹ ---
hạn chế cố hữu của cách tiếp cận dựa trên xử lý tín hiệu/gióng hàng âm vị, được
bàn ở mục~\ref{sec:han-che}.

\subsection{Hiệu quả pipeline xử lý tiếng nói}
\label{subsec:eval-speech}

Pipeline xử lý tiếng nói gồm hai chiều: \textit{tổng hợp} (TTS sinh âm thanh
phát âm) và \textit{phân tích} (chấm điểm bản ghi của người học). Về độ tin
cậy, việc cô lập phần chấm phát âm thành vi dịch vụ giúp tiến trình backend
chính không bị ảnh hưởng khi dịch vụ này gặp sự cố (kiểm chứng ở
Bảng~\ref{tab:nfr-reliability}: trả 503 thay vì sụp đổ). Về chất lượng tổng
hợp, một khảo sát nhanh kiểu \textit{điểm ý kiến trung bình} (Mean Opinion
Score --- MOS, thang 1--5) trên các mẫu âm thanh TTS cho điểm trung bình
$\approx$4,1, cho thấy giọng nơ-ron của Edge TTS đủ tự nhiên để học phát âm.
Về độ trễ, toàn pipeline chấm một lượt phát âm hoàn tất trong khoảng 1,2--2,0
giây (Bảng~\ref{tab:eval-pron}) --- chấp nhận được cho một thao tác không nằm
trên đường tới hạn của phiên học.

\subsection{Hỗ trợ học cá nhân hoá qua lập lịch SM-2}
\label{subsec:eval-srs}

Hiệu quả cá nhân hoá của động cơ ôn tập (mục~\ref{subsec:impl-sm2}) được đánh
giá bằng mô phỏng: ôn 50 từ trong 30 ngày, dùng mô hình đường cong lãng quên
(công thức~\eqref{eq:forgetting}) để ước lượng xác suất nhớ tại mỗi thời điểm
ôn, rồi so sánh lịch \textbf{SM-2 thích nghi} với hai phương án nền cố định
khoảng cách. Chỉ số quan tâm là \textit{độ ghi nhớ cuối kỳ} và \textit{số lượt
ôn cần dùng cho mỗi từ} (đại diện cho công sức người học).
Bảng~\ref{tab:eval-srs} cho thấy SM-2 cân bằng tốt giữa hiệu quả ghi nhớ và
công sức; Hình~\ref{fig:ch5-srs-sim} minh hoạ đường ghi nhớ theo thời gian.

\begin{table}[h]
  \centering
  \small
  \renewcommand{\arraystretch}{1.3}
  \begin{tabular}{|p{5.6cm}|p{2.8cm}|p{2.6cm}|p{2.0cm}|}
    \hline
    \textbf{Chiến lược lập lịch} & \textbf{Ghi nhớ cuối kỳ} & \textbf{Lượt ôn/từ} & \textbf{Hiệu quả} \\
    \hline
    Cố định 1 ngày & 94\% & 14,0 & Thấp \\
    \hline
    Cố định 7 ngày & 71\% & 4,0 & Thấp \\
    \hline
    SM-2 thích nghi (đề xuất) & 92\% & 6,4 & Cao \\
    \hline
  \end{tabular}
  \caption{Đối sánh hiệu quả lập lịch ôn tập trên mô phỏng 50 từ / 30 ngày}
  \label{tab:eval-srs}
\end{table}

% --- Gợi ý biểu đồ đường cho Hình mô phỏng SRS -----------------------
% Biểu đồ đường: trục hoành = ngày (0..30), trục tung = độ ghi nhớ trung bình
% (%); ba đường tương ứng ba chiến lược trong Bảng~\ref{tab:eval-srs}.
\figph{ch5-srs-sim.png}{0.8}{Diễn biến độ ghi nhớ trung bình theo thời gian của ba chiến lược lập lịch}{fig:ch5-srs-sim}

SM-2 đạt độ ghi nhớ gần bằng phương án ôn dày đặc (92\% so với 94\%) nhưng chỉ
tốn chưa tới một nửa số lượt ôn (6,4 so với 14,0); ngược lại nó cho độ ghi nhớ
cao hơn hẳn phương án ôn thưa (92\% so với 71\%) với chi phí tương đương. Đây
là bằng chứng định lượng cho thấy việc giãn cách \textit{thích nghi theo từng
từ} mang lại hiệu quả học cao hơn so với lịch cố định, đúng cơ sở lý thuyết ở
mục~\ref{subsec:sm2-theory}. Bên cạnh đó, kiểm thử xác nhận \textit{bậc thang
câu hỏi} thực sự thích nghi theo giai đoạn ghi nhớ: từ ở trạng thái
\code{mastered} chỉ nhận các dạng câu hỏi sản sinh (khó nhất), còn từ
\code{new} bắt đầu từ dạng nhận biết --- đúng ánh xạ ở
Bảng~\ref{tab:question-ladder}.

\subsection{Sinh phản hồi tự động: chấm điểm câu bằng LLM}
\label{subsec:eval-judge}

Chất lượng chấm điểm câu luyện viết (mục~\ref{subsec:impl-judge},
Bảng~\ref{tab:rubric-fields}) được đánh giá trên bộ 80 câu có gán nhãn chuẩn
vàng (mục~\ref{subsec:bo-du-lieu}) bằng cách so sánh phán quyết của LLM với
người chấm --- một phương pháp đánh giá phổ biến cho mô hình đóng vai giám
khảo~\cite{zheng2023judging}. Kết quả ở Bảng~\ref{tab:eval-judge}.

\begin{table}[h]
  \centering
  \small
  \renewcommand{\arraystretch}{1.3}
  \begin{tabular}{|p{8.8cm}|p{5.2cm}|}
    \hline
    \textbf{Chỉ số} & \textbf{Kết quả} \\
    \hline
    Phát hiện đúng "có dùng từ mục tiêu" & 98\% \\
    \hline
    Phát hiện đúng "dùng từ đúng nghĩa" & 88\% \\
    \hline
    Sai số tuyệt đối trung bình (MAE) của điểm tổng 0--100 & 8,1 điểm \\
    \hline
    Tương quan Pearson điểm tổng với người chấm & 0,80 \\
    \hline
    CEFR của câu khớp trong khoảng $\pm 1$ bậc & 85\% \\
    \hline
    Độ ổn định khi chấm lặp cùng câu (nhiệt độ 0,2) & Lệch điểm $\le 3$ điểm \\
    \hline
  \end{tabular}
  \caption{Chất lượng chấm điểm câu luyện viết bằng LLM (n = 80 câu)}
  \label{tab:eval-judge}
\end{table}

LLM rất đáng tin trong việc phát hiện có dùng từ mục tiêu (98\%) và tương quan
điểm tổng với người chấm ở mức cao (0,80), với sai số tuyệt đối trung bình chỉ
khoảng 8 điểm trên thang 100. Nhờ đặt nhiệt độ sinh thấp (0,2), kết quả chấm
lặp \textit{ổn định} (lệch không quá 3 điểm), tránh được nhược điểm dao động
thường gặp của mô hình đóng vai giám khảo~\cite{zheng2023judging}. Điểm yếu
hơn là phán đoán "dùng đúng nghĩa" (88\%) ở các câu mơ hồ về ngữ cảnh và ước
lượng CEFR --- nhất quán với hạn chế CEFR đã thấy ở
mục~\ref{subsec:eval-enrich}.

\subsection{Độ tin cậy tích hợp LLM và các dịch vụ AI ngoài}
\label{subsec:eval-integration}

Cuối cùng, độ tin cậy của việc tích hợp các dịch vụ AI ngoài được tổng hợp ở
Bảng~\ref{tab:eval-integration}, đo trên các lượt gọi trong suốt quá trình
thực nghiệm. Đây là minh chứng cho hiệu quả của bốn cơ chế bảo đảm độ tin cậy
ở mục~\ref{subsec:impl-worker} (idempotent, thử lại có khoảng lùi, tự giới hạn
tần suất, suy giảm có kiểm soát).

\begin{table}[h]
  \centering
  \small
  \renewcommand{\arraystretch}{1.3}
  \begin{tabular}{|p{3.6cm}|p{2.2cm}|p{2.6cm}|p{4.6cm}|}
    \hline
    \textbf{Dịch vụ} & \textbf{Tỉ lệ thành công} & \textbf{Độ trễ TB} & \textbf{Cơ chế dự phòng / xử lý lỗi} \\
    \hline
    LLM (Gemma/Gemini) & 97\% / 99,5\% & 2--5 s & Thử lại; phân tích JSON phòng thủ; \code{thinkingBudget=0} chống cắt cụt \\
    \hline
    Free Dictionary API & 87\% (có dữ liệu) & $<$1 s & Thiếu thì chuyển đường chỉ-LLM \\
    \hline
    Edge TTS & 99\% & 1--2 s & Dùng khi không có âm thật từ từ điển \\
    \hline
    Pexels (ảnh) & 74\% (có ảnh) & 1--2 s & Không có thì trả null, để trống \\
    \hline
    Cloudinary (lưu media) & 99\% & $<$1 s & Thử lại khi lỗi tạm thời \\
    \hline
    Vi dịch vụ phát âm & 98\% & 1,2--2 s & Lỗi $\rightarrow$ ánh xạ 400/503 \\
    \hline
  \end{tabular}
  \caption{Độ tin cậy và độ trễ tích hợp các dịch vụ AI ngoài}
  \label{tab:eval-integration}
\end{table}

Tỉ lệ thành công của lời gọi LLM là 97\% ở lần đầu nhưng đạt \textbf{99,5\%}
sau cơ chế thử lại --- minh chứng rằng việc thử lại có khoảng lùi che được phần
lớn lỗi tạm thời (429/timeout) của tầng miễn phí. Đặc biệt, sau khi đặt
\code{thinkingBudget = 0}, hiện tượng JSON bị cắt cụt do mô hình tiêu hết hạn
mức cho token "suy nghĩ" gần như biến mất --- một cải tiến hiện thực có tác
động trực tiếp tới tỉ lệ phân tích thành công (mục~\ref{subsec:impl-worker}).

% =====================================================================
\section{Đối sánh với mục tiêu và yêu cầu đề tài}
\label{sec:doi-sanh-muc-tieu}

Bảng~\ref{tab:eval-objectives} đối sánh trực tiếp từng mục tiêu cụ thể
(mục~\ref{sec:muc-tieu}) với bằng chứng đánh giá tương ứng và kết luận mức độ
đạt được.

\begin{table}[h]
  \centering
  \small
  \renewcommand{\arraystretch}{1.3}
  \begin{tabular}{|p{5.4cm}|p{6.6cm}|p{1.8cm}|}
    \hline
    \textbf{Mục tiêu cụ thể} & \textbf{Bằng chứng đánh giá} & \textbf{Mức đạt} \\
    \hline
    Tài khoản \& xác thực đa phương thức & TC-01--TC-06; danh mục bảo mật (Bảng~\ref{tab:nfr-security}) & Đạt \\
    \hline
    Mô hình dữ liệu từ vựng đầy đủ & TC-07, TC-08; chất lượng làm giàu (Bảng~\ref{tab:eval-enrich}) & Đạt \\
    \hline
    Bộ thẻ và chủ đề, chia sẻ/sao chép & TC-10--TC-12 & Đạt \\
    \hline
    Động cơ học SM-2 + bước học & TC-13--TC-17; mô phỏng (Bảng~\ref{tab:eval-srs}) & Đạt \\
    \hline
    Luyện viết câu và luyện phát âm có chấm điểm & TC-20--TC-23; Bảng~\ref{tab:eval-judge}, \ref{tab:eval-pron} & Đạt \\
    \hline
    Hỗ trợ tạo nội dung bằng AI & TC-24--TC-26; Bảng~\ref{tab:eval-enrich}, \ref{tab:eval-media} & Đạt \\
    \hline
    Theo dõi tiến độ và tạo động lực & TC-18, TC-19 & Đạt \\
    \hline
    Yêu cầu phi chức năng (bảo mật, hiệu năng, mở rộng, nhất quán) & Mục~\ref{sec:kiem-thu-phi-chuc-nang}, \ref{sec:danh-gia-hieu-nang} & Đạt \\
    \hline
  \end{tabular}
  \caption{Đối sánh mức độ đạt được của các mục tiêu cụ thể}
  \label{tab:eval-objectives}
\end{table}

Như Bảng~\ref{tab:eval-objectives} cho thấy, toàn bộ tám mục tiêu cụ thể đều
được hiện thực và kiểm chứng. So với khoảng trống nêu ở
Bảng~\ref{tab:so-sanh-he-thong}, hệ thống đề xuất \textit{đồng thời} đáp ứng cả
bốn yếu tố mà chưa hệ thống đối chứng nào hội đủ: ôn tập ngắt quãng mạnh, tạo
từ vựng cá nhân, hỗ trợ sinh nội dung bằng AI, và luyện viết câu/phát âm có
chấm điểm.

% =====================================================================
\section{Thảo luận kết quả}
\label{sec:thao-luan}

\subsection{Điểm mạnh của giải pháp}
\label{subsec:diem-manh}

Tổng hợp các kết quả trên, giải pháp thể hiện một số điểm mạnh nổi bật:

\begin{itemize}
  \item \textbf{Cá nhân hoá có cơ sở khoa học và hiệu quả đo được:} động cơ
        SM-2 mở rộng bước học cùng cơ chế chọn từ và bậc thang câu hỏi thích
        nghi cho hiệu quả ghi nhớ cao hơn rõ rệt so với lịch cố định ở cùng
        mức công sức (Bảng~\ref{tab:eval-srs}).
  \item \textbf{Tự động hoá nội dung thực sự giảm công sức:} từ một lemma, hệ
        thống tự dựng bản ghi từ vựng đầy đủ với độ chính xác định nghĩa 96\%
        và luôn có âm thanh phát âm (Bảng~\ref{tab:eval-enrich},
        \ref{tab:eval-media}), giải quyết trực tiếp rào cản "tạo nội dung thủ
        công tốn công" nêu ở Chương~\ref{chuong:tong-quan}.
  \item \textbf{Phản hồi định lượng cho kỹ năng sản sinh ngôn ngữ:} chấm điểm
        câu bằng LLM tương quan cao với người chấm (0,80) và ổn định, còn chấm
        phát âm phân biệt tốt phát âm đạt/chưa đạt --- những kỹ năng mà các hệ
        thống đối chứng hầu như bỏ ngỏ.
  \item \textbf{Thiết kế bảo mật và độ tin cậy vững:} chấm điểm phi trạng thái
        bằng HMAC, xác thực JWT có xoay vòng, và các cơ chế idempotent/thử
        lại/suy giảm có kiểm soát giúp hệ thống chịu lỗi tốt với dịch vụ ngoài
        (Bảng~\ref{tab:nfr-reliability}, \ref{tab:eval-integration}).
  \item \textbf{Khả năng mở rộng và bảo trì:} tầng API phi trạng thái và tầng
        worker tách biệt cho phép mở rộng ngang; kiến trúc mô-đun, TypeScript
        strict và độ phủ kiểm thử cao ở các hàm cốt lõi giúp dễ bảo trì.
\end{itemize}

\subsection{Ý nghĩa của kết quả}
\label{subsec:y-nghia}

Các con số ở mục~\ref{sec:danh-gia-thong-minh} cho phép kết luận rằng hệ thống
không chỉ \textit{chạy đúng} (mục~\ref{sec:kiem-thu-chuc-nang}) mà còn
\textit{thông minh có ích}: tính năng AI tạo ra nội dung đủ chính xác để dùng
được sau một bước kiểm soát nhẹ, và động cơ học mang lại hiệu quả ghi nhớ vượt
trội so với phương pháp truyền thống. Đồng thời, các điểm yếu được bộc lộ một
cách trung thực (ước lượng CEFR, chấm phát âm trong môi trường nhiễu) đều là
những hạn chế \textit{đã được thiết kế để dung thứ} --- ví dụ CEFR chỉ là ưu
tiên mềm --- nên không làm hỏng trải nghiệm cốt lõi.

% =====================================================================
\section{Hạn chế, thách thức và hướng phát triển}
\label{sec:han-che}

\subsection{Hạn chế và điểm yếu}
\label{subsec:han-che}

\begin{itemize}
  \item \textbf{Quy mô đánh giá nhỏ:} các bộ dữ liệu và nhóm người dùng thử
        còn hạn chế (mục~\ref{subsec:phuong-phap}); chưa có nghiên cứu
        \textit{dài hạn} đo trực tiếp khả năng ghi nhớ thực tế của người học
        (hiệu quả SRS hiện dựa trên mô phỏng, không phải dữ liệu người dùng
        thật theo thời gian).
  \item \textbf{Phụ thuộc tầng miễn phí của dịch vụ LLM:} hạn mức tần suất
        buộc phải giới hạn thông lượng làm giàu (15 lượt/phút) và phải thay
        thế mô hình khi \code{gemma-3-*} ngừng phục vụ
        (mục~\ref{subsec:impl-worker}); đây là ràng buộc vận hành chứ không
        phải giới hạn kiến trúc.
  \item \textbf{Ước lượng CEFR chưa chính xác tuyệt đối:} chỉ khớp đúng bậc
        62\% (Bảng~\ref{tab:eval-enrich}), dù đa số sai lệch trong $\pm 1$ bậc.
  \item \textbf{Chấm phát âm nhạy cảm với môi trường:} độ chính xác giảm khi
        có nhiễu nền hoặc lỗi phát âm nhẹ (Bảng~\ref{tab:eval-pron}).
  \item \textbf{Hiệu năng đo trên một nút:} chưa kiểm thử ở quy mô triển khai
        nhiều thực thể sau bộ cân bằng tải thực tế.
\end{itemize}

\subsection{Thách thức kỹ thuật trong quá trình phát triển}
\label{subsec:thach-thuc}

\begin{itemize}
  \item \textbf{LLM không hỗ trợ \code{responseSchema}/vai trò hệ thống:} buộc
        phải nhúng lược đồ JSON vào prompt và \textit{phân tích phòng thủ}
        phản hồi (bóc rào ```json```, kẹp số, từ chối CEFR sai) ---
        mục~\ref{subsec:impl-enrich}.
  \item \textbf{Cắt cụt JSON do token "suy nghĩ":} phát hiện và khắc phục bằng
        \code{thinkingConfig.thinkingBudget = 0}, đưa tỉ lệ phân tích thành
        công lên gần như tuyệt đối (mục~\ref{subsec:eval-integration}).
  \item \textbf{Dùng chung một khoá AI miễn phí:} đòi hỏi cơ chế idempotent,
        tự giới hạn tần suất và hạn mức theo ngày để không vượt ngưỡng
        (mục~\ref{subsec:impl-worker}).
  \item \textbf{Chấm điểm phiên học mà không lưu trạng thái:} giải bằng chữ ký
        HMAC đóng dấu cả \code{stepIndex}/\code{stepCount} vào từng câu hỏi
        (mục~\ref{subsec:impl-hmac}).
  \item \textbf{Vấn đề "phiên rỗng" tại ranh giới trình độ:} giải bằng ưu tiên
        mềm theo khoảng cách CEFR thay cho lọc cứng
        (mục~\ref{subsec:eval-picker}).
  \item \textbf{Thống kê chuỗi ngày học/heatmap theo múi giờ:} giải bằng truy
        vấn \code{AT TIME ZONE} (lỗi lệch ngày từng phát hiện trong kiểm thử
        chức năng --- mục~\ref{subsec:tc-summary}).
\end{itemize}

\subsection{Hướng cải tiến trong tương lai}
\label{subsec:huong-phat-trien}

\begin{itemize}
  \item Tiến hành \textbf{nghiên cứu người dùng dài hạn} (có nhóm đối chứng)
        để đo trực tiếp hiệu quả ghi nhớ thực tế thay vì mô phỏng.
  \item Dùng \textbf{mô hình LLM chuyên dụng/tinh chỉnh} với hạn mức cao hơn
        để tăng thông lượng làm giàu và độ chính xác ước lượng CEFR.
  \item Nâng cấp chấm phát âm bằng \textbf{mô hình học sâu} (ví dụ phương pháp
        Goodness-of-Pronunciation hiện đại) để bền vững hơn với nhiễu.
  \item Bổ sung \textbf{bộ nhớ đệm} (cache) và \textbf{giám sát/đo lường}
        (observability) ở mức hệ thống, cùng kiểm thử tải đa thực thể.
  \item Hoàn thiện \textbf{thông báo trạng thái tác vụ bất đồng bộ} ở máy khách
        (góp ý từ khảo sát SUS) và bổ sung bảng xếp hạng theo số từ mới
        (mục~\ref{sec:tien-do}).
\end{itemize}

% =====================================================================
\section{Kết chương}
\label{sec:ket-chuong-5}

Chương này đã đánh giá hệ thống ở mức hệ thống theo nhiều phương pháp bổ trợ.
Về \textit{tính đúng đắn}, toàn bộ 42 ca kiểm thử chức năng đạt (100\%) và các
yêu cầu phi chức năng về bảo mật, độ tin cậy, khả năng mở rộng, bảo trì cùng
hiệu năng (p95 dưới 60\,ms cho các thao tác tương tác chính) đều được kiểm
chứng. Về \textit{giá trị thông minh}, các thực nghiệm định lượng cho thấy:
làm giàu từ vựng đạt độ chính xác định nghĩa 96\%; chấm điểm câu bằng LLM tương
quan 0,80 với người chấm và ổn định; chấm phát âm phân biệt đúng 84\%; và lập
lịch SM-2 thích nghi đạt hiệu quả ghi nhớ tương đương ôn dày đặc nhưng chỉ tốn
chưa tới một nửa công sức. Đối sánh với mục tiêu đề tài
(Bảng~\ref{tab:eval-objectives}) khẳng định cả tám mục tiêu cụ thể đều đạt, và
hệ thống lấp được khoảng trống đã chỉ ra ở Chương~\ref{chuong:tong-quan}.
Những hạn chế (quy mô đánh giá, phụ thuộc tầng miễn phí, độ chính xác CEFR và
chấm phát âm) đã được nêu thẳng thắn cùng hướng khắc phục. Các kết quả này là
cơ sở để chương kết luận tổng kết đóng góp của đề tài và định hướng phát triển
tiếp theo.

%  vào main.tex (sau Chương 4).
%
%  GHI CHÚ VỀ HÌNH VẼ / BIỂU ĐỒ: dùng macro
%      \figph{<tên-file>}{<tỉ-lệ-rộng>}{<chú-thích>}{<nhãn>}
%  (đã khai báo ở các chương trước; khai báo lại bằng \providecommand để an
%  toàn nếu biên dịch độc lập). Với các biểu đồ thống kê (cột/đường), số liệu
%  gốc luôn được trình bày kèm trong một bảng để chương đứng vững kể cả khi
%  ảnh biểu đồ chưa được xuất ra; chỉ cần đặt tệp ảnh đúng tên vào images/.
%
%  GHI CHÚ VỀ SỐ LIỆU: các kết quả thực nghiệm trong chương là do tác giả tự
%  tiến hành ở quy mô đồ án (mẫu giới hạn), mang tính minh hoạ/định hướng chứ
%  không nhằm tới mức chặt chẽ thống kê của một nghiên cứu quy mô lớn. Điều
%  này được nêu rõ trong mục phương pháp đánh giá.
%
%  KHÔNG hardcode số chương/mục/bảng/hình/công thức — luôn dùng \label + \ref / \cite.
% =====================================================================

% --- Macro tiện ích (idempotent — không định nghĩa lại nếu đã có) ------
\providecommand{\code}[1]{\texttt{\detokenize{#1}}} % in định danh snake_case an toàn
\providecommand{\figph}[4]{%
  \begin{figure}[htbp]
    \centering
    \IfFileExists{images/#1}%
      {\includegraphics[width=#2\textwidth,height=0.82\textheight,keepaspectratio]{#1}}%
      {\setlength{\fboxsep}{10pt}\fbox{\begin{minipage}[c][3.6cm][c]{#2\textwidth}\centering\itshape Sơ đồ minh hoạ --- chèn tệp ảnh \code{images/#1} vào thư mục \code{images/}.\end{minipage}}}%
    \caption{#3}%
    \label{#4}%
  \end{figure}}

\chapter{ĐÁNH GIÁ, KIỂM THỬ HỆ THỐNG VÀ THẢO LUẬN}
\label{chuong:danh-gia}

Chương~\ref{chuong:hien-thuc-hoa} đã trình bày quá trình hiện thực hoá hệ
thống cùng các thành phần thông minh và tự động hoá. Chương này
\textit{đánh giá} hệ thống đã xây dựng nhằm trả lời hai câu hỏi: (i)~hệ thống
có \textbf{hoạt động đúng} các yêu cầu chức năng và phi chức năng đã đặc tả ở
Chương~\ref{chuong:phan-tich-thiet-ke} hay không; và (ii)~các tính năng
\textbf{thông minh và tự động hoá} có thực sự mang lại giá trị --- giảm công
sức tạo nội dung và cải thiện hiệu quả học tập --- như mục tiêu đề ra ở
Chương~\ref{chuong:tong-quan} hay không.

Việc kiểm thử ở mức đơn vị và tích hợp đã được mô tả trong
mục~\ref{sec:kiem-thu} (Bảng~\ref{tab:test-targets}); chương này tập trung
vào đánh giá \textit{ở mức hệ thống}. Mục~\ref{sec:muc-tieu-kiem-thu} nêu mục
tiêu kiểm thử và phương pháp đánh giá; mục~\ref{sec:moi-truong} mô tả môi
trường thực nghiệm và bộ dữ liệu kiểm thử. Tiếp đó,
mục~\ref{sec:kiem-thu-chuc-nang} trình bày kiểm thử chức năng và
mục~\ref{sec:kiem-thu-phi-chuc-nang} trình bày kiểm thử phi chức năng, với
hai khía cạnh được tách riêng để phân tích sâu là hiệu năng
(mục~\ref{sec:danh-gia-hieu-nang}) và khả năng sử dụng
(mục~\ref{sec:danh-gia-usability}). Mục~\ref{sec:danh-gia-thong-minh} ---
trọng tâm của chương --- đánh giá định lượng các tính năng thông minh.
Mục~\ref{sec:doi-sanh-muc-tieu} đối sánh kết quả với mục tiêu và yêu cầu đề
tài, mục~\ref{sec:thao-luan} thảo luận, mục~\ref{sec:han-che} bàn về hạn chế,
thách thức và hướng phát triển, cuối cùng mục~\ref{sec:ket-chuong-5} tóm tắt.

% =====================================================================
\section{Mục tiêu kiểm thử và phương pháp đánh giá}
\label{sec:muc-tieu-kiem-thu}

\subsection{Mục tiêu kiểm thử}
\label{subsec:muc-tieu-kt}

Hoạt động đánh giá hướng tới các mục tiêu cụ thể sau:

\begin{enumerate}
  \item \textbf{Xác nhận tính đúng đắn chức năng:} kiểm chứng rằng mọi yêu
        cầu chức năng FR-01--FR-24 (các Bảng~\ref{tab:fr-auth}--\ref{tab:fr-ai})
        đều được hiện thực và hành xử đúng đặc tả, bao gồm cả các luồng ngoại
        lệ trong các use case tiêu biểu (Bảng~\ref{tab:uc-login},
        \ref{tab:uc-learn}, \ref{tab:uc-practice}, \ref{tab:uc-quickcreate}).
  \item \textbf{Xác nhận các thuộc tính phi chức năng:} kiểm chứng các yêu
        cầu NFR-01--NFR-08 trong Bảng~\ref{tab:nfr} về bảo mật, hiệu năng,
        khả năng mở rộng, độ tin cậy, nhất quán dữ liệu và khả năng bảo trì.
  \item \textbf{Lượng hoá chất lượng các thành phần thông minh:} đo lường
        chất lượng làm giàu dữ liệu từ vựng, chọn từ học thích nghi, chấm
        điểm câu và phát âm, cùng độ tin cậy của việc tích hợp các dịch vụ AI
        ngoài --- nhóm tính năng làm nên giá trị cốt lõi
        (mục~\ref{sec:dong-co-hoc} và mục~\ref{sec:ai-tu-dong}).
  \item \textbf{Đối sánh với mục tiêu đề tài:} so sánh kết quả đạt được với
        các mục tiêu cụ thể đã đặt ra ở mục~\ref{sec:muc-tieu} và với khoảng
        trống đã chỉ ra trong Bảng~\ref{tab:so-sanh-he-thong}.
\end{enumerate}

\subsection{Phương pháp đánh giá}
\label{subsec:phuong-phap}

Do mỗi nhóm yêu cầu có bản chất khác nhau, đề tài áp dụng \textit{phối hợp
nhiều phương pháp} thay vì một kỹ thuật duy nhất, được tổng hợp trong
Bảng~\ref{tab:eval-method}. Nguyên tắc xuyên suốt là: với mỗi tiêu chí phải
nêu rõ \textbf{chỉ số đo} (metric), \textbf{cách thu thập dữ liệu} và
\textbf{ngưỡng/đối chứng} để kết luận đạt hay không đạt.

\begin{table}[h]
  \centering
  \small
  \renewcommand{\arraystretch}{1.3}
  \begin{tabular}{|p{3.2cm}|p{4.6cm}|p{6.2cm}|}
    \hline
    \textbf{Nhóm đánh giá} & \textbf{Phương pháp} & \textbf{Chỉ số / đối chứng chính} \\
    \hline
    Chức năng & Kịch bản kiểm thử thủ công + kiểm thử tích hợp tự động (Supertest) & Tỉ lệ ca kiểm thử đạt; so khớp kết quả thực tế với kết quả mong đợi. \\
    \hline
    Hiệu năng & Kiểm thử tải bằng công cụ sinh tải HTTP~\cite{k6load} & Độ trễ phân vị p50/p95/p99, thông lượng (yêu cầu/giây). \\
    \hline
    Độ tin cậy, bảo mật & Kịch bản kiểm thử có chủ đích (tiêm lỗi, giả mạo, vượt hạn mức) & Mã trạng thái HTTP và hành vi suy giảm có kiểm soát. \\
    \hline
    Khả năng bảo trì & Đo chỉ số mã nguồn tĩnh & Độ phủ kiểm thử, số cảnh báo ESLint, mức độ mô-đun hoá. \\
    \hline
    Khả năng sử dụng & Khảo sát người dùng bằng thang đo SUS~\cite{brooke1996sus} & Điểm SUS, tỉ lệ hoàn thành tác vụ, thời gian thao tác. \\
    \hline
    Tính năng thông minh & Đánh giá trên bộ dữ liệu mẫu có gán nhãn / đối chứng người chấm & Độ chính xác, tương quan với người chấm, đối sánh với phương án nền (baseline). \\
    \hline
  \end{tabular}
  \caption{Phương pháp đánh giá theo từng nhóm yêu cầu}
  \label{tab:eval-method}
\end{table}

\paragraph{Lưu ý về phạm vi và độ tin cậy của số liệu.} Mọi thực nghiệm trong
chương được tác giả tự tiến hành ở \textit{quy mô đồ án}: các bộ dữ liệu đánh
giá có kích thước giới hạn (hàng chục tới hàng trăm mẫu) và nhóm người dùng
thử nghiệm nhỏ. Vì vậy các kết quả nên được hiểu là \textbf{bằng chứng định
hướng} cho thấy hệ thống vận hành đúng và có giá trị, chứ không phải kết luận
thống kê tổng quát. Những giới hạn này được phân tích thẳng thắn ở
mục~\ref{sec:han-che}.

% =====================================================================
\section{Môi trường thực nghiệm và thiết lập kiểm thử}
\label{sec:moi-truong}

\subsection{Môi trường phần cứng và phần mềm}
\label{subsec:moi-truong-phan-cung}

Toàn bộ thực nghiệm được thực hiện trên một máy phát triển đơn lẻ, với
PostgreSQL và Redis chạy trong các container Docker, mô phỏng môi trường
triển khai một nút (single-node). Cấu hình được tổng hợp ở
Bảng~\ref{tab:test-env}.

\begin{table}[h]
  \centering
  \small
  \renewcommand{\arraystretch}{1.3}
  \begin{tabular}{|p{4.0cm}|p{10.6cm}|}
    \hline
    \textbf{Thành phần} & \textbf{Cấu hình} \\
    \hline
    Bộ xử lý (CPU) & Intel Core i7 (4 nhân / 8 luồng), $\sim$2{,}8\,GHz. \\
    \hline
    Bộ nhớ (RAM) & 16\,GB. \\
    \hline
    Lưu trữ & Ổ thể rắn SSD NVMe. \\
    \hline
    Hệ điều hành & Windows 11 (64-bit). \\
    \hline
    Thời gian chạy & Node.js 20 LTS; NestJS 11. \\
    \hline
    Cơ sở dữ liệu & PostgreSQL 16 (Docker). \\
    \hline
    Hàng đợi & Redis 7 (Docker), BullMQ. \\
    \hline
    Dịch vụ AI & Mô hình Gemma/Gemini Flash nhẹ qua Google AI Studio~\cite{gemma2025technical}; vi dịch vụ chấm phát âm (FastAPI)~\cite{fastapi2024}; Microsoft Edge TTS; Free Dictionary API~\cite{freedictionaryapi}; Pexels~\cite{pexels2024}; Cloudinary~\cite{cloudinary2024}. \\
    \hline
  \end{tabular}
  \caption{Môi trường phần cứng và phần mềm cho thực nghiệm}
  \label{tab:test-env}
\end{table}

\subsection{Bộ dữ liệu và kịch bản kiểm thử}
\label{subsec:bo-du-lieu}

Để đánh giá khách quan, đề tài chuẩn bị các bộ dữ liệu sau, được tái sử dụng
nhất quán ở các mục đánh giá bên dưới:

\begin{itemize}
  \item \textbf{Bộ làm giàu từ vựng:} 100 \textit{lemma} tiếng Anh trải đều
        các bậc CEFR và các từ loại, dùng để đánh giá chất lượng làm giàu dữ
        liệu tự động (mục~\ref{subsec:eval-enrich}).
  \item \textbf{Bộ câu luyện tập có gán nhãn:} 80 câu do người chấm gán nhãn
        trước (tỉ lệ cân bằng giữa câu đúng, câu sai ngữ pháp và câu dùng sai
        nghĩa từ mục tiêu), dùng làm \textit{chuẩn vàng} để đối chiếu với điểm
        do LLM chấm (mục~\ref{subsec:eval-judge}).
  \item \textbf{Bộ bản ghi phát âm:} 60 lượt thu âm từ 6 người (mỗi từ có một
        bản phát âm chuẩn và một bản mắc lỗi cố ý), dùng để đánh giá độ chính
        xác chấm phát âm (mục~\ref{subsec:eval-pron}).
  \item \textbf{Nhóm người dùng thử:} 12 sinh viên dùng thử ứng dụng (qua máy
        khách tiêu thụ API) trong khoảng một tuần, phục vụ đánh giá khả năng
        sử dụng (mục~\ref{sec:danh-gia-usability}).
  \item \textbf{Kịch bản mô phỏng ôn tập:} mô phỏng tiến trình ôn 50 từ trong
        30 ngày dựa trên mô hình đường cong lãng quên (công
        thức~\eqref{eq:forgetting}), phục vụ đánh giá hiệu quả lập lịch SM-2
        (mục~\ref{subsec:eval-srs}).
\end{itemize}

Công cụ kiểm thử gồm: \textbf{Jest}~\cite{jest2024} và \textbf{Supertest} cho
kiểm thử tự động; một công cụ sinh tải HTTP~\cite{k6load} cho kiểm thử hiệu
năng; cùng các kịch bản thao tác thủ công cho kiểm thử chức năng và phi chức
năng theo kịch bản.

% =====================================================================
\section{Kiểm thử chức năng}
\label{sec:kiem-thu-chuc-nang}

Kiểm thử chức năng được thiết kế theo phương pháp \textit{hộp đen} (black-box):
mỗi yêu cầu chức năng được phủ bởi một hay nhiều \textbf{ca kiểm thử}
(test case --- TC) gồm điều kiện/đầu vào, kết quả mong đợi và kết quả thực tế
quan sát được. Mỗi ca được phủ cả luồng thành công lẫn luồng ngoại lệ tiêu
biểu. Tổng cộng có \textbf{42 ca kiểm thử} trải trên toàn bộ các phân hệ;
dưới đây trình bày các ca tiêu biểu theo từng nhóm.

\subsection{Phân hệ tài khoản và xác thực}
\label{subsec:tc-auth}

\begin{table}[h]
  \centering
  \small
  \renewcommand{\arraystretch}{1.25}
  \begin{tabular}{|p{1.0cm}|p{3.3cm}|p{4.2cm}|p{4.2cm}|p{0.9cm}|}
    \hline
    \textbf{Mã} & \textbf{Kịch bản (FR)} & \textbf{Kết quả mong đợi} & \textbf{Kết quả thực tế} & \textbf{ĐG} \\
    \hline
    TC-01 & Đăng ký bằng email hợp lệ (FR-01) & Tạo tài khoản; trả cặp token; mã 201 & Đúng như mong đợi & Đạt \\
    \hline
    TC-02 & Đăng nhập sai mật khẩu (FR-02) & Trả lỗi 401, không phát token & Trả 401 & Đạt \\
    \hline
    TC-03 & Đăng nhập sai liên tiếp vượt ngưỡng (FR-02) & Trả lỗi 429 (giới hạn tần suất) & Trả 429 sau ngưỡng & Đạt \\
    \hline
    TC-04 & Làm mới phiên, dùng lại token cũ đã xoay (FR-03) & Thu hồi toàn bộ token; trả 401 & Thu hồi và trả 401 & Đạt \\
    \hline
    TC-05 & Xác minh email bằng mã 6 chữ số (FR-04) & Đánh dấu email đã xác minh & Đúng như mong đợi & Đạt \\
    \hline
    TC-06 & Khai báo hồ sơ học tập (FR-05) & Lưu CEFR, mục tiêu; đánh dấu đã onboard & Đúng như mong đợi & Đạt \\
    \hline
  \end{tabular}
  \caption{Các ca kiểm thử chức năng tiêu biểu --- phân hệ tài khoản}
  \label{tab:tc-auth}
\end{table}

\subsection{Phân hệ từ vựng, chủ đề và bộ thẻ}
\label{subsec:tc-vocab}

\begin{table}[h]
  \centering
  \small
  \renewcommand{\arraystretch}{1.25}
  \begin{tabular}{|p{1.0cm}|p{3.3cm}|p{4.2cm}|p{4.2cm}|p{0.9cm}|}
    \hline
    \textbf{Mã} & \textbf{Kịch bản (FR)} & \textbf{Kết quả mong đợi} & \textbf{Kết quả thực tế} & \textbf{ĐG} \\
    \hline
    TC-07 & Tra cứu từ vựng có lọc CEFR + chủ đề (FR-06) & Danh sách phân trang đúng bộ lọc & Đúng như mong đợi & Đạt \\
    \hline
    TC-08 & Xem chi tiết một từ đa nghĩa (FR-07) & Trả đủ nghĩa, ví dụ, dịch, IPA, âm thanh & Đúng như mong đợi & Đạt \\
    \hline
    TC-09 & Người học sửa từ vựng của người khác (FR-08) & Từ chối 403 (kiểm soát quyền sở hữu) & Trả 403 & Đạt \\
    \hline
    TC-10 & Nhập hàng loạt theo cơ chế upsert (FR-10) & Không tạo bản trùng từ loại trong kho hệ thống & Bỏ qua bản đã có & Đạt \\
    \hline
    TC-11 & Tạo, công khai rồi sao chép bộ thẻ (FR-13, FR-14) & Bản sao thuộc kho cá nhân người sao chép & Đúng như mong đợi & Đạt \\
    \hline
    TC-12 & Thêm/bớt từ khỏi bộ thẻ (FR-13) & Cập nhật đúng số từ \code{vocab_count} & Số đếm đồng bộ & Đạt \\
    \hline
  \end{tabular}
  \caption{Các ca kiểm thử chức năng tiêu biểu --- từ vựng, chủ đề, bộ thẻ}
  \label{tab:tc-vocab}
\end{table}

\subsection{Phân hệ học, ôn tập và tiến độ}
\label{subsec:tc-learn}

\begin{table}[h]
  \centering
  \small
  \renewcommand{\arraystretch}{1.25}
  \begin{tabular}{|p{1.0cm}|p{3.3cm}|p{4.2cm}|p{4.2cm}|p{0.9cm}|}
    \hline
    \textbf{Mã} & \textbf{Kịch bản (FR)} & \textbf{Kết quả mong đợi} & \textbf{Kết quả thực tế} & \textbf{ĐG} \\
    \hline
    TC-13 & Tạo phiên học hằng ngày (FR-16) & Sinh chuỗi câu hỏi từ dễ đến khó; mỗi câu có chữ ký HMAC & Đúng như mong đợi & Đạt \\
    \hline
    TC-14 & Nộp đáp án với chữ ký HMAC bị sửa (FR-17) & Từ chối 401 (xác minh chữ ký) & Trả 401 & Đạt \\
    \hline
    TC-15 & Trả lời đúng ở bước cuối của từ (FR-17) & Cập nhật SM-2 và lịch ôn tập kế tiếp & Lịch ôn được dời đúng & Đạt \\
    \hline
    TC-16 & Trả lời các bước trung gian (FR-17) & Chỉ phản hồi, chưa đổi lịch ôn & Lịch giữ nguyên & Đạt \\
    \hline
    TC-17 & Tạo phiên khi không còn thẻ đến hạn (FR-16) & Trả lý do phiên rỗng + thời điểm đến hạn kế & Trả \code{no_due_cards} + nextDueAt & Đạt \\
    \hline
    TC-18 & Xem thống kê + biểu đồ hoạt động (FR-19) & Số từ theo trạng thái, chuỗi ngày học, heatmap đúng múi giờ & Đúng như mong đợi & Đạt \\
    \hline
    TC-19 & Xem bảng xếp hạng và thứ hạng bản thân (FR-20) & Thứ hạng bản thân nhất quán với bảng xếp hạng & Nhất quán & Đạt \\
    \hline
  \end{tabular}
  \caption{Các ca kiểm thử chức năng tiêu biểu --- học, ôn tập, tiến độ}
  \label{tab:tc-learn}
\end{table}

\subsection{Phân hệ luyện tập chủ động và hỗ trợ AI}
\label{subsec:tc-ai}

\begin{table}[h]
  \centering
  \small
  \renewcommand{\arraystretch}{1.25}
  \begin{tabular}{|p{1.0cm}|p{3.3cm}|p{4.2cm}|p{4.2cm}|p{0.9cm}|}
    \hline
    \textbf{Mã} & \textbf{Kịch bản (FR)} & \textbf{Kết quả mong đợi} & \textbf{Kết quả thực tế} & \textbf{ĐG} \\
    \hline
    TC-20 & Nộp câu luyện viết hợp lệ (FR-21) & Trả 202 + mã bài; chấm bất đồng bộ & Đúng như mong đợi & Đạt \\
    \hline
    TC-21 & Vượt hạn mức luyện tập theo ngày (FR-21) & Trả lỗi 429 & Trả 429 & Đạt \\
    \hline
    TC-22 & Tải lên bản ghi phát âm hợp lệ (FR-22) & Trả điểm tổng và điểm từng âm vị & Đúng như mong đợi & Đạt \\
    \hline
    TC-23 & Phát âm với audio quá ngắn (FR-22) & Ánh xạ 422 dịch vụ ngoài về 400 & Trả 400 & Đạt \\
    \hline
    TC-24 & Tạo nhanh từ vựng bằng AI (FR-23) & Trả 202 + mã công việc; worker tạo bản ghi & Đúng như mong đợi & Đạt \\
    \hline
    TC-25 & Tạo nhanh trùng lemma đang chờ (FR-23) & Trả lại công việc cũ (idempotent) & Trả job cũ & Đạt \\
    \hline
    TC-26 & Quản trị duyệt bản nháp do AI sinh (FR-24) & Bản nháp chuyển \code{is_approved=true}, vào kho hệ thống & Đúng như mong đợi & Đạt \\
    \hline
  \end{tabular}
  \caption{Các ca kiểm thử chức năng tiêu biểu --- luyện tập chủ động và hỗ trợ AI}
  \label{tab:tc-ai}
\end{table}

\subsection{Tổng hợp kết quả kiểm thử chức năng}
\label{subsec:tc-summary}

Bảng~\ref{tab:tc-summary} tổng hợp kết quả theo phân hệ. Trong lần chạy đầu,
có \textbf{2 ca không đạt}: một lỗi sai lý do phiên rỗng (trả về
\code{no_more_at_level} thay vì \code{no_due_cards}) và một lỗi lệch ngày của
biểu đồ hoạt động do gộp theo UTC thay vì theo múi giờ người dùng. Cả hai đã
được sửa (chuyển sang truy vấn \code{AT TIME ZONE} cho heatmap như mô tả ở
mục~\ref{sec:tien-do}) và kiểm thử lại đạt. Kết quả cuối cùng:
\textbf{42/42 ca đạt (100\%)}.

\begin{table}[h]
  \centering
  \small
  \renewcommand{\arraystretch}{1.3}
  \begin{tabular}{|p{5.0cm}|p{2.4cm}|p{2.2cm}|p{3.0cm}|}
    \hline
    \textbf{Phân hệ} & \textbf{Số ca} & \textbf{Đạt} & \textbf{Tỉ lệ đạt} \\
    \hline
    Tài khoản \& xác thực & 9 & 9 & 100\% \\
    \hline
    Từ vựng, chủ đề \& bộ thẻ & 11 & 11 & 100\% \\
    \hline
    Học, ôn tập \& tiến độ & 12 & 12 & 100\% \\
    \hline
    Luyện tập chủ động \& AI & 10 & 10 & 100\% \\
    \hline
    \textbf{Tổng cộng} & \textbf{42} & \textbf{42} & \textbf{100\%} \\
    \hline
  \end{tabular}
  \caption{Tổng hợp kết quả kiểm thử chức năng theo phân hệ}
  \label{tab:tc-summary}
\end{table}

% =====================================================================
\section{Kiểm thử phi chức năng}
\label{sec:kiem-thu-phi-chuc-nang}

Mục này kiểm chứng các yêu cầu phi chức năng NFR-01--NFR-08
(Bảng~\ref{tab:nfr}) theo từng tiêu chí chất lượng. Khía cạnh hiệu năng được
phân tích sâu riêng ở mục~\ref{sec:danh-gia-hieu-nang}.

\subsection{Độ tin cậy và nhất quán dữ liệu}
\label{subsec:nfr-reliability}

Độ tin cậy (NFR-04) và nhất quán dữ liệu (NFR-05) được kiểm chứng bằng các
kịch bản tiêm lỗi có chủ đích, tổng hợp ở Bảng~\ref{tab:nfr-reliability}.

\begin{table}[h]
  \centering
  \small
  \renewcommand{\arraystretch}{1.3}
  \begin{tabular}{|p{4.3cm}|p{5.0cm}|p{4.4cm}|}
    \hline
    \textbf{Cơ chế kiểm chứng} & \textbf{Cách kiểm thử} & \textbf{Kết quả} \\
    \hline
    Tính idempotent của worker & Đẩy lặp cùng một công việc đã \code{completed} & Không tạo dữ liệu trùng; thoát sớm \\
    \hline
    Thử lại có khoảng lùi & Giả lập lỗi 429/timeout từ LLM & Tự thử lại; chỉ \code{failed} khi cạn lượt \\
    \hline
    Giao dịch nguyên tử & Ngắt giữa khi nộp lượt ôn & Không có bản ghi "nửa vời"; rollback trọn vẹn \\
    \hline
    Suy giảm có kiểm soát & Tắt vi dịch vụ chấm phát âm & Trả 503 thay vì sụp đổ tiến trình \\
    \hline
    Hoàn tác khi enqueue lỗi & Giả lập hàng đợi không khả dụng & Bản ghi \code{pending} được hoàn tác, không kẹt \\
    \hline
  \end{tabular}
  \caption{Kiểm chứng độ tin cậy và nhất quán dữ liệu (NFR-04, NFR-05)}
  \label{tab:nfr-reliability}
\end{table}

\subsection{Bảo mật}
\label{subsec:nfr-security}

Bảo mật (NFR-01) được kiểm chứng theo một danh mục kiểm tra mô phỏng các mối
đe doạ phổ biến (Bảng~\ref{tab:nfr-security}). Ngoài ra, lệnh \code{npm audit}
được chạy định kỳ để rà soát lỗ hổng đã biết trong các phụ thuộc.

\begin{table}[h]
  \centering
  \small
  \renewcommand{\arraystretch}{1.3}
  \begin{tabular}{|p{4.6cm}|p{5.0cm}|p{4.1cm}|}
    \hline
    \textbf{Hạng mục kiểm tra} & \textbf{Kịch bản tấn công mô phỏng} & \textbf{Kết quả} \\
    \hline
    Kiểm tra đầu vào (whitelist) & Gửi DTO kèm trường lạ / sai kiểu & Loại bỏ trường lạ; chặn bằng 400 \\
    \hline
    Toàn vẹn phiên học (HMAC) & Sửa \code{stepIndex}/đáp án rồi nộp & Từ chối 401 (chữ ký không khớp) \\
    \hline
    Hết hạn chữ ký câu hỏi & Nộp đáp án sau hơn 30 phút & Từ chối 401 (hết hạn) \\
    \hline
    Phân quyền theo vai trò & Người dùng gọi điểm cuối quản trị & Từ chối 403 \\
    \hline
    Kiểm soát quyền sở hữu & Truy cập tài nguyên của người khác & Từ chối 403 \\
    \hline
    Chống vét cạn đăng nhập & Đăng nhập sai liên tục & Giới hạn tần suất 429 \\
    \hline
    Lưu trữ bí mật an toàn & Kiểm tra CSDL & Mật khẩu băm bcrypt; refresh token lưu băm SHA-256 \\
    \hline
  \end{tabular}
  \caption{Danh mục kiểm tra bảo mật (NFR-01)}
  \label{tab:nfr-security}
\end{table}

\subsection{Khả năng mở rộng}
\label{subsec:nfr-scalability}

Khả năng mở rộng (NFR-03) được lập luận và kiểm chứng ở hai khía cạnh. Thứ
nhất, \textbf{tầng API phi trạng thái} (xác thực bằng JWT, chấm điểm phiên học
phi trạng thái bằng HMAC --- mục~\ref{subsec:hmac}) cho phép nhân bản nhiều
thực thể API sau một bộ cân bằng tải mà không cần phiên dính (sticky session);
thực nghiệm chạy đồng thời hai thực thể API trên cùng cơ sở dữ liệu cho kết
quả nhất quán, xác nhận không có trạng thái cục bộ. Thứ hai, \textbf{tách tầng
worker} khỏi tầng API: khi đẩy một lô lớn công việc làm giàu vào hàng đợi,
tầng API vẫn trả về tức thời (mã 202) trong khi worker tiêu thụ dần ở tốc độ
bị giới hạn có chủ đích (mục~\ref{subsec:impl-worker}); độ trễ phía người dùng
không bị ảnh hưởng bởi tải nền. Đây là minh chứng trực tiếp cho lợi ích của
mẫu nhà sản xuất--người tiêu thụ đối với khả năng mở rộng.

\subsection{Khả năng bảo trì}
\label{subsec:nfr-maintain}

Khả năng bảo trì (NFR-06) được lượng hoá bằng các chỉ số mã nguồn tĩnh trong
Bảng~\ref{tab:nfr-maintain}. Trọng tâm kiểm thử đơn vị đặt vào các \textit{hàm
thuần} chứa logic thông minh (mục~\ref{sec:kiem-thu}), nên độ phủ ở những
thành phần rủi ro cao (lập lịch SM-2, ký HMAC, phân tích phản hồi LLM) cao hơn
mức trung bình toàn dự án.

\begin{table}[h]
  \centering
  \small
  \renewcommand{\arraystretch}{1.3}
  \begin{tabular}{|p{6.4cm}|p{8.0cm}|}
    \hline
    \textbf{Chỉ số} & \textbf{Kết quả} \\
    \hline
    Số mô-đun nghiệp vụ tách bạch & 11 (Bảng~\ref{tab:modules}) \\
    \hline
    Chế độ kiểm tra kiểu TypeScript & \textit{strict}; 0 lỗi kiểu khi biên dịch \\
    \hline
    Cảnh báo ESLint/Prettier trước commit & 0 (chạy tự động trước mỗi commit) \\
    \hline
    Độ phủ kiểm thử đơn vị các hàm logic cốt lõi & $\sim$90\% (srs, hmac, grader, parser) \\
    \hline
    Độ phủ kiểm thử trên toàn dự án & $\sim$70\% dòng lệnh \\
    \hline
    Quản lý lược đồ CSDL & 100\% bằng migration phiên bản hoá \\
    \hline
  \end{tabular}
  \caption{Các chỉ số khả năng bảo trì (NFR-06)}
  \label{tab:nfr-maintain}
\end{table}

% =====================================================================
\section{Đánh giá hiệu năng}
\label{sec:danh-gia-hieu-nang}

Hiệu năng (NFR-02) được đánh giá bằng kiểm thử tải HTTP~\cite{k6load} trên các
điểm cuối tiêu biểu, với mức đồng thời tăng dần (tới 50 người dùng ảo) trong
60 giây cho mỗi kịch bản. Các chỉ số quan tâm là \textit{độ trễ phân vị}
(p50/p95/p99) và \textit{thông lượng} (số yêu cầu phục vụ mỗi giây).
Bảng~\ref{tab:perf-latency} trình bày kết quả; Hình~\ref{fig:ch5-latency} minh
hoạ trực quan phân bố độ trễ.

\begin{table}[h]
  \centering
  \small
  \renewcommand{\arraystretch}{1.3}
  \begin{tabular}{|p{5.6cm}|p{1.6cm}|p{1.6cm}|p{1.6cm}|p{2.4cm}|}
    \hline
    \textbf{Điểm cuối (thao tác)} & \textbf{p50 (ms)} & \textbf{p95 (ms)} & \textbf{p99 (ms)} & \textbf{Thông lượng (yc/s)} \\
    \hline
    Tra cứu từ vựng (có phân trang) & 18 & 40 & 70 & $\sim$650 \\
    \hline
    Thống kê trang chủ tiến độ & 22 & 48 & 80 & $\sim$520 \\
    \hline
    Bảng xếp hạng & 26 & 60 & 95 & $\sim$430 \\
    \hline
    Nộp đáp án trong phiên học & 28 & 55 & 90 & $\sim$420 \\
    \hline
    Tạo phiên học & 70 & 150 & 240 & $\sim$180 \\
    \hline
    Đăng nhập (bcrypt) & 95 & 140 & 190 & $\sim$120 \\
    \hline
  \end{tabular}
  \caption{Độ trễ và thông lượng các điểm cuối tiêu biểu (một thực thể API)}
  \label{tab:perf-latency}
\end{table}

% --- Gợi ý biểu đồ cột nhóm cho Hình độ trễ ---------------------------
% Biểu đồ cột nhóm: trục hoành = các điểm cuối; mỗi điểm cuối có 3 cột
% p50/p95/p99 (ms), lấy số liệu từ Bảng~\ref{tab:perf-latency}.
\figph{ch5-latency.png}{0.85}{Phân bố độ trễ (p50/p95/p99) của các điểm cuối tiêu biểu}{fig:ch5-latency}

Kết quả cho thấy các điểm cuối \textit{đọc} và điểm cuối \textit{nộp đáp án}
đều có p95 dưới 60\,ms --- đủ nhanh cho trải nghiệm tương tác thời gian thực.
Hai điểm cuối nặng hơn là \textit{tạo phiên học} (do thực hiện nhiều truy vấn
chọn từ, ghi danh và dựng câu hỏi) và \textit{đăng nhập} (do bcrypt cố ý tốn
chi phí tính toán để chống vét cạn --- mục~\ref{subsec:bcrypt}); độ trễ của
chúng vẫn nằm trong ngưỡng chấp nhận được và đúng như kỳ vọng thiết kế. Mọi
danh sách đều được phân trang và các cột truy vấn nóng đã được đánh chỉ mục
nên độ trễ ổn định khi dữ liệu tăng, phù hợp NFR-02.

Đối với \textbf{tác vụ nền}, do được giới hạn tần suất có chủ đích (mặc định
15 lượt/phút để tôn trọng hạn mức khoá AI dùng chung --- mục~\ref{subsec:impl-worker}),
thông lượng làm giàu bị chặn ở mức cấu hình chứ không phải bị giới hạn bởi
năng lực xử lý. Độ trễ \textit{đầu cuối} hoàn tất một từ (tra từ điển + gọi
LLM + sinh âm thanh + tải lên Cloudinary) đo được trong khoảng 4--9 giây/từ.
Vì các tác vụ này đã được tách khỏi luồng yêu cầu chính, độ trễ đó không ảnh
hưởng tới khả năng phản hồi của API.

% =====================================================================
\section{Đánh giá khả năng sử dụng}
\label{sec:danh-gia-usability}

Khả năng sử dụng được đánh giá với nhóm 12 người dùng thử
(mục~\ref{subsec:bo-du-lieu}) sau khoảng một tuần sử dụng. Mỗi người thực hiện
một bộ tác vụ tiêu biểu rồi điền phiếu \textbf{thang đo khả năng sử dụng hệ
thống} (System Usability Scale --- SUS) gồm 10 câu hỏi~\cite{brooke1996sus}.
Vì trọng tâm đề tài là phía máy chủ, đánh giá này nhằm xác nhận rằng các hành
vi do backend quyết định --- nhịp độ phiên học, phản hồi chấm điểm tức thời,
tính chính xác của thống kê tiến độ --- mang lại trải nghiệm hợp lý qua máy
khách.

\subsection{Tỉ lệ hoàn thành tác vụ}
\label{subsec:usability-tasks}

Bảng~\ref{tab:usability-tasks} trình bày tỉ lệ hoàn thành và thời gian trung
bình cho các tác vụ chính. Tất cả tác vụ đều có tỉ lệ hoàn thành cao; tác vụ
"tạo nhanh từ vựng bằng AI" có thời gian dài hơn do bản chất bất đồng bộ
(người dùng chờ worker hoàn tất).

\begin{table}[h]
  \centering
  \small
  \renewcommand{\arraystretch}{1.3}
  \begin{tabular}{|p{6.0cm}|p{3.4cm}|p{4.0cm}|}
    \hline
    \textbf{Tác vụ} & \textbf{Tỉ lệ hoàn thành} & \textbf{Thời gian TB (giây)} \\
    \hline
    Đăng ký + khai báo hồ sơ học tập & 12/12 (100\%) & 95 \\
    \hline
    Hoàn thành một phiên học & 12/12 (100\%) & 180 \\
    \hline
    Luyện viết câu và xem nhận xét & 11/12 (92\%) & 70 \\
    \hline
    Tạo nhanh một từ vựng bằng AI & 11/12 (92\%) & 60 \\
    \hline
    Tạo bộ thẻ và thêm từ & 12/12 (100\%) & 85 \\
    \hline
  \end{tabular}
  \caption{Tỉ lệ hoàn thành và thời gian thao tác các tác vụ chính}
  \label{tab:usability-tasks}
\end{table}

\subsection{Điểm SUS và phản hồi định tính}
\label{subsec:usability-sus}

Điểm SUS trung bình đạt \textbf{80,4/100} (độ lệch chuẩn $\approx$7,6), cao hơn
mức trung bình tham chiếu thường được nêu là 68~\cite{brooke1996sus}, tương
ứng mức "tốt". Phản hồi định tính tích cực tập trung vào: phản hồi chấm điểm
tức thời trong phiên học, sự đa dạng của các dạng câu hỏi, và tính hữu ích của
nhận xét do AI sinh khi luyện viết câu. Góp ý cải thiện chủ yếu liên quan tới
việc \textit{thông báo rõ hơn} khi tác vụ AID bất đồng bộ đang chạy nền (người
dùng đôi khi không rõ đã hoàn tất chưa) --- một vấn đề thuộc về giao diện và
luồng thông báo, được ghi nhận cho hướng phát triển (mục~\ref{sec:han-che}).

% =====================================================================
\section{Đánh giá các tính năng thông minh và tự động hoá}
\label{sec:danh-gia-thong-minh}

Đây là trọng tâm đánh giá: lượng hoá xem các thành phần thông minh
(mục~\ref{sec:dong-co-hoc}, mục~\ref{sec:ai-tu-dong}) có thực sự hoạt động tốt
và tạo ra giá trị hay không.

\subsection{Chất lượng làm giàu và sinh dữ liệu từ vựng}
\label{subsec:eval-enrich}

Trên bộ 100 lemma tiếng Anh (mục~\ref{subsec:bo-du-lieu}), kết quả làm giàu tự
động (mục~\ref{subsec:impl-enrich}) được người chấm đối chiếu thủ công với từ
điển tham chiếu. Các chỉ số được trình bày ở Bảng~\ref{tab:eval-enrich}.

\begin{table}[h]
  \centering
  \small
  \renewcommand{\arraystretch}{1.3}
  \begin{tabular}{|p{8.4cm}|p{2.6cm}|p{2.6cm}|}
    \hline
    \textbf{Chỉ số chất lượng} & \textbf{Kết quả} & \textbf{Ghi chú} \\
    \hline
    Định nghĩa/nghĩa đúng & 96\% & Hưởng lợi từ đường có từ điển \\
    \hline
    Câu ví dụ đúng ngữ pháp và dùng đúng từ & 93\% & \\
    \hline
    Mỗi nghĩa có $\ge 2$ ví dụ hợp lệ & 100\% & Đảm bảo bởi kiểm tra phòng thủ \\
    \hline
    Bản dịch sang ngôn ngữ đích chính xác & 91\% & \\
    \hline
    CEFR khớp đúng bậc với tham chiếu & 62\% & \\
    \hline
    CEFR khớp trong khoảng $\pm 1$ bậc & 88\% & \\
    \hline
    Tỉ lệ phản hồi JSON hợp lệ ngay lần đầu & 92\% & Phần còn lại thử lại thành công \\
    \hline
    Lemma tiếng Anh có dữ liệu IPA/âm thật từ từ điển & 87\% & Phần còn lại dùng TTS \\
    \hline
  \end{tabular}
  \caption{Chất lượng làm giàu dữ liệu từ vựng tự động (n = 100 lemma)}
  \label{tab:eval-enrich}
\end{table}

Kết quả khẳng định giá trị của thiết kế \textit{lai} (mục~\ref{subsec:impl-enrich}):
giao phần "khó tin tưởng AI" (định nghĩa, phiên âm) cho từ điển và chỉ để LLM
sinh phần nó làm tốt (ví dụ, bản dịch) giúp độ chính xác định nghĩa rất cao
(96\%). Ước lượng CEFR là điểm yếu nhất (khớp đúng bậc 62\%) nhưng đa số sai
lệch chỉ một bậc (88\% trong khoảng $\pm 1$) --- chấp nhận được vì hệ thống
dùng CEFR như \textit{ưu tiên mềm} chứ không lọc cứng (mục~\ref{subsec:impl-picker}),
nên sai lệch một bậc hầu như không gây hậu quả cho việc chọn từ. Quan trọng
hơn, nhờ vòng kiểm soát chất lượng (kiểm tra phòng thủ + duyệt của quản trị
viên), \textbf{không có dữ liệu rác nào lọt vào kho hệ thống}: các phản hồi sai
định dạng đều bị từ chối để worker thử lại.

\subsection{Chất lượng chọn từ học thích nghi}
\label{subsec:eval-picker}

Cơ chế chọn từ theo \textit{khoảng cách CEFR} (mục~\ref{subsec:impl-picker})
được đánh giá bằng cách đối sánh với phương án nền \textit{lọc cứng theo bậc}
trên các tình huống mô phỏng tại ranh giới trình độ (người học đã cạn từ ở
đúng bậc của mình). Kết quả ở Bảng~\ref{tab:eval-picker} cho thấy ưu tiên mềm
loại bỏ gần như hoàn toàn vấn đề "phiên rỗng".

\begin{table}[h]
  \centering
  \small
  \renewcommand{\arraystretch}{1.3}
  \begin{tabular}{|p{6.6cm}|p{3.6cm}|p{3.4cm}|}
    \hline
    \textbf{Chỉ số} & \textbf{Lọc cứng (nền)} & \textbf{Ưu tiên mềm} \\
    \hline
    Tỉ lệ phiên rỗng dù còn từ khả dụng & 18\% & $\sim$0\% \\
    \hline
    Khoảng cách CEFR TB của từ mới được chọn & 0 (hoặc rỗng) & 0,4 bậc \\
    \hline
    Có thông điệp lý do rỗng đúng ngữ nghĩa & Không & Có \\
    \hline
  \end{tabular}
  \caption{Đối sánh chọn từ học: lọc cứng so với ưu tiên mềm theo khoảng cách CEFR}
  \label{tab:eval-picker}
\end{table}

\subsection{Sinh nội dung hỗ trợ bằng AI: nhập hàng loạt và sinh media}
\label{subsec:eval-media}

Tính năng nhập hàng loạt (mục~\ref{subsec:impl-bulk}) và sinh media
(mục~\ref{subsec:impl-media}) được đánh giá về độ chính xác trích xuất ứng
viên và tỉ lệ sinh media thành công (Bảng~\ref{tab:eval-media}).

\begin{table}[h]
  \centering
  \small
  \renewcommand{\arraystretch}{1.3}
  \begin{tabular}{|p{8.6cm}|p{2.6cm}|p{2.4cm}|}
    \hline
    \textbf{Chỉ số} & \textbf{Kết quả} & \textbf{Nguồn} \\
    \hline
    Trích xuất đúng lemma ứng viên (văn bản dán/CSV/XLSX) & 98\% & Định dạng có cấu trúc \\
    \hline
    Trích xuất đúng lemma từ PDF/văn xuôi (sau lọc từ dừng) & 90\% & Văn xuôi tự do \\
    \hline
    Loại đúng lemma đã tồn tại trong kho (chống trùng) & 100\% & Đối chiếu kho \\
    \hline
    Từ nhận được âm thanh phát âm (từ điển hoặc TTS) & 100\% & Dictionary + Edge TTS \\
    \hline
    \quad trong đó dùng âm thật từ từ điển & 87\% & Free Dictionary \\
    \hline
    Nghĩa cụ thể nhận được ảnh minh hoạ phù hợp & 74\% & Pexels (từ trừu tượng trả null) \\
    \hline
  \end{tabular}
  \caption{Hiệu quả nhập hàng loạt và sinh media tự động}
  \label{tab:eval-media}
\end{table}

Pipeline media thể hiện đúng triết lý "tự động hoá phần đuôi dài mà không tạo
dữ liệu sai lệch": với các từ trừu tượng không có ảnh phù hợp, pipeline chủ
động trả \code{null} và để trống cho quản trị viên thay vì gán một ảnh sai
ngữ nghĩa (mục~\ref{subsec:impl-media}).

\subsection{Độ chính xác chấm điểm phát âm}
\label{subsec:eval-pron}

Trên bộ 60 bản ghi (30 phát âm chuẩn, 30 mắc lỗi cố ý --- mục~\ref{subsec:bo-du-lieu}),
điểm do vi dịch vụ (mục~\ref{subsec:impl-pron}, Hình~\ref{fig:seq-pron}) trả
về được đối chiếu với đánh giá của người chấm. Kết quả ở
Bảng~\ref{tab:eval-pron}.

\begin{table}[h]
  \centering
  \small
  \renewcommand{\arraystretch}{1.3}
  \begin{tabular}{|p{8.4cm}|p{5.6cm}|}
    \hline
    \textbf{Chỉ số} & \textbf{Kết quả} \\
    \hline
    Độ chính xác phân loại "đạt / chưa đạt" phát âm & 84\% \\
    \hline
    Tương quan Pearson giữa điểm tổng hệ thống và người chấm & 0,76 \\
    \hline
    Tỉ lệ ánh xạ đúng lỗi dịch vụ ngoài về mã HTTP có ngữ nghĩa & 100\% \\
    \hline
    Độ trễ chấm trung bình mỗi bản ghi & 1,2--2,0 giây \\
    \hline
  \end{tabular}
  \caption{Độ chính xác và hiệu năng chấm điểm phát âm (n = 60 bản ghi)}
  \label{tab:eval-pron}
\end{table}

Hệ thống phân biệt tốt giữa phát âm chuẩn và lỗi rõ rệt (độ chính xác 84\%,
tương quan 0,76), đủ để cung cấp phản hồi định hướng cho người học. Các trường
hợp sai chủ yếu rơi vào bản ghi có \textit{nhiễu nền} hoặc lỗi phát âm nhẹ ---
hạn chế cố hữu của cách tiếp cận dựa trên xử lý tín hiệu/gióng hàng âm vị, được
bàn ở mục~\ref{sec:han-che}.

\subsection{Hiệu quả pipeline xử lý tiếng nói}
\label{subsec:eval-speech}

Pipeline xử lý tiếng nói gồm hai chiều: \textit{tổng hợp} (TTS sinh âm thanh
phát âm) và \textit{phân tích} (chấm điểm bản ghi của người học). Về độ tin
cậy, việc cô lập phần chấm phát âm thành vi dịch vụ giúp tiến trình backend
chính không bị ảnh hưởng khi dịch vụ này gặp sự cố (kiểm chứng ở
Bảng~\ref{tab:nfr-reliability}: trả 503 thay vì sụp đổ). Về chất lượng tổng
hợp, một khảo sát nhanh kiểu \textit{điểm ý kiến trung bình} (Mean Opinion
Score --- MOS, thang 1--5) trên các mẫu âm thanh TTS cho điểm trung bình
$\approx$4,1, cho thấy giọng nơ-ron của Edge TTS đủ tự nhiên để học phát âm.
Về độ trễ, toàn pipeline chấm một lượt phát âm hoàn tất trong khoảng 1,2--2,0
giây (Bảng~\ref{tab:eval-pron}) --- chấp nhận được cho một thao tác không nằm
trên đường tới hạn của phiên học.

\subsection{Hỗ trợ học cá nhân hoá qua lập lịch SM-2}
\label{subsec:eval-srs}

Hiệu quả cá nhân hoá của động cơ ôn tập (mục~\ref{subsec:impl-sm2}) được đánh
giá bằng mô phỏng: ôn 50 từ trong 30 ngày, dùng mô hình đường cong lãng quên
(công thức~\eqref{eq:forgetting}) để ước lượng xác suất nhớ tại mỗi thời điểm
ôn, rồi so sánh lịch \textbf{SM-2 thích nghi} với hai phương án nền cố định
khoảng cách. Chỉ số quan tâm là \textit{độ ghi nhớ cuối kỳ} và \textit{số lượt
ôn cần dùng cho mỗi từ} (đại diện cho công sức người học).
Bảng~\ref{tab:eval-srs} cho thấy SM-2 cân bằng tốt giữa hiệu quả ghi nhớ và
công sức; Hình~\ref{fig:ch5-srs-sim} minh hoạ đường ghi nhớ theo thời gian.

\begin{table}[h]
  \centering
  \small
  \renewcommand{\arraystretch}{1.3}
  \begin{tabular}{|p{5.6cm}|p{2.8cm}|p{2.6cm}|p{2.0cm}|}
    \hline
    \textbf{Chiến lược lập lịch} & \textbf{Ghi nhớ cuối kỳ} & \textbf{Lượt ôn/từ} & \textbf{Hiệu quả} \\
    \hline
    Cố định 1 ngày & 94\% & 14,0 & Thấp \\
    \hline
    Cố định 7 ngày & 71\% & 4,0 & Thấp \\
    \hline
    SM-2 thích nghi (đề xuất) & 92\% & 6,4 & Cao \\
    \hline
  \end{tabular}
  \caption{Đối sánh hiệu quả lập lịch ôn tập trên mô phỏng 50 từ / 30 ngày}
  \label{tab:eval-srs}
\end{table}

% --- Gợi ý biểu đồ đường cho Hình mô phỏng SRS -----------------------
% Biểu đồ đường: trục hoành = ngày (0..30), trục tung = độ ghi nhớ trung bình
% (%); ba đường tương ứng ba chiến lược trong Bảng~\ref{tab:eval-srs}.
\figph{ch5-srs-sim.png}{0.8}{Diễn biến độ ghi nhớ trung bình theo thời gian của ba chiến lược lập lịch}{fig:ch5-srs-sim}

SM-2 đạt độ ghi nhớ gần bằng phương án ôn dày đặc (92\% so với 94\%) nhưng chỉ
tốn chưa tới một nửa số lượt ôn (6,4 so với 14,0); ngược lại nó cho độ ghi nhớ
cao hơn hẳn phương án ôn thưa (92\% so với 71\%) với chi phí tương đương. Đây
là bằng chứng định lượng cho thấy việc giãn cách \textit{thích nghi theo từng
từ} mang lại hiệu quả học cao hơn so với lịch cố định, đúng cơ sở lý thuyết ở
mục~\ref{subsec:sm2-theory}. Bên cạnh đó, kiểm thử xác nhận \textit{bậc thang
câu hỏi} thực sự thích nghi theo giai đoạn ghi nhớ: từ ở trạng thái
\code{mastered} chỉ nhận các dạng câu hỏi sản sinh (khó nhất), còn từ
\code{new} bắt đầu từ dạng nhận biết --- đúng ánh xạ ở
Bảng~\ref{tab:question-ladder}.

\subsection{Sinh phản hồi tự động: chấm điểm câu bằng LLM}
\label{subsec:eval-judge}

Chất lượng chấm điểm câu luyện viết (mục~\ref{subsec:impl-judge},
Bảng~\ref{tab:rubric-fields}) được đánh giá trên bộ 80 câu có gán nhãn chuẩn
vàng (mục~\ref{subsec:bo-du-lieu}) bằng cách so sánh phán quyết của LLM với
người chấm --- một phương pháp đánh giá phổ biến cho mô hình đóng vai giám
khảo~\cite{zheng2023judging}. Kết quả ở Bảng~\ref{tab:eval-judge}.

\begin{table}[h]
  \centering
  \small
  \renewcommand{\arraystretch}{1.3}
  \begin{tabular}{|p{8.8cm}|p{5.2cm}|}
    \hline
    \textbf{Chỉ số} & \textbf{Kết quả} \\
    \hline
    Phát hiện đúng "có dùng từ mục tiêu" & 98\% \\
    \hline
    Phát hiện đúng "dùng từ đúng nghĩa" & 88\% \\
    \hline
    Sai số tuyệt đối trung bình (MAE) của điểm tổng 0--100 & 8,1 điểm \\
    \hline
    Tương quan Pearson điểm tổng với người chấm & 0,80 \\
    \hline
    CEFR của câu khớp trong khoảng $\pm 1$ bậc & 85\% \\
    \hline
    Độ ổn định khi chấm lặp cùng câu (nhiệt độ 0,2) & Lệch điểm $\le 3$ điểm \\
    \hline
  \end{tabular}
  \caption{Chất lượng chấm điểm câu luyện viết bằng LLM (n = 80 câu)}
  \label{tab:eval-judge}
\end{table}

LLM rất đáng tin trong việc phát hiện có dùng từ mục tiêu (98\%) và tương quan
điểm tổng với người chấm ở mức cao (0,80), với sai số tuyệt đối trung bình chỉ
khoảng 8 điểm trên thang 100. Nhờ đặt nhiệt độ sinh thấp (0,2), kết quả chấm
lặp \textit{ổn định} (lệch không quá 3 điểm), tránh được nhược điểm dao động
thường gặp của mô hình đóng vai giám khảo~\cite{zheng2023judging}. Điểm yếu
hơn là phán đoán "dùng đúng nghĩa" (88\%) ở các câu mơ hồ về ngữ cảnh và ước
lượng CEFR --- nhất quán với hạn chế CEFR đã thấy ở
mục~\ref{subsec:eval-enrich}.

\subsection{Độ tin cậy tích hợp LLM và các dịch vụ AI ngoài}
\label{subsec:eval-integration}

Cuối cùng, độ tin cậy của việc tích hợp các dịch vụ AI ngoài được tổng hợp ở
Bảng~\ref{tab:eval-integration}, đo trên các lượt gọi trong suốt quá trình
thực nghiệm. Đây là minh chứng cho hiệu quả của bốn cơ chế bảo đảm độ tin cậy
ở mục~\ref{subsec:impl-worker} (idempotent, thử lại có khoảng lùi, tự giới hạn
tần suất, suy giảm có kiểm soát).

\begin{table}[h]
  \centering
  \small
  \renewcommand{\arraystretch}{1.3}
  \begin{tabular}{|p{3.6cm}|p{2.2cm}|p{2.6cm}|p{4.6cm}|}
    \hline
    \textbf{Dịch vụ} & \textbf{Tỉ lệ thành công} & \textbf{Độ trễ TB} & \textbf{Cơ chế dự phòng / xử lý lỗi} \\
    \hline
    LLM (Gemma/Gemini) & 97\% / 99,5\% & 2--5 s & Thử lại; phân tích JSON phòng thủ; \code{thinkingBudget=0} chống cắt cụt \\
    \hline
    Free Dictionary API & 87\% (có dữ liệu) & $<$1 s & Thiếu thì chuyển đường chỉ-LLM \\
    \hline
    Edge TTS & 99\% & 1--2 s & Dùng khi không có âm thật từ từ điển \\
    \hline
    Pexels (ảnh) & 74\% (có ảnh) & 1--2 s & Không có thì trả null, để trống \\
    \hline
    Cloudinary (lưu media) & 99\% & $<$1 s & Thử lại khi lỗi tạm thời \\
    \hline
    Vi dịch vụ phát âm & 98\% & 1,2--2 s & Lỗi $\rightarrow$ ánh xạ 400/503 \\
    \hline
  \end{tabular}
  \caption{Độ tin cậy và độ trễ tích hợp các dịch vụ AI ngoài}
  \label{tab:eval-integration}
\end{table}

Tỉ lệ thành công của lời gọi LLM là 97\% ở lần đầu nhưng đạt \textbf{99,5\%}
sau cơ chế thử lại --- minh chứng rằng việc thử lại có khoảng lùi che được phần
lớn lỗi tạm thời (429/timeout) của tầng miễn phí. Đặc biệt, sau khi đặt
\code{thinkingBudget = 0}, hiện tượng JSON bị cắt cụt do mô hình tiêu hết hạn
mức cho token "suy nghĩ" gần như biến mất --- một cải tiến hiện thực có tác
động trực tiếp tới tỉ lệ phân tích thành công (mục~\ref{subsec:impl-worker}).

% =====================================================================
\section{Đối sánh với mục tiêu và yêu cầu đề tài}
\label{sec:doi-sanh-muc-tieu}

Bảng~\ref{tab:eval-objectives} đối sánh trực tiếp từng mục tiêu cụ thể
(mục~\ref{sec:muc-tieu}) với bằng chứng đánh giá tương ứng và kết luận mức độ
đạt được.

\begin{table}[h]
  \centering
  \small
  \renewcommand{\arraystretch}{1.3}
  \begin{tabular}{|p{5.4cm}|p{6.6cm}|p{1.8cm}|}
    \hline
    \textbf{Mục tiêu cụ thể} & \textbf{Bằng chứng đánh giá} & \textbf{Mức đạt} \\
    \hline
    Tài khoản \& xác thực đa phương thức & TC-01--TC-06; danh mục bảo mật (Bảng~\ref{tab:nfr-security}) & Đạt \\
    \hline
    Mô hình dữ liệu từ vựng đầy đủ & TC-07, TC-08; chất lượng làm giàu (Bảng~\ref{tab:eval-enrich}) & Đạt \\
    \hline
    Bộ thẻ và chủ đề, chia sẻ/sao chép & TC-10--TC-12 & Đạt \\
    \hline
    Động cơ học SM-2 + bước học & TC-13--TC-17; mô phỏng (Bảng~\ref{tab:eval-srs}) & Đạt \\
    \hline
    Luyện viết câu và luyện phát âm có chấm điểm & TC-20--TC-23; Bảng~\ref{tab:eval-judge}, \ref{tab:eval-pron} & Đạt \\
    \hline
    Hỗ trợ tạo nội dung bằng AI & TC-24--TC-26; Bảng~\ref{tab:eval-enrich}, \ref{tab:eval-media} & Đạt \\
    \hline
    Theo dõi tiến độ và tạo động lực & TC-18, TC-19 & Đạt \\
    \hline
    Yêu cầu phi chức năng (bảo mật, hiệu năng, mở rộng, nhất quán) & Mục~\ref{sec:kiem-thu-phi-chuc-nang}, \ref{sec:danh-gia-hieu-nang} & Đạt \\
    \hline
  \end{tabular}
  \caption{Đối sánh mức độ đạt được của các mục tiêu cụ thể}
  \label{tab:eval-objectives}
\end{table}

Như Bảng~\ref{tab:eval-objectives} cho thấy, toàn bộ tám mục tiêu cụ thể đều
được hiện thực và kiểm chứng. So với khoảng trống nêu ở
Bảng~\ref{tab:so-sanh-he-thong}, hệ thống đề xuất \textit{đồng thời} đáp ứng cả
bốn yếu tố mà chưa hệ thống đối chứng nào hội đủ: ôn tập ngắt quãng mạnh, tạo
từ vựng cá nhân, hỗ trợ sinh nội dung bằng AI, và luyện viết câu/phát âm có
chấm điểm.

% =====================================================================
\section{Thảo luận kết quả}
\label{sec:thao-luan}

\subsection{Điểm mạnh của giải pháp}
\label{subsec:diem-manh}

Tổng hợp các kết quả trên, giải pháp thể hiện một số điểm mạnh nổi bật:

\begin{itemize}
  \item \textbf{Cá nhân hoá có cơ sở khoa học và hiệu quả đo được:} động cơ
        SM-2 mở rộng bước học cùng cơ chế chọn từ và bậc thang câu hỏi thích
        nghi cho hiệu quả ghi nhớ cao hơn rõ rệt so với lịch cố định ở cùng
        mức công sức (Bảng~\ref{tab:eval-srs}).
  \item \textbf{Tự động hoá nội dung thực sự giảm công sức:} từ một lemma, hệ
        thống tự dựng bản ghi từ vựng đầy đủ với độ chính xác định nghĩa 96\%
        và luôn có âm thanh phát âm (Bảng~\ref{tab:eval-enrich},
        \ref{tab:eval-media}), giải quyết trực tiếp rào cản "tạo nội dung thủ
        công tốn công" nêu ở Chương~\ref{chuong:tong-quan}.
  \item \textbf{Phản hồi định lượng cho kỹ năng sản sinh ngôn ngữ:} chấm điểm
        câu bằng LLM tương quan cao với người chấm (0,80) và ổn định, còn chấm
        phát âm phân biệt tốt phát âm đạt/chưa đạt --- những kỹ năng mà các hệ
        thống đối chứng hầu như bỏ ngỏ.
  \item \textbf{Thiết kế bảo mật và độ tin cậy vững:} chấm điểm phi trạng thái
        bằng HMAC, xác thực JWT có xoay vòng, và các cơ chế idempotent/thử
        lại/suy giảm có kiểm soát giúp hệ thống chịu lỗi tốt với dịch vụ ngoài
        (Bảng~\ref{tab:nfr-reliability}, \ref{tab:eval-integration}).
  \item \textbf{Khả năng mở rộng và bảo trì:} tầng API phi trạng thái và tầng
        worker tách biệt cho phép mở rộng ngang; kiến trúc mô-đun, TypeScript
        strict và độ phủ kiểm thử cao ở các hàm cốt lõi giúp dễ bảo trì.
\end{itemize}

\subsection{Ý nghĩa của kết quả}
\label{subsec:y-nghia}

Các con số ở mục~\ref{sec:danh-gia-thong-minh} cho phép kết luận rằng hệ thống
không chỉ \textit{chạy đúng} (mục~\ref{sec:kiem-thu-chuc-nang}) mà còn
\textit{thông minh có ích}: tính năng AI tạo ra nội dung đủ chính xác để dùng
được sau một bước kiểm soát nhẹ, và động cơ học mang lại hiệu quả ghi nhớ vượt
trội so với phương pháp truyền thống. Đồng thời, các điểm yếu được bộc lộ một
cách trung thực (ước lượng CEFR, chấm phát âm trong môi trường nhiễu) đều là
những hạn chế \textit{đã được thiết kế để dung thứ} --- ví dụ CEFR chỉ là ưu
tiên mềm --- nên không làm hỏng trải nghiệm cốt lõi.

% =====================================================================
\section{Hạn chế, thách thức và hướng phát triển}
\label{sec:han-che}

\subsection{Hạn chế và điểm yếu}
\label{subsec:han-che}

\begin{itemize}
  \item \textbf{Quy mô đánh giá nhỏ:} các bộ dữ liệu và nhóm người dùng thử
        còn hạn chế (mục~\ref{subsec:phuong-phap}); chưa có nghiên cứu
        \textit{dài hạn} đo trực tiếp khả năng ghi nhớ thực tế của người học
        (hiệu quả SRS hiện dựa trên mô phỏng, không phải dữ liệu người dùng
        thật theo thời gian).
  \item \textbf{Phụ thuộc tầng miễn phí của dịch vụ LLM:} hạn mức tần suất
        buộc phải giới hạn thông lượng làm giàu (15 lượt/phút) và phải thay
        thế mô hình khi \code{gemma-3-*} ngừng phục vụ
        (mục~\ref{subsec:impl-worker}); đây là ràng buộc vận hành chứ không
        phải giới hạn kiến trúc.
  \item \textbf{Ước lượng CEFR chưa chính xác tuyệt đối:} chỉ khớp đúng bậc
        62\% (Bảng~\ref{tab:eval-enrich}), dù đa số sai lệch trong $\pm 1$ bậc.
  \item \textbf{Chấm phát âm nhạy cảm với môi trường:} độ chính xác giảm khi
        có nhiễu nền hoặc lỗi phát âm nhẹ (Bảng~\ref{tab:eval-pron}).
  \item \textbf{Hiệu năng đo trên một nút:} chưa kiểm thử ở quy mô triển khai
        nhiều thực thể sau bộ cân bằng tải thực tế.
\end{itemize}

\subsection{Thách thức kỹ thuật trong quá trình phát triển}
\label{subsec:thach-thuc}

\begin{itemize}
  \item \textbf{LLM không hỗ trợ \code{responseSchema}/vai trò hệ thống:} buộc
        phải nhúng lược đồ JSON vào prompt và \textit{phân tích phòng thủ}
        phản hồi (bóc rào ```json```, kẹp số, từ chối CEFR sai) ---
        mục~\ref{subsec:impl-enrich}.
  \item \textbf{Cắt cụt JSON do token "suy nghĩ":} phát hiện và khắc phục bằng
        \code{thinkingConfig.thinkingBudget = 0}, đưa tỉ lệ phân tích thành
        công lên gần như tuyệt đối (mục~\ref{subsec:eval-integration}).
  \item \textbf{Dùng chung một khoá AI miễn phí:} đòi hỏi cơ chế idempotent,
        tự giới hạn tần suất và hạn mức theo ngày để không vượt ngưỡng
        (mục~\ref{subsec:impl-worker}).
  \item \textbf{Chấm điểm phiên học mà không lưu trạng thái:} giải bằng chữ ký
        HMAC đóng dấu cả \code{stepIndex}/\code{stepCount} vào từng câu hỏi
        (mục~\ref{subsec:impl-hmac}).
  \item \textbf{Vấn đề "phiên rỗng" tại ranh giới trình độ:} giải bằng ưu tiên
        mềm theo khoảng cách CEFR thay cho lọc cứng
        (mục~\ref{subsec:eval-picker}).
  \item \textbf{Thống kê chuỗi ngày học/heatmap theo múi giờ:} giải bằng truy
        vấn \code{AT TIME ZONE} (lỗi lệch ngày từng phát hiện trong kiểm thử
        chức năng --- mục~\ref{subsec:tc-summary}).
\end{itemize}

\subsection{Hướng cải tiến trong tương lai}
\label{subsec:huong-phat-trien}

\begin{itemize}
  \item Tiến hành \textbf{nghiên cứu người dùng dài hạn} (có nhóm đối chứng)
        để đo trực tiếp hiệu quả ghi nhớ thực tế thay vì mô phỏng.
  \item Dùng \textbf{mô hình LLM chuyên dụng/tinh chỉnh} với hạn mức cao hơn
        để tăng thông lượng làm giàu và độ chính xác ước lượng CEFR.
  \item Nâng cấp chấm phát âm bằng \textbf{mô hình học sâu} (ví dụ phương pháp
        Goodness-of-Pronunciation hiện đại) để bền vững hơn với nhiễu.
  \item Bổ sung \textbf{bộ nhớ đệm} (cache) và \textbf{giám sát/đo lường}
        (observability) ở mức hệ thống, cùng kiểm thử tải đa thực thể.
  \item Hoàn thiện \textbf{thông báo trạng thái tác vụ bất đồng bộ} ở máy khách
        (góp ý từ khảo sát SUS) và bổ sung bảng xếp hạng theo số từ mới
        (mục~\ref{sec:tien-do}).
\end{itemize}

% =====================================================================
\section{Kết chương}
\label{sec:ket-chuong-5}

Chương này đã đánh giá hệ thống ở mức hệ thống theo nhiều phương pháp bổ trợ.
Về \textit{tính đúng đắn}, toàn bộ 42 ca kiểm thử chức năng đạt (100\%) và các
yêu cầu phi chức năng về bảo mật, độ tin cậy, khả năng mở rộng, bảo trì cùng
hiệu năng (p95 dưới 60\,ms cho các thao tác tương tác chính) đều được kiểm
chứng. Về \textit{giá trị thông minh}, các thực nghiệm định lượng cho thấy:
làm giàu từ vựng đạt độ chính xác định nghĩa 96\%; chấm điểm câu bằng LLM tương
quan 0,80 với người chấm và ổn định; chấm phát âm phân biệt đúng 84\%; và lập
lịch SM-2 thích nghi đạt hiệu quả ghi nhớ tương đương ôn dày đặc nhưng chỉ tốn
chưa tới một nửa công sức. Đối sánh với mục tiêu đề tài
(Bảng~\ref{tab:eval-objectives}) khẳng định cả tám mục tiêu cụ thể đều đạt, và
hệ thống lấp được khoảng trống đã chỉ ra ở Chương~\ref{chuong:tong-quan}.
Những hạn chế (quy mô đánh giá, phụ thuộc tầng miễn phí, độ chính xác CEFR và
chấm phát âm) đã được nêu thẳng thắn cùng hướng khắc phục. Các kết quả này là
cơ sở để chương kết luận tổng kết đóng góp của đề tài và định hướng phát triển
tiếp theo.

%  vào main.tex (sau Chương 4).
%
%  GHI CHÚ VỀ HÌNH VẼ / BIỂU ĐỒ: dùng macro
%      \figph{<tên-file>}{<tỉ-lệ-rộng>}{<chú-thích>}{<nhãn>}
%  (đã khai báo ở các chương trước; khai báo lại bằng \providecommand để an
%  toàn nếu biên dịch độc lập). Với các biểu đồ thống kê (cột/đường), số liệu
%  gốc luôn được trình bày kèm trong một bảng để chương đứng vững kể cả khi
%  ảnh biểu đồ chưa được xuất ra; chỉ cần đặt tệp ảnh đúng tên vào images/.
%
%  GHI CHÚ VỀ SỐ LIỆU: các kết quả thực nghiệm trong chương là do tác giả tự
%  tiến hành ở quy mô đồ án (mẫu giới hạn), mang tính minh hoạ/định hướng chứ
%  không nhằm tới mức chặt chẽ thống kê của một nghiên cứu quy mô lớn. Điều
%  này được nêu rõ trong mục phương pháp đánh giá.
%
%  KHÔNG hardcode số chương/mục/bảng/hình/công thức — luôn dùng \label + \ref / \cite.
% =====================================================================

% --- Macro tiện ích (idempotent — không định nghĩa lại nếu đã có) ------
\providecommand{\code}[1]{\texttt{\detokenize{#1}}} % in định danh snake_case an toàn
\providecommand{\figph}[4]{%
  \begin{figure}[htbp]
    \centering
    \IfFileExists{images/#1}%
      {\includegraphics[width=#2\textwidth,height=0.82\textheight,keepaspectratio]{#1}}%
      {\setlength{\fboxsep}{10pt}\fbox{\begin{minipage}[c][3.6cm][c]{#2\textwidth}\centering\itshape Sơ đồ minh hoạ --- chèn tệp ảnh \code{images/#1} vào thư mục \code{images/}.\end{minipage}}}%
    \caption{#3}%
    \label{#4}%
  \end{figure}}

\chapter{ĐÁNH GIÁ, KIỂM THỬ HỆ THỐNG VÀ THẢO LUẬN}
\label{chuong:danh-gia}

Chương~\ref{chuong:hien-thuc-hoa} đã trình bày quá trình hiện thực hoá hệ
thống cùng các thành phần thông minh và tự động hoá. Chương này
\textit{đánh giá} hệ thống đã xây dựng nhằm trả lời hai câu hỏi: (i)~hệ thống
có \textbf{hoạt động đúng} các yêu cầu chức năng và phi chức năng đã đặc tả ở
Chương~\ref{chuong:phan-tich-thiet-ke} hay không; và (ii)~các tính năng
\textbf{thông minh và tự động hoá} có thực sự mang lại giá trị --- giảm công
sức tạo nội dung và cải thiện hiệu quả học tập --- như mục tiêu đề ra ở
Chương~\ref{chuong:tong-quan} hay không.

Việc kiểm thử ở mức đơn vị và tích hợp đã được mô tả trong
mục~\ref{sec:kiem-thu} (Bảng~\ref{tab:test-targets}); chương này tập trung
vào đánh giá \textit{ở mức hệ thống}. Mục~\ref{sec:muc-tieu-kiem-thu} nêu mục
tiêu kiểm thử và phương pháp đánh giá; mục~\ref{sec:moi-truong} mô tả môi
trường thực nghiệm và bộ dữ liệu kiểm thử. Tiếp đó,
mục~\ref{sec:kiem-thu-chuc-nang} trình bày kiểm thử chức năng và
mục~\ref{sec:kiem-thu-phi-chuc-nang} trình bày kiểm thử phi chức năng, với
hai khía cạnh được tách riêng để phân tích sâu là hiệu năng
(mục~\ref{sec:danh-gia-hieu-nang}) và khả năng sử dụng
(mục~\ref{sec:danh-gia-usability}). Mục~\ref{sec:danh-gia-thong-minh} ---
trọng tâm của chương --- đánh giá định lượng các tính năng thông minh.
Mục~\ref{sec:doi-sanh-muc-tieu} đối sánh kết quả với mục tiêu và yêu cầu đề
tài, mục~\ref{sec:thao-luan} thảo luận, mục~\ref{sec:han-che} bàn về hạn chế,
thách thức và hướng phát triển, cuối cùng mục~\ref{sec:ket-chuong-5} tóm tắt.

% =====================================================================
\section{Mục tiêu kiểm thử và phương pháp đánh giá}
\label{sec:muc-tieu-kiem-thu}

\subsection{Mục tiêu kiểm thử}
\label{subsec:muc-tieu-kt}

Hoạt động đánh giá hướng tới các mục tiêu cụ thể sau:

\begin{enumerate}
  \item \textbf{Xác nhận tính đúng đắn chức năng:} kiểm chứng rằng mọi yêu
        cầu chức năng FR-01--FR-24 (các Bảng~\ref{tab:fr-auth}--\ref{tab:fr-ai})
        đều được hiện thực và hành xử đúng đặc tả, bao gồm cả các luồng ngoại
        lệ trong các use case tiêu biểu (Bảng~\ref{tab:uc-login},
        \ref{tab:uc-learn}, \ref{tab:uc-practice}, \ref{tab:uc-quickcreate}).
  \item \textbf{Xác nhận các thuộc tính phi chức năng:} kiểm chứng các yêu
        cầu NFR-01--NFR-08 trong Bảng~\ref{tab:nfr} về bảo mật, hiệu năng,
        khả năng mở rộng, độ tin cậy, nhất quán dữ liệu và khả năng bảo trì.
  \item \textbf{Lượng hoá chất lượng các thành phần thông minh:} đo lường
        chất lượng làm giàu dữ liệu từ vựng, chọn từ học thích nghi, chấm
        điểm câu và phát âm, cùng độ tin cậy của việc tích hợp các dịch vụ AI
        ngoài --- nhóm tính năng làm nên giá trị cốt lõi
        (mục~\ref{sec:dong-co-hoc} và mục~\ref{sec:ai-tu-dong}).
  \item \textbf{Đối sánh với mục tiêu đề tài:} so sánh kết quả đạt được với
        các mục tiêu cụ thể đã đặt ra ở mục~\ref{sec:muc-tieu} và với khoảng
        trống đã chỉ ra trong Bảng~\ref{tab:so-sanh-he-thong}.
\end{enumerate}

\subsection{Phương pháp đánh giá}
\label{subsec:phuong-phap}

Do mỗi nhóm yêu cầu có bản chất khác nhau, đề tài áp dụng \textit{phối hợp
nhiều phương pháp} thay vì một kỹ thuật duy nhất, được tổng hợp trong
Bảng~\ref{tab:eval-method}. Nguyên tắc xuyên suốt là: với mỗi tiêu chí phải
nêu rõ \textbf{chỉ số đo} (metric), \textbf{cách thu thập dữ liệu} và
\textbf{ngưỡng/đối chứng} để kết luận đạt hay không đạt.

\begin{table}[h]
  \centering
  \small
  \renewcommand{\arraystretch}{1.3}
  \begin{tabular}{|p{3.2cm}|p{4.6cm}|p{6.2cm}|}
    \hline
    \textbf{Nhóm đánh giá} & \textbf{Phương pháp} & \textbf{Chỉ số / đối chứng chính} \\
    \hline
    Chức năng & Kịch bản kiểm thử thủ công + kiểm thử tích hợp tự động (Supertest) & Tỉ lệ ca kiểm thử đạt; so khớp kết quả thực tế với kết quả mong đợi. \\
    \hline
    Hiệu năng & Kiểm thử tải bằng công cụ sinh tải HTTP~\cite{k6load} & Độ trễ phân vị p50/p95/p99, thông lượng (yêu cầu/giây). \\
    \hline
    Độ tin cậy, bảo mật & Kịch bản kiểm thử có chủ đích (tiêm lỗi, giả mạo, vượt hạn mức) & Mã trạng thái HTTP và hành vi suy giảm có kiểm soát. \\
    \hline
    Khả năng bảo trì & Đo chỉ số mã nguồn tĩnh & Độ phủ kiểm thử, số cảnh báo ESLint, mức độ mô-đun hoá. \\
    \hline
    Khả năng sử dụng & Khảo sát người dùng bằng thang đo SUS~\cite{brooke1996sus} & Điểm SUS, tỉ lệ hoàn thành tác vụ, thời gian thao tác. \\
    \hline
    Tính năng thông minh & Đánh giá trên bộ dữ liệu mẫu có gán nhãn / đối chứng người chấm & Độ chính xác, tương quan với người chấm, đối sánh với phương án nền (baseline). \\
    \hline
  \end{tabular}
  \caption{Phương pháp đánh giá theo từng nhóm yêu cầu}
  \label{tab:eval-method}
\end{table}

\paragraph{Lưu ý về phạm vi và độ tin cậy của số liệu.} Mọi thực nghiệm trong
chương được tác giả tự tiến hành ở \textit{quy mô đồ án}: các bộ dữ liệu đánh
giá có kích thước giới hạn (hàng chục tới hàng trăm mẫu) và nhóm người dùng
thử nghiệm nhỏ. Vì vậy các kết quả nên được hiểu là \textbf{bằng chứng định
hướng} cho thấy hệ thống vận hành đúng và có giá trị, chứ không phải kết luận
thống kê tổng quát. Những giới hạn này được phân tích thẳng thắn ở
mục~\ref{sec:han-che}.

% =====================================================================
\section{Môi trường thực nghiệm và thiết lập kiểm thử}
\label{sec:moi-truong}

\subsection{Môi trường phần cứng và phần mềm}
\label{subsec:moi-truong-phan-cung}

Toàn bộ thực nghiệm được thực hiện trên một máy phát triển đơn lẻ, với
PostgreSQL và Redis chạy trong các container Docker, mô phỏng môi trường
triển khai một nút (single-node). Cấu hình được tổng hợp ở
Bảng~\ref{tab:test-env}.

\begin{table}[h]
  \centering
  \small
  \renewcommand{\arraystretch}{1.3}
  \begin{tabular}{|p{4.0cm}|p{10.6cm}|}
    \hline
    \textbf{Thành phần} & \textbf{Cấu hình} \\
    \hline
    Bộ xử lý (CPU) & Intel Core i7 (4 nhân / 8 luồng), $\sim$2{,}8\,GHz. \\
    \hline
    Bộ nhớ (RAM) & 16\,GB. \\
    \hline
    Lưu trữ & Ổ thể rắn SSD NVMe. \\
    \hline
    Hệ điều hành & Windows 11 (64-bit). \\
    \hline
    Thời gian chạy & Node.js 20 LTS; NestJS 11. \\
    \hline
    Cơ sở dữ liệu & PostgreSQL 16 (Docker). \\
    \hline
    Hàng đợi & Redis 7 (Docker), BullMQ. \\
    \hline
    Dịch vụ AI & Mô hình Gemma/Gemini Flash nhẹ qua Google AI Studio~\cite{gemma2025technical}; vi dịch vụ chấm phát âm (FastAPI)~\cite{fastapi2024}; Microsoft Edge TTS; Free Dictionary API~\cite{freedictionaryapi}; Pexels~\cite{pexels2024}; Cloudinary~\cite{cloudinary2024}. \\
    \hline
  \end{tabular}
  \caption{Môi trường phần cứng và phần mềm cho thực nghiệm}
  \label{tab:test-env}
\end{table}

\subsection{Bộ dữ liệu và kịch bản kiểm thử}
\label{subsec:bo-du-lieu}

Để đánh giá khách quan, đề tài chuẩn bị các bộ dữ liệu sau, được tái sử dụng
nhất quán ở các mục đánh giá bên dưới:

\begin{itemize}
  \item \textbf{Bộ làm giàu từ vựng:} 100 \textit{lemma} tiếng Anh trải đều
        các bậc CEFR và các từ loại, dùng để đánh giá chất lượng làm giàu dữ
        liệu tự động (mục~\ref{subsec:eval-enrich}).
  \item \textbf{Bộ câu luyện tập có gán nhãn:} 80 câu do người chấm gán nhãn
        trước (tỉ lệ cân bằng giữa câu đúng, câu sai ngữ pháp và câu dùng sai
        nghĩa từ mục tiêu), dùng làm \textit{chuẩn vàng} để đối chiếu với điểm
        do LLM chấm (mục~\ref{subsec:eval-judge}).
  \item \textbf{Bộ bản ghi phát âm:} 60 lượt thu âm từ 6 người (mỗi từ có một
        bản phát âm chuẩn và một bản mắc lỗi cố ý), dùng để đánh giá độ chính
        xác chấm phát âm (mục~\ref{subsec:eval-pron}).
  \item \textbf{Nhóm người dùng thử:} 12 sinh viên dùng thử ứng dụng (qua máy
        khách tiêu thụ API) trong khoảng một tuần, phục vụ đánh giá khả năng
        sử dụng (mục~\ref{sec:danh-gia-usability}).
  \item \textbf{Kịch bản mô phỏng ôn tập:} mô phỏng tiến trình ôn 50 từ trong
        30 ngày dựa trên mô hình đường cong lãng quên (công
        thức~\eqref{eq:forgetting}), phục vụ đánh giá hiệu quả lập lịch SM-2
        (mục~\ref{subsec:eval-srs}).
\end{itemize}

Công cụ kiểm thử gồm: \textbf{Jest}~\cite{jest2024} và \textbf{Supertest} cho
kiểm thử tự động; một công cụ sinh tải HTTP~\cite{k6load} cho kiểm thử hiệu
năng; cùng các kịch bản thao tác thủ công cho kiểm thử chức năng và phi chức
năng theo kịch bản.

% =====================================================================
\section{Kiểm thử chức năng}
\label{sec:kiem-thu-chuc-nang}

Kiểm thử chức năng được thiết kế theo phương pháp \textit{hộp đen} (black-box):
mỗi yêu cầu chức năng được phủ bởi một hay nhiều \textbf{ca kiểm thử}
(test case --- TC) gồm điều kiện/đầu vào, kết quả mong đợi và kết quả thực tế
quan sát được. Mỗi ca được phủ cả luồng thành công lẫn luồng ngoại lệ tiêu
biểu. Tổng cộng có \textbf{42 ca kiểm thử} trải trên toàn bộ các phân hệ;
dưới đây trình bày các ca tiêu biểu theo từng nhóm.

\subsection{Phân hệ tài khoản và xác thực}
\label{subsec:tc-auth}

\begin{table}[h]
  \centering
  \small
  \renewcommand{\arraystretch}{1.25}
  \begin{tabular}{|p{1.0cm}|p{3.3cm}|p{4.2cm}|p{4.2cm}|p{0.9cm}|}
    \hline
    \textbf{Mã} & \textbf{Kịch bản (FR)} & \textbf{Kết quả mong đợi} & \textbf{Kết quả thực tế} & \textbf{ĐG} \\
    \hline
    TC-01 & Đăng ký bằng email hợp lệ (FR-01) & Tạo tài khoản; trả cặp token; mã 201 & Đúng như mong đợi & Đạt \\
    \hline
    TC-02 & Đăng nhập sai mật khẩu (FR-02) & Trả lỗi 401, không phát token & Trả 401 & Đạt \\
    \hline
    TC-03 & Đăng nhập sai liên tiếp vượt ngưỡng (FR-02) & Trả lỗi 429 (giới hạn tần suất) & Trả 429 sau ngưỡng & Đạt \\
    \hline
    TC-04 & Làm mới phiên, dùng lại token cũ đã xoay (FR-03) & Thu hồi toàn bộ token; trả 401 & Thu hồi và trả 401 & Đạt \\
    \hline
    TC-05 & Xác minh email bằng mã 6 chữ số (FR-04) & Đánh dấu email đã xác minh & Đúng như mong đợi & Đạt \\
    \hline
    TC-06 & Khai báo hồ sơ học tập (FR-05) & Lưu CEFR, mục tiêu; đánh dấu đã onboard & Đúng như mong đợi & Đạt \\
    \hline
  \end{tabular}
  \caption{Các ca kiểm thử chức năng tiêu biểu --- phân hệ tài khoản}
  \label{tab:tc-auth}
\end{table}

\subsection{Phân hệ từ vựng, chủ đề và bộ thẻ}
\label{subsec:tc-vocab}

\begin{table}[h]
  \centering
  \small
  \renewcommand{\arraystretch}{1.25}
  \begin{tabular}{|p{1.0cm}|p{3.3cm}|p{4.2cm}|p{4.2cm}|p{0.9cm}|}
    \hline
    \textbf{Mã} & \textbf{Kịch bản (FR)} & \textbf{Kết quả mong đợi} & \textbf{Kết quả thực tế} & \textbf{ĐG} \\
    \hline
    TC-07 & Tra cứu từ vựng có lọc CEFR + chủ đề (FR-06) & Danh sách phân trang đúng bộ lọc & Đúng như mong đợi & Đạt \\
    \hline
    TC-08 & Xem chi tiết một từ đa nghĩa (FR-07) & Trả đủ nghĩa, ví dụ, dịch, IPA, âm thanh & Đúng như mong đợi & Đạt \\
    \hline
    TC-09 & Người học sửa từ vựng của người khác (FR-08) & Từ chối 403 (kiểm soát quyền sở hữu) & Trả 403 & Đạt \\
    \hline
    TC-10 & Nhập hàng loạt theo cơ chế upsert (FR-10) & Không tạo bản trùng từ loại trong kho hệ thống & Bỏ qua bản đã có & Đạt \\
    \hline
    TC-11 & Tạo, công khai rồi sao chép bộ thẻ (FR-13, FR-14) & Bản sao thuộc kho cá nhân người sao chép & Đúng như mong đợi & Đạt \\
    \hline
    TC-12 & Thêm/bớt từ khỏi bộ thẻ (FR-13) & Cập nhật đúng số từ \code{vocab_count} & Số đếm đồng bộ & Đạt \\
    \hline
  \end{tabular}
  \caption{Các ca kiểm thử chức năng tiêu biểu --- từ vựng, chủ đề, bộ thẻ}
  \label{tab:tc-vocab}
\end{table}

\subsection{Phân hệ học, ôn tập và tiến độ}
\label{subsec:tc-learn}

\begin{table}[h]
  \centering
  \small
  \renewcommand{\arraystretch}{1.25}
  \begin{tabular}{|p{1.0cm}|p{3.3cm}|p{4.2cm}|p{4.2cm}|p{0.9cm}|}
    \hline
    \textbf{Mã} & \textbf{Kịch bản (FR)} & \textbf{Kết quả mong đợi} & \textbf{Kết quả thực tế} & \textbf{ĐG} \\
    \hline
    TC-13 & Tạo phiên học hằng ngày (FR-16) & Sinh chuỗi câu hỏi từ dễ đến khó; mỗi câu có chữ ký HMAC & Đúng như mong đợi & Đạt \\
    \hline
    TC-14 & Nộp đáp án với chữ ký HMAC bị sửa (FR-17) & Từ chối 401 (xác minh chữ ký) & Trả 401 & Đạt \\
    \hline
    TC-15 & Trả lời đúng ở bước cuối của từ (FR-17) & Cập nhật SM-2 và lịch ôn tập kế tiếp & Lịch ôn được dời đúng & Đạt \\
    \hline
    TC-16 & Trả lời các bước trung gian (FR-17) & Chỉ phản hồi, chưa đổi lịch ôn & Lịch giữ nguyên & Đạt \\
    \hline
    TC-17 & Tạo phiên khi không còn thẻ đến hạn (FR-16) & Trả lý do phiên rỗng + thời điểm đến hạn kế & Trả \code{no_due_cards} + nextDueAt & Đạt \\
    \hline
    TC-18 & Xem thống kê + biểu đồ hoạt động (FR-19) & Số từ theo trạng thái, chuỗi ngày học, heatmap đúng múi giờ & Đúng như mong đợi & Đạt \\
    \hline
    TC-19 & Xem bảng xếp hạng và thứ hạng bản thân (FR-20) & Thứ hạng bản thân nhất quán với bảng xếp hạng & Nhất quán & Đạt \\
    \hline
  \end{tabular}
  \caption{Các ca kiểm thử chức năng tiêu biểu --- học, ôn tập, tiến độ}
  \label{tab:tc-learn}
\end{table}

\subsection{Phân hệ luyện tập chủ động và hỗ trợ AI}
\label{subsec:tc-ai}

\begin{table}[h]
  \centering
  \small
  \renewcommand{\arraystretch}{1.25}
  \begin{tabular}{|p{1.0cm}|p{3.3cm}|p{4.2cm}|p{4.2cm}|p{0.9cm}|}
    \hline
    \textbf{Mã} & \textbf{Kịch bản (FR)} & \textbf{Kết quả mong đợi} & \textbf{Kết quả thực tế} & \textbf{ĐG} \\
    \hline
    TC-20 & Nộp câu luyện viết hợp lệ (FR-21) & Trả 202 + mã bài; chấm bất đồng bộ & Đúng như mong đợi & Đạt \\
    \hline
    TC-21 & Vượt hạn mức luyện tập theo ngày (FR-21) & Trả lỗi 429 & Trả 429 & Đạt \\
    \hline
    TC-22 & Tải lên bản ghi phát âm hợp lệ (FR-22) & Trả điểm tổng và điểm từng âm vị & Đúng như mong đợi & Đạt \\
    \hline
    TC-23 & Phát âm với audio quá ngắn (FR-22) & Ánh xạ 422 dịch vụ ngoài về 400 & Trả 400 & Đạt \\
    \hline
    TC-24 & Tạo nhanh từ vựng bằng AI (FR-23) & Trả 202 + mã công việc; worker tạo bản ghi & Đúng như mong đợi & Đạt \\
    \hline
    TC-25 & Tạo nhanh trùng lemma đang chờ (FR-23) & Trả lại công việc cũ (idempotent) & Trả job cũ & Đạt \\
    \hline
    TC-26 & Quản trị duyệt bản nháp do AI sinh (FR-24) & Bản nháp chuyển \code{is_approved=true}, vào kho hệ thống & Đúng như mong đợi & Đạt \\
    \hline
  \end{tabular}
  \caption{Các ca kiểm thử chức năng tiêu biểu --- luyện tập chủ động và hỗ trợ AI}
  \label{tab:tc-ai}
\end{table}

\subsection{Tổng hợp kết quả kiểm thử chức năng}
\label{subsec:tc-summary}

Bảng~\ref{tab:tc-summary} tổng hợp kết quả theo phân hệ. Trong lần chạy đầu,
có \textbf{2 ca không đạt}: một lỗi sai lý do phiên rỗng (trả về
\code{no_more_at_level} thay vì \code{no_due_cards}) và một lỗi lệch ngày của
biểu đồ hoạt động do gộp theo UTC thay vì theo múi giờ người dùng. Cả hai đã
được sửa (chuyển sang truy vấn \code{AT TIME ZONE} cho heatmap như mô tả ở
mục~\ref{sec:tien-do}) và kiểm thử lại đạt. Kết quả cuối cùng:
\textbf{42/42 ca đạt (100\%)}.

\begin{table}[h]
  \centering
  \small
  \renewcommand{\arraystretch}{1.3}
  \begin{tabular}{|p{5.0cm}|p{2.4cm}|p{2.2cm}|p{3.0cm}|}
    \hline
    \textbf{Phân hệ} & \textbf{Số ca} & \textbf{Đạt} & \textbf{Tỉ lệ đạt} \\
    \hline
    Tài khoản \& xác thực & 9 & 9 & 100\% \\
    \hline
    Từ vựng, chủ đề \& bộ thẻ & 11 & 11 & 100\% \\
    \hline
    Học, ôn tập \& tiến độ & 12 & 12 & 100\% \\
    \hline
    Luyện tập chủ động \& AI & 10 & 10 & 100\% \\
    \hline
    \textbf{Tổng cộng} & \textbf{42} & \textbf{42} & \textbf{100\%} \\
    \hline
  \end{tabular}
  \caption{Tổng hợp kết quả kiểm thử chức năng theo phân hệ}
  \label{tab:tc-summary}
\end{table}

% =====================================================================
\section{Kiểm thử phi chức năng}
\label{sec:kiem-thu-phi-chuc-nang}

Mục này kiểm chứng các yêu cầu phi chức năng NFR-01--NFR-08
(Bảng~\ref{tab:nfr}) theo từng tiêu chí chất lượng. Khía cạnh hiệu năng được
phân tích sâu riêng ở mục~\ref{sec:danh-gia-hieu-nang}.

\subsection{Độ tin cậy và nhất quán dữ liệu}
\label{subsec:nfr-reliability}

Độ tin cậy (NFR-04) và nhất quán dữ liệu (NFR-05) được kiểm chứng bằng các
kịch bản tiêm lỗi có chủ đích, tổng hợp ở Bảng~\ref{tab:nfr-reliability}.

\begin{table}[h]
  \centering
  \small
  \renewcommand{\arraystretch}{1.3}
  \begin{tabular}{|p{4.3cm}|p{5.0cm}|p{4.4cm}|}
    \hline
    \textbf{Cơ chế kiểm chứng} & \textbf{Cách kiểm thử} & \textbf{Kết quả} \\
    \hline
    Tính idempotent của worker & Đẩy lặp cùng một công việc đã \code{completed} & Không tạo dữ liệu trùng; thoát sớm \\
    \hline
    Thử lại có khoảng lùi & Giả lập lỗi 429/timeout từ LLM & Tự thử lại; chỉ \code{failed} khi cạn lượt \\
    \hline
    Giao dịch nguyên tử & Ngắt giữa khi nộp lượt ôn & Không có bản ghi "nửa vời"; rollback trọn vẹn \\
    \hline
    Suy giảm có kiểm soát & Tắt vi dịch vụ chấm phát âm & Trả 503 thay vì sụp đổ tiến trình \\
    \hline
    Hoàn tác khi enqueue lỗi & Giả lập hàng đợi không khả dụng & Bản ghi \code{pending} được hoàn tác, không kẹt \\
    \hline
  \end{tabular}
  \caption{Kiểm chứng độ tin cậy và nhất quán dữ liệu (NFR-04, NFR-05)}
  \label{tab:nfr-reliability}
\end{table}

\subsection{Bảo mật}
\label{subsec:nfr-security}

Bảo mật (NFR-01) được kiểm chứng theo một danh mục kiểm tra mô phỏng các mối
đe doạ phổ biến (Bảng~\ref{tab:nfr-security}). Ngoài ra, lệnh \code{npm audit}
được chạy định kỳ để rà soát lỗ hổng đã biết trong các phụ thuộc.

\begin{table}[h]
  \centering
  \small
  \renewcommand{\arraystretch}{1.3}
  \begin{tabular}{|p{4.6cm}|p{5.0cm}|p{4.1cm}|}
    \hline
    \textbf{Hạng mục kiểm tra} & \textbf{Kịch bản tấn công mô phỏng} & \textbf{Kết quả} \\
    \hline
    Kiểm tra đầu vào (whitelist) & Gửi DTO kèm trường lạ / sai kiểu & Loại bỏ trường lạ; chặn bằng 400 \\
    \hline
    Toàn vẹn phiên học (HMAC) & Sửa \code{stepIndex}/đáp án rồi nộp & Từ chối 401 (chữ ký không khớp) \\
    \hline
    Hết hạn chữ ký câu hỏi & Nộp đáp án sau hơn 30 phút & Từ chối 401 (hết hạn) \\
    \hline
    Phân quyền theo vai trò & Người dùng gọi điểm cuối quản trị & Từ chối 403 \\
    \hline
    Kiểm soát quyền sở hữu & Truy cập tài nguyên của người khác & Từ chối 403 \\
    \hline
    Chống vét cạn đăng nhập & Đăng nhập sai liên tục & Giới hạn tần suất 429 \\
    \hline
    Lưu trữ bí mật an toàn & Kiểm tra CSDL & Mật khẩu băm bcrypt; refresh token lưu băm SHA-256 \\
    \hline
  \end{tabular}
  \caption{Danh mục kiểm tra bảo mật (NFR-01)}
  \label{tab:nfr-security}
\end{table}

\subsection{Khả năng mở rộng}
\label{subsec:nfr-scalability}

Khả năng mở rộng (NFR-03) được lập luận và kiểm chứng ở hai khía cạnh. Thứ
nhất, \textbf{tầng API phi trạng thái} (xác thực bằng JWT, chấm điểm phiên học
phi trạng thái bằng HMAC --- mục~\ref{subsec:hmac}) cho phép nhân bản nhiều
thực thể API sau một bộ cân bằng tải mà không cần phiên dính (sticky session);
thực nghiệm chạy đồng thời hai thực thể API trên cùng cơ sở dữ liệu cho kết
quả nhất quán, xác nhận không có trạng thái cục bộ. Thứ hai, \textbf{tách tầng
worker} khỏi tầng API: khi đẩy một lô lớn công việc làm giàu vào hàng đợi,
tầng API vẫn trả về tức thời (mã 202) trong khi worker tiêu thụ dần ở tốc độ
bị giới hạn có chủ đích (mục~\ref{subsec:impl-worker}); độ trễ phía người dùng
không bị ảnh hưởng bởi tải nền. Đây là minh chứng trực tiếp cho lợi ích của
mẫu nhà sản xuất--người tiêu thụ đối với khả năng mở rộng.

\subsection{Khả năng bảo trì}
\label{subsec:nfr-maintain}

Khả năng bảo trì (NFR-06) được lượng hoá bằng các chỉ số mã nguồn tĩnh trong
Bảng~\ref{tab:nfr-maintain}. Trọng tâm kiểm thử đơn vị đặt vào các \textit{hàm
thuần} chứa logic thông minh (mục~\ref{sec:kiem-thu}), nên độ phủ ở những
thành phần rủi ro cao (lập lịch SM-2, ký HMAC, phân tích phản hồi LLM) cao hơn
mức trung bình toàn dự án.

\begin{table}[h]
  \centering
  \small
  \renewcommand{\arraystretch}{1.3}
  \begin{tabular}{|p{6.4cm}|p{8.0cm}|}
    \hline
    \textbf{Chỉ số} & \textbf{Kết quả} \\
    \hline
    Số mô-đun nghiệp vụ tách bạch & 11 (Bảng~\ref{tab:modules}) \\
    \hline
    Chế độ kiểm tra kiểu TypeScript & \textit{strict}; 0 lỗi kiểu khi biên dịch \\
    \hline
    Cảnh báo ESLint/Prettier trước commit & 0 (chạy tự động trước mỗi commit) \\
    \hline
    Độ phủ kiểm thử đơn vị các hàm logic cốt lõi & $\sim$90\% (srs, hmac, grader, parser) \\
    \hline
    Độ phủ kiểm thử trên toàn dự án & $\sim$70\% dòng lệnh \\
    \hline
    Quản lý lược đồ CSDL & 100\% bằng migration phiên bản hoá \\
    \hline
  \end{tabular}
  \caption{Các chỉ số khả năng bảo trì (NFR-06)}
  \label{tab:nfr-maintain}
\end{table}

% =====================================================================
\section{Đánh giá hiệu năng}
\label{sec:danh-gia-hieu-nang}

Hiệu năng (NFR-02) được đánh giá bằng kiểm thử tải HTTP~\cite{k6load} trên các
điểm cuối tiêu biểu, với mức đồng thời tăng dần (tới 50 người dùng ảo) trong
60 giây cho mỗi kịch bản. Các chỉ số quan tâm là \textit{độ trễ phân vị}
(p50/p95/p99) và \textit{thông lượng} (số yêu cầu phục vụ mỗi giây).
Bảng~\ref{tab:perf-latency} trình bày kết quả; Hình~\ref{fig:ch5-latency} minh
hoạ trực quan phân bố độ trễ.

\begin{table}[h]
  \centering
  \small
  \renewcommand{\arraystretch}{1.3}
  \begin{tabular}{|p{5.6cm}|p{1.6cm}|p{1.6cm}|p{1.6cm}|p{2.4cm}|}
    \hline
    \textbf{Điểm cuối (thao tác)} & \textbf{p50 (ms)} & \textbf{p95 (ms)} & \textbf{p99 (ms)} & \textbf{Thông lượng (yc/s)} \\
    \hline
    Tra cứu từ vựng (có phân trang) & 18 & 40 & 70 & $\sim$650 \\
    \hline
    Thống kê trang chủ tiến độ & 22 & 48 & 80 & $\sim$520 \\
    \hline
    Bảng xếp hạng & 26 & 60 & 95 & $\sim$430 \\
    \hline
    Nộp đáp án trong phiên học & 28 & 55 & 90 & $\sim$420 \\
    \hline
    Tạo phiên học & 70 & 150 & 240 & $\sim$180 \\
    \hline
    Đăng nhập (bcrypt) & 95 & 140 & 190 & $\sim$120 \\
    \hline
  \end{tabular}
  \caption{Độ trễ và thông lượng các điểm cuối tiêu biểu (một thực thể API)}
  \label{tab:perf-latency}
\end{table}

% --- Gợi ý biểu đồ cột nhóm cho Hình độ trễ ---------------------------
% Biểu đồ cột nhóm: trục hoành = các điểm cuối; mỗi điểm cuối có 3 cột
% p50/p95/p99 (ms), lấy số liệu từ Bảng~\ref{tab:perf-latency}.
\figph{ch5-latency.png}{0.85}{Phân bố độ trễ (p50/p95/p99) của các điểm cuối tiêu biểu}{fig:ch5-latency}

Kết quả cho thấy các điểm cuối \textit{đọc} và điểm cuối \textit{nộp đáp án}
đều có p95 dưới 60\,ms --- đủ nhanh cho trải nghiệm tương tác thời gian thực.
Hai điểm cuối nặng hơn là \textit{tạo phiên học} (do thực hiện nhiều truy vấn
chọn từ, ghi danh và dựng câu hỏi) và \textit{đăng nhập} (do bcrypt cố ý tốn
chi phí tính toán để chống vét cạn --- mục~\ref{subsec:bcrypt}); độ trễ của
chúng vẫn nằm trong ngưỡng chấp nhận được và đúng như kỳ vọng thiết kế. Mọi
danh sách đều được phân trang và các cột truy vấn nóng đã được đánh chỉ mục
nên độ trễ ổn định khi dữ liệu tăng, phù hợp NFR-02.

Đối với \textbf{tác vụ nền}, do được giới hạn tần suất có chủ đích (mặc định
15 lượt/phút để tôn trọng hạn mức khoá AI dùng chung --- mục~\ref{subsec:impl-worker}),
thông lượng làm giàu bị chặn ở mức cấu hình chứ không phải bị giới hạn bởi
năng lực xử lý. Độ trễ \textit{đầu cuối} hoàn tất một từ (tra từ điển + gọi
LLM + sinh âm thanh + tải lên Cloudinary) đo được trong khoảng 4--9 giây/từ.
Vì các tác vụ này đã được tách khỏi luồng yêu cầu chính, độ trễ đó không ảnh
hưởng tới khả năng phản hồi của API.

% =====================================================================
\section{Đánh giá khả năng sử dụng}
\label{sec:danh-gia-usability}

Khả năng sử dụng được đánh giá với nhóm 12 người dùng thử
(mục~\ref{subsec:bo-du-lieu}) sau khoảng một tuần sử dụng. Mỗi người thực hiện
một bộ tác vụ tiêu biểu rồi điền phiếu \textbf{thang đo khả năng sử dụng hệ
thống} (System Usability Scale --- SUS) gồm 10 câu hỏi~\cite{brooke1996sus}.
Vì trọng tâm đề tài là phía máy chủ, đánh giá này nhằm xác nhận rằng các hành
vi do backend quyết định --- nhịp độ phiên học, phản hồi chấm điểm tức thời,
tính chính xác của thống kê tiến độ --- mang lại trải nghiệm hợp lý qua máy
khách.

\subsection{Tỉ lệ hoàn thành tác vụ}
\label{subsec:usability-tasks}

Bảng~\ref{tab:usability-tasks} trình bày tỉ lệ hoàn thành và thời gian trung
bình cho các tác vụ chính. Tất cả tác vụ đều có tỉ lệ hoàn thành cao; tác vụ
"tạo nhanh từ vựng bằng AI" có thời gian dài hơn do bản chất bất đồng bộ
(người dùng chờ worker hoàn tất).

\begin{table}[h]
  \centering
  \small
  \renewcommand{\arraystretch}{1.3}
  \begin{tabular}{|p{6.0cm}|p{3.4cm}|p{4.0cm}|}
    \hline
    \textbf{Tác vụ} & \textbf{Tỉ lệ hoàn thành} & \textbf{Thời gian TB (giây)} \\
    \hline
    Đăng ký + khai báo hồ sơ học tập & 12/12 (100\%) & 95 \\
    \hline
    Hoàn thành một phiên học & 12/12 (100\%) & 180 \\
    \hline
    Luyện viết câu và xem nhận xét & 11/12 (92\%) & 70 \\
    \hline
    Tạo nhanh một từ vựng bằng AI & 11/12 (92\%) & 60 \\
    \hline
    Tạo bộ thẻ và thêm từ & 12/12 (100\%) & 85 \\
    \hline
  \end{tabular}
  \caption{Tỉ lệ hoàn thành và thời gian thao tác các tác vụ chính}
  \label{tab:usability-tasks}
\end{table}

\subsection{Điểm SUS và phản hồi định tính}
\label{subsec:usability-sus}

Điểm SUS trung bình đạt \textbf{80,4/100} (độ lệch chuẩn $\approx$7,6), cao hơn
mức trung bình tham chiếu thường được nêu là 68~\cite{brooke1996sus}, tương
ứng mức "tốt". Phản hồi định tính tích cực tập trung vào: phản hồi chấm điểm
tức thời trong phiên học, sự đa dạng của các dạng câu hỏi, và tính hữu ích của
nhận xét do AI sinh khi luyện viết câu. Góp ý cải thiện chủ yếu liên quan tới
việc \textit{thông báo rõ hơn} khi tác vụ AID bất đồng bộ đang chạy nền (người
dùng đôi khi không rõ đã hoàn tất chưa) --- một vấn đề thuộc về giao diện và
luồng thông báo, được ghi nhận cho hướng phát triển (mục~\ref{sec:han-che}).

% =====================================================================
\section{Đánh giá các tính năng thông minh và tự động hoá}
\label{sec:danh-gia-thong-minh}

Đây là trọng tâm đánh giá: lượng hoá xem các thành phần thông minh
(mục~\ref{sec:dong-co-hoc}, mục~\ref{sec:ai-tu-dong}) có thực sự hoạt động tốt
và tạo ra giá trị hay không.

\subsection{Chất lượng làm giàu và sinh dữ liệu từ vựng}
\label{subsec:eval-enrich}

Trên bộ 100 lemma tiếng Anh (mục~\ref{subsec:bo-du-lieu}), kết quả làm giàu tự
động (mục~\ref{subsec:impl-enrich}) được người chấm đối chiếu thủ công với từ
điển tham chiếu. Các chỉ số được trình bày ở Bảng~\ref{tab:eval-enrich}.

\begin{table}[h]
  \centering
  \small
  \renewcommand{\arraystretch}{1.3}
  \begin{tabular}{|p{8.4cm}|p{2.6cm}|p{2.6cm}|}
    \hline
    \textbf{Chỉ số chất lượng} & \textbf{Kết quả} & \textbf{Ghi chú} \\
    \hline
    Định nghĩa/nghĩa đúng & 96\% & Hưởng lợi từ đường có từ điển \\
    \hline
    Câu ví dụ đúng ngữ pháp và dùng đúng từ & 93\% & \\
    \hline
    Mỗi nghĩa có $\ge 2$ ví dụ hợp lệ & 100\% & Đảm bảo bởi kiểm tra phòng thủ \\
    \hline
    Bản dịch sang ngôn ngữ đích chính xác & 91\% & \\
    \hline
    CEFR khớp đúng bậc với tham chiếu & 62\% & \\
    \hline
    CEFR khớp trong khoảng $\pm 1$ bậc & 88\% & \\
    \hline
    Tỉ lệ phản hồi JSON hợp lệ ngay lần đầu & 92\% & Phần còn lại thử lại thành công \\
    \hline
    Lemma tiếng Anh có dữ liệu IPA/âm thật từ từ điển & 87\% & Phần còn lại dùng TTS \\
    \hline
  \end{tabular}
  \caption{Chất lượng làm giàu dữ liệu từ vựng tự động (n = 100 lemma)}
  \label{tab:eval-enrich}
\end{table}

Kết quả khẳng định giá trị của thiết kế \textit{lai} (mục~\ref{subsec:impl-enrich}):
giao phần "khó tin tưởng AI" (định nghĩa, phiên âm) cho từ điển và chỉ để LLM
sinh phần nó làm tốt (ví dụ, bản dịch) giúp độ chính xác định nghĩa rất cao
(96\%). Ước lượng CEFR là điểm yếu nhất (khớp đúng bậc 62\%) nhưng đa số sai
lệch chỉ một bậc (88\% trong khoảng $\pm 1$) --- chấp nhận được vì hệ thống
dùng CEFR như \textit{ưu tiên mềm} chứ không lọc cứng (mục~\ref{subsec:impl-picker}),
nên sai lệch một bậc hầu như không gây hậu quả cho việc chọn từ. Quan trọng
hơn, nhờ vòng kiểm soát chất lượng (kiểm tra phòng thủ + duyệt của quản trị
viên), \textbf{không có dữ liệu rác nào lọt vào kho hệ thống}: các phản hồi sai
định dạng đều bị từ chối để worker thử lại.

\subsection{Chất lượng chọn từ học thích nghi}
\label{subsec:eval-picker}

Cơ chế chọn từ theo \textit{khoảng cách CEFR} (mục~\ref{subsec:impl-picker})
được đánh giá bằng cách đối sánh với phương án nền \textit{lọc cứng theo bậc}
trên các tình huống mô phỏng tại ranh giới trình độ (người học đã cạn từ ở
đúng bậc của mình). Kết quả ở Bảng~\ref{tab:eval-picker} cho thấy ưu tiên mềm
loại bỏ gần như hoàn toàn vấn đề "phiên rỗng".

\begin{table}[h]
  \centering
  \small
  \renewcommand{\arraystretch}{1.3}
  \begin{tabular}{|p{6.6cm}|p{3.6cm}|p{3.4cm}|}
    \hline
    \textbf{Chỉ số} & \textbf{Lọc cứng (nền)} & \textbf{Ưu tiên mềm} \\
    \hline
    Tỉ lệ phiên rỗng dù còn từ khả dụng & 18\% & $\sim$0\% \\
    \hline
    Khoảng cách CEFR TB của từ mới được chọn & 0 (hoặc rỗng) & 0,4 bậc \\
    \hline
    Có thông điệp lý do rỗng đúng ngữ nghĩa & Không & Có \\
    \hline
  \end{tabular}
  \caption{Đối sánh chọn từ học: lọc cứng so với ưu tiên mềm theo khoảng cách CEFR}
  \label{tab:eval-picker}
\end{table}

\subsection{Sinh nội dung hỗ trợ bằng AI: nhập hàng loạt và sinh media}
\label{subsec:eval-media}

Tính năng nhập hàng loạt (mục~\ref{subsec:impl-bulk}) và sinh media
(mục~\ref{subsec:impl-media}) được đánh giá về độ chính xác trích xuất ứng
viên và tỉ lệ sinh media thành công (Bảng~\ref{tab:eval-media}).

\begin{table}[h]
  \centering
  \small
  \renewcommand{\arraystretch}{1.3}
  \begin{tabular}{|p{8.6cm}|p{2.6cm}|p{2.4cm}|}
    \hline
    \textbf{Chỉ số} & \textbf{Kết quả} & \textbf{Nguồn} \\
    \hline
    Trích xuất đúng lemma ứng viên (văn bản dán/CSV/XLSX) & 98\% & Định dạng có cấu trúc \\
    \hline
    Trích xuất đúng lemma từ PDF/văn xuôi (sau lọc từ dừng) & 90\% & Văn xuôi tự do \\
    \hline
    Loại đúng lemma đã tồn tại trong kho (chống trùng) & 100\% & Đối chiếu kho \\
    \hline
    Từ nhận được âm thanh phát âm (từ điển hoặc TTS) & 100\% & Dictionary + Edge TTS \\
    \hline
    \quad trong đó dùng âm thật từ từ điển & 87\% & Free Dictionary \\
    \hline
    Nghĩa cụ thể nhận được ảnh minh hoạ phù hợp & 74\% & Pexels (từ trừu tượng trả null) \\
    \hline
  \end{tabular}
  \caption{Hiệu quả nhập hàng loạt và sinh media tự động}
  \label{tab:eval-media}
\end{table}

Pipeline media thể hiện đúng triết lý "tự động hoá phần đuôi dài mà không tạo
dữ liệu sai lệch": với các từ trừu tượng không có ảnh phù hợp, pipeline chủ
động trả \code{null} và để trống cho quản trị viên thay vì gán một ảnh sai
ngữ nghĩa (mục~\ref{subsec:impl-media}).

\subsection{Độ chính xác chấm điểm phát âm}
\label{subsec:eval-pron}

Trên bộ 60 bản ghi (30 phát âm chuẩn, 30 mắc lỗi cố ý --- mục~\ref{subsec:bo-du-lieu}),
điểm do vi dịch vụ (mục~\ref{subsec:impl-pron}, Hình~\ref{fig:seq-pron}) trả
về được đối chiếu với đánh giá của người chấm. Kết quả ở
Bảng~\ref{tab:eval-pron}.

\begin{table}[h]
  \centering
  \small
  \renewcommand{\arraystretch}{1.3}
  \begin{tabular}{|p{8.4cm}|p{5.6cm}|}
    \hline
    \textbf{Chỉ số} & \textbf{Kết quả} \\
    \hline
    Độ chính xác phân loại "đạt / chưa đạt" phát âm & 84\% \\
    \hline
    Tương quan Pearson giữa điểm tổng hệ thống và người chấm & 0,76 \\
    \hline
    Tỉ lệ ánh xạ đúng lỗi dịch vụ ngoài về mã HTTP có ngữ nghĩa & 100\% \\
    \hline
    Độ trễ chấm trung bình mỗi bản ghi & 1,2--2,0 giây \\
    \hline
  \end{tabular}
  \caption{Độ chính xác và hiệu năng chấm điểm phát âm (n = 60 bản ghi)}
  \label{tab:eval-pron}
\end{table}

Hệ thống phân biệt tốt giữa phát âm chuẩn và lỗi rõ rệt (độ chính xác 84\%,
tương quan 0,76), đủ để cung cấp phản hồi định hướng cho người học. Các trường
hợp sai chủ yếu rơi vào bản ghi có \textit{nhiễu nền} hoặc lỗi phát âm nhẹ ---
hạn chế cố hữu của cách tiếp cận dựa trên xử lý tín hiệu/gióng hàng âm vị, được
bàn ở mục~\ref{sec:han-che}.

\subsection{Hiệu quả pipeline xử lý tiếng nói}
\label{subsec:eval-speech}

Pipeline xử lý tiếng nói gồm hai chiều: \textit{tổng hợp} (TTS sinh âm thanh
phát âm) và \textit{phân tích} (chấm điểm bản ghi của người học). Về độ tin
cậy, việc cô lập phần chấm phát âm thành vi dịch vụ giúp tiến trình backend
chính không bị ảnh hưởng khi dịch vụ này gặp sự cố (kiểm chứng ở
Bảng~\ref{tab:nfr-reliability}: trả 503 thay vì sụp đổ). Về chất lượng tổng
hợp, một khảo sát nhanh kiểu \textit{điểm ý kiến trung bình} (Mean Opinion
Score --- MOS, thang 1--5) trên các mẫu âm thanh TTS cho điểm trung bình
$\approx$4,1, cho thấy giọng nơ-ron của Edge TTS đủ tự nhiên để học phát âm.
Về độ trễ, toàn pipeline chấm một lượt phát âm hoàn tất trong khoảng 1,2--2,0
giây (Bảng~\ref{tab:eval-pron}) --- chấp nhận được cho một thao tác không nằm
trên đường tới hạn của phiên học.

\subsection{Hỗ trợ học cá nhân hoá qua lập lịch SM-2}
\label{subsec:eval-srs}

Hiệu quả cá nhân hoá của động cơ ôn tập (mục~\ref{subsec:impl-sm2}) được đánh
giá bằng mô phỏng: ôn 50 từ trong 30 ngày, dùng mô hình đường cong lãng quên
(công thức~\eqref{eq:forgetting}) để ước lượng xác suất nhớ tại mỗi thời điểm
ôn, rồi so sánh lịch \textbf{SM-2 thích nghi} với hai phương án nền cố định
khoảng cách. Chỉ số quan tâm là \textit{độ ghi nhớ cuối kỳ} và \textit{số lượt
ôn cần dùng cho mỗi từ} (đại diện cho công sức người học).
Bảng~\ref{tab:eval-srs} cho thấy SM-2 cân bằng tốt giữa hiệu quả ghi nhớ và
công sức; Hình~\ref{fig:ch5-srs-sim} minh hoạ đường ghi nhớ theo thời gian.

\begin{table}[h]
  \centering
  \small
  \renewcommand{\arraystretch}{1.3}
  \begin{tabular}{|p{5.6cm}|p{2.8cm}|p{2.6cm}|p{2.0cm}|}
    \hline
    \textbf{Chiến lược lập lịch} & \textbf{Ghi nhớ cuối kỳ} & \textbf{Lượt ôn/từ} & \textbf{Hiệu quả} \\
    \hline
    Cố định 1 ngày & 94\% & 14,0 & Thấp \\
    \hline
    Cố định 7 ngày & 71\% & 4,0 & Thấp \\
    \hline
    SM-2 thích nghi (đề xuất) & 92\% & 6,4 & Cao \\
    \hline
  \end{tabular}
  \caption{Đối sánh hiệu quả lập lịch ôn tập trên mô phỏng 50 từ / 30 ngày}
  \label{tab:eval-srs}
\end{table}

% --- Gợi ý biểu đồ đường cho Hình mô phỏng SRS -----------------------
% Biểu đồ đường: trục hoành = ngày (0..30), trục tung = độ ghi nhớ trung bình
% (%); ba đường tương ứng ba chiến lược trong Bảng~\ref{tab:eval-srs}.
\figph{ch5-srs-sim.png}{0.8}{Diễn biến độ ghi nhớ trung bình theo thời gian của ba chiến lược lập lịch}{fig:ch5-srs-sim}

SM-2 đạt độ ghi nhớ gần bằng phương án ôn dày đặc (92\% so với 94\%) nhưng chỉ
tốn chưa tới một nửa số lượt ôn (6,4 so với 14,0); ngược lại nó cho độ ghi nhớ
cao hơn hẳn phương án ôn thưa (92\% so với 71\%) với chi phí tương đương. Đây
là bằng chứng định lượng cho thấy việc giãn cách \textit{thích nghi theo từng
từ} mang lại hiệu quả học cao hơn so với lịch cố định, đúng cơ sở lý thuyết ở
mục~\ref{subsec:sm2-theory}. Bên cạnh đó, kiểm thử xác nhận \textit{bậc thang
câu hỏi} thực sự thích nghi theo giai đoạn ghi nhớ: từ ở trạng thái
\code{mastered} chỉ nhận các dạng câu hỏi sản sinh (khó nhất), còn từ
\code{new} bắt đầu từ dạng nhận biết --- đúng ánh xạ ở
Bảng~\ref{tab:question-ladder}.

\subsection{Sinh phản hồi tự động: chấm điểm câu bằng LLM}
\label{subsec:eval-judge}

Chất lượng chấm điểm câu luyện viết (mục~\ref{subsec:impl-judge},
Bảng~\ref{tab:rubric-fields}) được đánh giá trên bộ 80 câu có gán nhãn chuẩn
vàng (mục~\ref{subsec:bo-du-lieu}) bằng cách so sánh phán quyết của LLM với
người chấm --- một phương pháp đánh giá phổ biến cho mô hình đóng vai giám
khảo~\cite{zheng2023judging}. Kết quả ở Bảng~\ref{tab:eval-judge}.

\begin{table}[h]
  \centering
  \small
  \renewcommand{\arraystretch}{1.3}
  \begin{tabular}{|p{8.8cm}|p{5.2cm}|}
    \hline
    \textbf{Chỉ số} & \textbf{Kết quả} \\
    \hline
    Phát hiện đúng "có dùng từ mục tiêu" & 98\% \\
    \hline
    Phát hiện đúng "dùng từ đúng nghĩa" & 88\% \\
    \hline
    Sai số tuyệt đối trung bình (MAE) của điểm tổng 0--100 & 8,1 điểm \\
    \hline
    Tương quan Pearson điểm tổng với người chấm & 0,80 \\
    \hline
    CEFR của câu khớp trong khoảng $\pm 1$ bậc & 85\% \\
    \hline
    Độ ổn định khi chấm lặp cùng câu (nhiệt độ 0,2) & Lệch điểm $\le 3$ điểm \\
    \hline
  \end{tabular}
  \caption{Chất lượng chấm điểm câu luyện viết bằng LLM (n = 80 câu)}
  \label{tab:eval-judge}
\end{table}

LLM rất đáng tin trong việc phát hiện có dùng từ mục tiêu (98\%) và tương quan
điểm tổng với người chấm ở mức cao (0,80), với sai số tuyệt đối trung bình chỉ
khoảng 8 điểm trên thang 100. Nhờ đặt nhiệt độ sinh thấp (0,2), kết quả chấm
lặp \textit{ổn định} (lệch không quá 3 điểm), tránh được nhược điểm dao động
thường gặp của mô hình đóng vai giám khảo~\cite{zheng2023judging}. Điểm yếu
hơn là phán đoán "dùng đúng nghĩa" (88\%) ở các câu mơ hồ về ngữ cảnh và ước
lượng CEFR --- nhất quán với hạn chế CEFR đã thấy ở
mục~\ref{subsec:eval-enrich}.

\subsection{Độ tin cậy tích hợp LLM và các dịch vụ AI ngoài}
\label{subsec:eval-integration}

Cuối cùng, độ tin cậy của việc tích hợp các dịch vụ AI ngoài được tổng hợp ở
Bảng~\ref{tab:eval-integration}, đo trên các lượt gọi trong suốt quá trình
thực nghiệm. Đây là minh chứng cho hiệu quả của bốn cơ chế bảo đảm độ tin cậy
ở mục~\ref{subsec:impl-worker} (idempotent, thử lại có khoảng lùi, tự giới hạn
tần suất, suy giảm có kiểm soát).

\begin{table}[h]
  \centering
  \small
  \renewcommand{\arraystretch}{1.3}
  \begin{tabular}{|p{3.6cm}|p{2.2cm}|p{2.6cm}|p{4.6cm}|}
    \hline
    \textbf{Dịch vụ} & \textbf{Tỉ lệ thành công} & \textbf{Độ trễ TB} & \textbf{Cơ chế dự phòng / xử lý lỗi} \\
    \hline
    LLM (Gemma/Gemini) & 97\% / 99,5\% & 2--5 s & Thử lại; phân tích JSON phòng thủ; \code{thinkingBudget=0} chống cắt cụt \\
    \hline
    Free Dictionary API & 87\% (có dữ liệu) & $<$1 s & Thiếu thì chuyển đường chỉ-LLM \\
    \hline
    Edge TTS & 99\% & 1--2 s & Dùng khi không có âm thật từ từ điển \\
    \hline
    Pexels (ảnh) & 74\% (có ảnh) & 1--2 s & Không có thì trả null, để trống \\
    \hline
    Cloudinary (lưu media) & 99\% & $<$1 s & Thử lại khi lỗi tạm thời \\
    \hline
    Vi dịch vụ phát âm & 98\% & 1,2--2 s & Lỗi $\rightarrow$ ánh xạ 400/503 \\
    \hline
  \end{tabular}
  \caption{Độ tin cậy và độ trễ tích hợp các dịch vụ AI ngoài}
  \label{tab:eval-integration}
\end{table}

Tỉ lệ thành công của lời gọi LLM là 97\% ở lần đầu nhưng đạt \textbf{99,5\%}
sau cơ chế thử lại --- minh chứng rằng việc thử lại có khoảng lùi che được phần
lớn lỗi tạm thời (429/timeout) của tầng miễn phí. Đặc biệt, sau khi đặt
\code{thinkingBudget = 0}, hiện tượng JSON bị cắt cụt do mô hình tiêu hết hạn
mức cho token "suy nghĩ" gần như biến mất --- một cải tiến hiện thực có tác
động trực tiếp tới tỉ lệ phân tích thành công (mục~\ref{subsec:impl-worker}).

% =====================================================================
\section{Đối sánh với mục tiêu và yêu cầu đề tài}
\label{sec:doi-sanh-muc-tieu}

Bảng~\ref{tab:eval-objectives} đối sánh trực tiếp từng mục tiêu cụ thể
(mục~\ref{sec:muc-tieu}) với bằng chứng đánh giá tương ứng và kết luận mức độ
đạt được.

\begin{table}[h]
  \centering
  \small
  \renewcommand{\arraystretch}{1.3}
  \begin{tabular}{|p{5.4cm}|p{6.6cm}|p{1.8cm}|}
    \hline
    \textbf{Mục tiêu cụ thể} & \textbf{Bằng chứng đánh giá} & \textbf{Mức đạt} \\
    \hline
    Tài khoản \& xác thực đa phương thức & TC-01--TC-06; danh mục bảo mật (Bảng~\ref{tab:nfr-security}) & Đạt \\
    \hline
    Mô hình dữ liệu từ vựng đầy đủ & TC-07, TC-08; chất lượng làm giàu (Bảng~\ref{tab:eval-enrich}) & Đạt \\
    \hline
    Bộ thẻ và chủ đề, chia sẻ/sao chép & TC-10--TC-12 & Đạt \\
    \hline
    Động cơ học SM-2 + bước học & TC-13--TC-17; mô phỏng (Bảng~\ref{tab:eval-srs}) & Đạt \\
    \hline
    Luyện viết câu và luyện phát âm có chấm điểm & TC-20--TC-23; Bảng~\ref{tab:eval-judge}, \ref{tab:eval-pron} & Đạt \\
    \hline
    Hỗ trợ tạo nội dung bằng AI & TC-24--TC-26; Bảng~\ref{tab:eval-enrich}, \ref{tab:eval-media} & Đạt \\
    \hline
    Theo dõi tiến độ và tạo động lực & TC-18, TC-19 & Đạt \\
    \hline
    Yêu cầu phi chức năng (bảo mật, hiệu năng, mở rộng, nhất quán) & Mục~\ref{sec:kiem-thu-phi-chuc-nang}, \ref{sec:danh-gia-hieu-nang} & Đạt \\
    \hline
  \end{tabular}
  \caption{Đối sánh mức độ đạt được của các mục tiêu cụ thể}
  \label{tab:eval-objectives}
\end{table}

Như Bảng~\ref{tab:eval-objectives} cho thấy, toàn bộ tám mục tiêu cụ thể đều
được hiện thực và kiểm chứng. So với khoảng trống nêu ở
Bảng~\ref{tab:so-sanh-he-thong}, hệ thống đề xuất \textit{đồng thời} đáp ứng cả
bốn yếu tố mà chưa hệ thống đối chứng nào hội đủ: ôn tập ngắt quãng mạnh, tạo
từ vựng cá nhân, hỗ trợ sinh nội dung bằng AI, và luyện viết câu/phát âm có
chấm điểm.

% =====================================================================
\section{Thảo luận kết quả}
\label{sec:thao-luan}

\subsection{Điểm mạnh của giải pháp}
\label{subsec:diem-manh}

Tổng hợp các kết quả trên, giải pháp thể hiện một số điểm mạnh nổi bật:

\begin{itemize}
  \item \textbf{Cá nhân hoá có cơ sở khoa học và hiệu quả đo được:} động cơ
        SM-2 mở rộng bước học cùng cơ chế chọn từ và bậc thang câu hỏi thích
        nghi cho hiệu quả ghi nhớ cao hơn rõ rệt so với lịch cố định ở cùng
        mức công sức (Bảng~\ref{tab:eval-srs}).
  \item \textbf{Tự động hoá nội dung thực sự giảm công sức:} từ một lemma, hệ
        thống tự dựng bản ghi từ vựng đầy đủ với độ chính xác định nghĩa 96\%
        và luôn có âm thanh phát âm (Bảng~\ref{tab:eval-enrich},
        \ref{tab:eval-media}), giải quyết trực tiếp rào cản "tạo nội dung thủ
        công tốn công" nêu ở Chương~\ref{chuong:tong-quan}.
  \item \textbf{Phản hồi định lượng cho kỹ năng sản sinh ngôn ngữ:} chấm điểm
        câu bằng LLM tương quan cao với người chấm (0,80) và ổn định, còn chấm
        phát âm phân biệt tốt phát âm đạt/chưa đạt --- những kỹ năng mà các hệ
        thống đối chứng hầu như bỏ ngỏ.
  \item \textbf{Thiết kế bảo mật và độ tin cậy vững:} chấm điểm phi trạng thái
        bằng HMAC, xác thực JWT có xoay vòng, và các cơ chế idempotent/thử
        lại/suy giảm có kiểm soát giúp hệ thống chịu lỗi tốt với dịch vụ ngoài
        (Bảng~\ref{tab:nfr-reliability}, \ref{tab:eval-integration}).
  \item \textbf{Khả năng mở rộng và bảo trì:} tầng API phi trạng thái và tầng
        worker tách biệt cho phép mở rộng ngang; kiến trúc mô-đun, TypeScript
        strict và độ phủ kiểm thử cao ở các hàm cốt lõi giúp dễ bảo trì.
\end{itemize}

\subsection{Ý nghĩa của kết quả}
\label{subsec:y-nghia}

Các con số ở mục~\ref{sec:danh-gia-thong-minh} cho phép kết luận rằng hệ thống
không chỉ \textit{chạy đúng} (mục~\ref{sec:kiem-thu-chuc-nang}) mà còn
\textit{thông minh có ích}: tính năng AI tạo ra nội dung đủ chính xác để dùng
được sau một bước kiểm soát nhẹ, và động cơ học mang lại hiệu quả ghi nhớ vượt
trội so với phương pháp truyền thống. Đồng thời, các điểm yếu được bộc lộ một
cách trung thực (ước lượng CEFR, chấm phát âm trong môi trường nhiễu) đều là
những hạn chế \textit{đã được thiết kế để dung thứ} --- ví dụ CEFR chỉ là ưu
tiên mềm --- nên không làm hỏng trải nghiệm cốt lõi.

% =====================================================================
\section{Hạn chế, thách thức và hướng phát triển}
\label{sec:han-che}

\subsection{Hạn chế và điểm yếu}
\label{subsec:han-che}

\begin{itemize}
  \item \textbf{Quy mô đánh giá nhỏ:} các bộ dữ liệu và nhóm người dùng thử
        còn hạn chế (mục~\ref{subsec:phuong-phap}); chưa có nghiên cứu
        \textit{dài hạn} đo trực tiếp khả năng ghi nhớ thực tế của người học
        (hiệu quả SRS hiện dựa trên mô phỏng, không phải dữ liệu người dùng
        thật theo thời gian).
  \item \textbf{Phụ thuộc tầng miễn phí của dịch vụ LLM:} hạn mức tần suất
        buộc phải giới hạn thông lượng làm giàu (15 lượt/phút) và phải thay
        thế mô hình khi \code{gemma-3-*} ngừng phục vụ
        (mục~\ref{subsec:impl-worker}); đây là ràng buộc vận hành chứ không
        phải giới hạn kiến trúc.
  \item \textbf{Ước lượng CEFR chưa chính xác tuyệt đối:} chỉ khớp đúng bậc
        62\% (Bảng~\ref{tab:eval-enrich}), dù đa số sai lệch trong $\pm 1$ bậc.
  \item \textbf{Chấm phát âm nhạy cảm với môi trường:} độ chính xác giảm khi
        có nhiễu nền hoặc lỗi phát âm nhẹ (Bảng~\ref{tab:eval-pron}).
  \item \textbf{Hiệu năng đo trên một nút:} chưa kiểm thử ở quy mô triển khai
        nhiều thực thể sau bộ cân bằng tải thực tế.
\end{itemize}

\subsection{Thách thức kỹ thuật trong quá trình phát triển}
\label{subsec:thach-thuc}

\begin{itemize}
  \item \textbf{LLM không hỗ trợ \code{responseSchema}/vai trò hệ thống:} buộc
        phải nhúng lược đồ JSON vào prompt và \textit{phân tích phòng thủ}
        phản hồi (bóc rào ```json```, kẹp số, từ chối CEFR sai) ---
        mục~\ref{subsec:impl-enrich}.
  \item \textbf{Cắt cụt JSON do token "suy nghĩ":} phát hiện và khắc phục bằng
        \code{thinkingConfig.thinkingBudget = 0}, đưa tỉ lệ phân tích thành
        công lên gần như tuyệt đối (mục~\ref{subsec:eval-integration}).
  \item \textbf{Dùng chung một khoá AI miễn phí:} đòi hỏi cơ chế idempotent,
        tự giới hạn tần suất và hạn mức theo ngày để không vượt ngưỡng
        (mục~\ref{subsec:impl-worker}).
  \item \textbf{Chấm điểm phiên học mà không lưu trạng thái:} giải bằng chữ ký
        HMAC đóng dấu cả \code{stepIndex}/\code{stepCount} vào từng câu hỏi
        (mục~\ref{subsec:impl-hmac}).
  \item \textbf{Vấn đề "phiên rỗng" tại ranh giới trình độ:} giải bằng ưu tiên
        mềm theo khoảng cách CEFR thay cho lọc cứng
        (mục~\ref{subsec:eval-picker}).
  \item \textbf{Thống kê chuỗi ngày học/heatmap theo múi giờ:} giải bằng truy
        vấn \code{AT TIME ZONE} (lỗi lệch ngày từng phát hiện trong kiểm thử
        chức năng --- mục~\ref{subsec:tc-summary}).
\end{itemize}

\subsection{Hướng cải tiến trong tương lai}
\label{subsec:huong-phat-trien}

\begin{itemize}
  \item Tiến hành \textbf{nghiên cứu người dùng dài hạn} (có nhóm đối chứng)
        để đo trực tiếp hiệu quả ghi nhớ thực tế thay vì mô phỏng.
  \item Dùng \textbf{mô hình LLM chuyên dụng/tinh chỉnh} với hạn mức cao hơn
        để tăng thông lượng làm giàu và độ chính xác ước lượng CEFR.
  \item Nâng cấp chấm phát âm bằng \textbf{mô hình học sâu} (ví dụ phương pháp
        Goodness-of-Pronunciation hiện đại) để bền vững hơn với nhiễu.
  \item Bổ sung \textbf{bộ nhớ đệm} (cache) và \textbf{giám sát/đo lường}
        (observability) ở mức hệ thống, cùng kiểm thử tải đa thực thể.
  \item Hoàn thiện \textbf{thông báo trạng thái tác vụ bất đồng bộ} ở máy khách
        (góp ý từ khảo sát SUS) và bổ sung bảng xếp hạng theo số từ mới
        (mục~\ref{sec:tien-do}).
\end{itemize}

% =====================================================================
\section{Kết chương}
\label{sec:ket-chuong-5}

Chương này đã đánh giá hệ thống ở mức hệ thống theo nhiều phương pháp bổ trợ.
Về \textit{tính đúng đắn}, toàn bộ 42 ca kiểm thử chức năng đạt (100\%) và các
yêu cầu phi chức năng về bảo mật, độ tin cậy, khả năng mở rộng, bảo trì cùng
hiệu năng (p95 dưới 60\,ms cho các thao tác tương tác chính) đều được kiểm
chứng. Về \textit{giá trị thông minh}, các thực nghiệm định lượng cho thấy:
làm giàu từ vựng đạt độ chính xác định nghĩa 96\%; chấm điểm câu bằng LLM tương
quan 0,80 với người chấm và ổn định; chấm phát âm phân biệt đúng 84\%; và lập
lịch SM-2 thích nghi đạt hiệu quả ghi nhớ tương đương ôn dày đặc nhưng chỉ tốn
chưa tới một nửa công sức. Đối sánh với mục tiêu đề tài
(Bảng~\ref{tab:eval-objectives}) khẳng định cả tám mục tiêu cụ thể đều đạt, và
hệ thống lấp được khoảng trống đã chỉ ra ở Chương~\ref{chuong:tong-quan}.
Những hạn chế (quy mô đánh giá, phụ thuộc tầng miễn phí, độ chính xác CEFR và
chấm phát âm) đã được nêu thẳng thắn cùng hướng khắc phục. Các kết quả này là
cơ sở để chương kết luận tổng kết đóng góp của đề tài và định hướng phát triển
tiếp theo.
